\PassOptionsToPackage{colorlinks=true,allcolors=black}{hyperref}
% REQ-FILE: The above line, \PassOptionsToPackage{} must be the first line.

\documentclass[11pt]{article}

\usepackage[T1]{fontenc}
\usepackage{lmodern}
\usepackage{setspace}
\onehalfspacing
\usepackage{amsmath,amssymb}
\usepackage{array}
\usepackage{booktabs}
\usepackage{amsthm}
\usepackage{url}
\usepackage{natbib}
\bibpunct{(}{)}{;}{a}{}{ }
\usepackage{hyperref}
\usepackage{bookmark}
\usepackage{authblk}
\usepackage{microtype}
\usepackage{enumitem}

\setlist[itemize]{itemsep=0.2em,topsep=0.2em,parsep=0em,partopsep=0em}
\setlist[enumerate]{itemsep=0.2em,topsep=0.2em,parsep=0em,partopsep=0em}
\setlength{\parskip}{0.75em}
\setlength{\parindent}{0em}

\theoremstyle{plain}
\newtheorem{theorem}{Theorem}[section]
\newtheorem{proposition}[theorem]{Proposition}
\newtheorem{lemma}[theorem]{Lemma}
\newtheorem{corollary}[theorem]{Corollary}

\theoremstyle{definition}
\newtheorem{definition}[theorem]{Definition}
\newtheorem{assumption}[theorem]{Assumption}
\newtheorem{example}[theorem]{Example}
\newtheorem{prediction}[theorem]{Prediction}

\theoremstyle{remark}
\newtheorem{remark}[theorem]{Remark}

\providecommand{\keywords}[1]{\vspace{0.5em}\noindent\textbf{Keywords: }\small #1}

% Transformation classifications
\newcommand{\PRS}{\ensuremath{\mathsf{PRS}}}
\newcommand{\BRK}{\ensuremath{\mathsf{BRK}}}
\newcommand{\BF}{\ensuremath{\mathsf{BF}}}
\newcommand{\NEU}{\ensuremath{\mathsf{NEU}}}
\newcommand{\NAcls}{\ensuremath{\mathsf{N/A}}}

% Split identity regimes imported from Referential Regimes
\newcommand{\LOC}{\ensuremath{\mathrm{LOC}}}
\newcommand{\OBJ}{\ensuremath{\mathrm{OBJ}}}
\newcommand{\SCOPEE}{\ensuremath{\mathrm{SCOPE\text{-}E}}}
\newcommand{\SCOPES}{\ensuremath{\mathrm{SCOPE\text{-}S}}}
\newcommand{\RULEC}{\ensuremath{\mathrm{RULE\text{-}C}}}
\newcommand{\RULES}{\ensuremath{\mathrm{RULE\text{-}S}}}

% Formal notation
\newcommand{\Fset}{\ensuremath{\mathcal{F}}}
\newcommand{\Class}{\ensuremath{\mathrm{Class}}}
\newcommand{\Sib}{\ensuremath{\mathrm{Sib}}}
\newcommand{\IT}{\ensuremath{\mathsf{IT}}}
\newcommand{\Ker}{\ensuremath{\mathrel{\approx_{d}}}}
\newcommand{\Dop}{\ensuremath{\Delta^{\mathrm{op}}_{d}}}
\newcommand{\Dsib}{\ensuremath{\Delta^{\mathrm{sib}}}}
\newcommand{\PASS}{\textnormal{\textsc{pass}}}
\newcommand{\FAIL}{\textnormal{\textsc{fail}}}
\newcommand{\INDET}{\textnormal{\textsc{indeterminate}}}

% =============================================================================
% PREFERRED CITATION SPECIFICATION
%
% @techreport{case2026kernel,
%   author      = {Case, Denise M.},
%    title       = {Interpretive Kernel: From Occurrence Equivalence to Frame Equivalence},
%   institution = {Preprint},
%   year        = {2026},
% }
% =============================================================================

\title{
Interpretive Kernel:\\
From Occurrence Equivalence \\
to Frame Equivalence
}

\author{Denise M. Case}
\affil{\small School of Computer Science and Information Systems\\
Northwest Missouri State University, Maryville, MO, USA}
\date{}

\begin{document}
\maketitle
\vspace{-2em}

% arxiv allows 1920 char for meta abstract
\begin{abstract}
  Many operational systems derive outputs
  by interpreting stable declarations against external frames
  that can change independently of those declarations.
  A recurring schedule, for example, may preserve the same stored rule
  while revisions to time-zone or calendar data alter its generated occurrences.
  Whether two frames are equivalent depends both on which interpretations are observed
  and on which output differences count as material.
  This paper formalizes that dependence.
  For finite frame and occurrence carriers, a finite probe family,
  and an interpretation function, the probe-indexed interpretation construction
  transports an occurrence equivalence relation to a frame equivalence relation
  by identifying two frame environments when every selected probe
  produces occurrence-equivalent outputs.
  The induced relation is the \emph{interpretive kernel}:
  the partition of frame environments that the selected probes cannot distinguish.
  A reverse construction generates, from any declared frame regime,
  the least occurrence regime required to support it.
  The two maps form an adjunction between the lattices of occurrence and frame equivalence relations.
  Consequently, interpretive kernels preserve meets, generated demands preserve joins,
  and admissible occurrence regimes form a principal upper set with a least threshold.
  Pairwise thresholds yield an equivalence-relation-valued generalized pseudoultrametric representation of the kernel family.
  Finite counterexamples show that unordered outputs require declared canonicalization,
  favorable matching need not preserve the required lattice structure,
  and frame-selection policies cannot be treated as identity regimes without equivalence-relation validation.
  A computational procedure constructs kernels, thresholds, and finite failure witnesses
  from a cached interpretation table using signatures, generating forests, and union-find.
  The framework is instantiated for recurring schedules using typed civil and resolved outputs,
  chronological canonicalization, finite observation horizons, and versioned frame environments.
  The result makes interpretive equivalence a declared, parameter-relative, and executable relation
  between rules of sameness on opposite sides of interpretation.
\end{abstract}

\keywords{
  structural explainability;
  interpretive kernels;
  occurrence equivalence;
  frame equivalence;
  observational equivalence;
  partition lattices;
  Galois connections;
  generalized ultrametrics;
  schedule semantics;
  time-zone data
}

% =============================================================================
\section{Introduction}
\label{ik.sec.intro}
% =============================================================================

Many operational systems derive outputs
by interpreting a stable declaration against an external frame.
The declaration and the frame may change independently.
A recurring schedule may retain the same rule,
identifier,
and stored representation
while a revised time-zone
or holiday-calendar artifact
changes the occurrences it produces.
Two deployments may therefore hold the same declaration
and produce different outputs.

The converse also occurs.
Distinct frame environments
may produce the same observed occurrences
for selected declarations
and finite observation horizons.
Whether those outputs count as equivalent
may itself depend on whether the comparison preserves
civil values,
resolved instants,
multiplicity,
terminal status,
or some coarser notion of occurrence.
The question
``do these frames produce the same output?''
is therefore incomplete
until the comparison declares
both what is observed
and which output differences count as material.

Scheduling is one instance
of a broader class of interpreted systems.

Whenever a stable declaration
is evaluated against an external,
mutable,
or versioned frame,
the system contains at least two domains of sameness:

\begin{enumerate}
  \item sameness among the interpretive frames; and
  \item sameness among the outputs produced under those frames.
\end{enumerate}

The two domains are related, but not interchangeable.
A rule declaring two frame environments equivalent
may require selected output differences to be ignored.
A rule distinguishing two outputs may separate frames
that another output regime would merge.
This paper focuses on the relation between those choices.

\emph{Neutral Substrates}~\citep{case2026neutral}
establishes the neutrality-by-design constraint
for systems operating under persistent interpretive disagreement.

Represented structure preserves stable reference
without requiring agreement
on every interpretation constructed over it.

\emph{Referential Regimes}~\citep{case2026referential}
makes identity regime-relative
and supplies the equivalence-relation framework
needed to state different rules of sameness
without collapsing them into one privileged answer.

\emph{Operational Identity}~\citep{case2026operational}
represents declared and implemented rules of sameness
as partitions of one finite carrier
and compares them
in the refinement lattice.

This work takes the next step.
It studies how equivalence is transported
across an interpretation function
between two different finite carriers.
One carrier contains frame environments;
the other contains interpreted occurrence environments.
A finite probe family selects which interpretations are observed.
An occurrence regime declares
which observed outputs count as equivalent.
Together, these declarations induce
an equivalence relation on the frame environments.
That induced relation is the \emph{interpretive kernel}.

\subsection{Output Equivalence is Regime-Relative}

Let
\[
  \mathcal{E}
\]
be a finite carrier
of frame environments,
let
\[
  \mathcal{A}
\]
be a finite carrier
of canonical occurrence environments,
and let
\[
  P
\]
be a finite probe family.

An interpretation function
\[
  I:
  P\times\mathcal{E}
  \longrightarrow
  \mathcal{A}
\]
assigns one occurrence environment
to each probe and frame environment.

An occurrence regime
\[
  \rho
  \in
  \operatorname{Eq}(\mathcal{A})
\]
determines which interpretation results
count as equivalent.

The induced frame relation is
\[
  F
  \mathrel{K_P(\rho)}
  F'
  \quad\Longleftrightarrow\quad
  \forall p\in P,\;
  I(p,F)
  \mathrel{\rho}
  I(p,F').
\]

The kernel is therefore relative
to both the probe family and the occurrence regime.
Changing either parameter may change the resulting frame partition.

Strict equality is one possible occurrence regime,
but it is not the only one.
A schedule comparison may distinguish
civil values,
resolved instants,
multiplicity,
terminal status,
or other typed features.
Another comparison may intentionally identify some of those differences.
Neither choice is supplied by the interpretation function alone.

The reverse question is equally important.
Suppose a system declares
a frame regime
\[
  \varphi
  \in
  \operatorname{Eq}(\mathcal{E}).
\]

Which occurrence distinctions
must be ignored
for every pair of frames
identified by \(\varphi\)
to remain indistinguishable
under the selected probes?

This paper constructs
a generated occurrence regime
\[
  L_P(\varphi)
  \in
  \operatorname{Eq}(\mathcal{A})
\]
and proves
\[
  L_P(\varphi)
  \sqsubseteq
  \rho
  \quad\Longleftrightarrow\quad
  \varphi
  \sqsubseteq
  K_P(\rho).
\]

Thus
\[
  L_P
  \dashv
  K_P.
\]

The left adjoint generates the least occurrence regime
supporting a declared frame regime.
The right adjoint returns the frame equivalence
induced by a proposed occurrence regime.
Their adjunction yields an exact threshold:
\[
  \rho
  \text{ supports }
  \varphi
  \quad\Longleftrightarrow\quad
  L_P(\varphi)
  \sqsubseteq
  \rho.
\]

The threshold does not choose which frame or occurrence regime
a system ought to adopt.
It identifies the precise compatibility condition between the two declarations.

The same construction yields a pairwise interpretive demand
\[
  \delta_P(F,F'),
\]
the least occurrence regime under which
\(F\) and \(F'\)
become interpretively equivalent.

These pairwise values form an
\(\operatorname{Eq}(\mathcal{A})\)-valued
generalized pseudoultrametric.

The representation is closely related to established correspondences
between meet-preserving equivalence families and lattice-valued ultrametrics.

The contribution here is the finite interpretive specialization:
the equivalence family is generated from
probes,
interpretation,
and occurrence regimes,
and each distance value
has the operational meaning
of a least supporting regime.

The construction also requires discipline about representation.
An unordered family of occurrences cannot be compared index-wise
without a declared canonicalization.
Replacing canonicalization with optimization over alternative matchings
does not repair the problem.
Finite counterexamples show that such matching constructions
may fail meet preservation and the strong triangle law.
For schedule interpretation, the paper therefore uses
typed chronological canonicalization over finite observation horizons.

\subsection{Contributions}

This paper makes six contributions.

\begin{enumerate}
  \item It defines a finite typed interpretation setting
        with separate frame
        and occurrence carriers,
        a finite probe family,
        and a regime-relative interpretive kernel.

  \item It constructs the generated occurrence-regime map
        \(L_P\)
        and proves the adjunction
        \[
          L_P
          \dashv
          K_P.
        \]
        The adjunction identifies
        the least occurrence regime
        supporting any declared frame regime
        and yields the corresponding
        meet-
        and join-preservation laws.

  \item It defines pairwise interpretive demand
        and gives the generalized-pseudoultrametric representation
        of the kernel family,
        while locating that representation
        within the established lattice-valued-ultrametric correspondence.

  \item It gives finite counterexamples
        showing why unordered presentations require
        declared canonicalization,
        why favorable matching does not preserve
        the required lattice structure,
        and why resolution policies
        cannot be admitted as frame regimes
        without equivalence-relation validation.

  \item It provides a finite computational procedure
        for constructing partitions,
        kernels,
        thresholds,
        and pairwise demands
        using signatures,
        generating forests,
        and union-find,
        together with finite witnesses
        for failed admissibility
        and refinement claims.

  \item It instantiates the framework
        for schedule interpretation
        using typed civil and resolved outputs,
        chronological canonicalization,
        finite recurrence horizons,
        and versioned frame environments,
        and states a preregistered evidence plan
        for testing nondegeneracy,
        regime sensitivity,
        threshold structure,
        and residuals.
\end{enumerate}

The central result is therefore both algebraic and diagnostic.
Equivalence declared on one side of an interpretation function
induces equivalence on the other.
The adjunction computes the exact boundary
between the occurrence regimes
that support a declared frame demand and those that do not.
Because the carriers and probes are finite, that boundary is executable,
and every failed admissibility claim has a finite witness.

Section~\ref{ik.sec.interface} states the imported interface from the prior
Structural Explainability papers and the preconditions on a finite comparison
artifact.
Section~\ref{ik.sec.setting} fixes the finite typed setting: occurrence and
frame carriers, probes, the interpretation function, and the two regime lattices.
Section~\ref{ik.sec.kernel} defines the interpretive kernel induced by a probe
family and occurrence regime, and establishes its monotonicity and probe-union
laws.
Section~\ref{ik.sec.adjoint} constructs the generated occurrence regime, proves
the adjunction \(L_P\dashv K_P\), and derives the meet- and join-preservation
laws.
Section~\ref{ik.sec.thresholds} interprets the adjunction as an exact
lattice-valued threshold between admissible and inadmissible occurrence regimes.
Section~\ref{ik.sec.representation} gives the pairwise interpretive demand and
the generalized-pseudoultrametric representation.
Section~\ref{ik.sec.finite-counterexamples} gives finite counterexamples
delimiting canonicalization, matching, and partition-hood requirements.
Section~\ref{ik.sec.schedule-interpretation} instantiates the construction for
finite schedule interpretation over versioned frame artifacts.
Section~\ref{ik.sec.computational-procedure} gives the near-linear partition
procedure and finite witness extraction.
Section~\ref{ik.sec.evidence-plan} states the preregistered evidence plan.
Section~\ref{ik.sec.related} situates the work,
Section~\ref{ik.sec.limitations} states its limits, and
Section~\ref{ik.sec.conclusion} concludes.

% =============================================================================
\section{Imported Interface and Preconditions}
\label{ik.sec.interface}
% =============================================================================

This paper builds on three results from the Structural Explainability series.

\emph{Neutral Substrates} separates the commitments required for stable shared
reference from the causal, normative, and domain-specific interpretations
that may be constructed over them.
The substrate fixes what is jointly represented
without selecting among the interpretations that may later be applied.

\emph{Referential Regimes} establishes that preservation and invariance
are relative to an explicitly declared basis.
A transformation may preserve an object under one regime
and distinguish it under another.
The relevant question is therefore not whether two representations
are invariant without qualification,
but under which declared regime they remain equivalent.

\emph{Operational Identity} represents declared and implemented rules of sameness
as partitions of a common finite carrier
and compares those partitions in the refinement lattice.
That construction demonstrates that a system's effective equivalence relation
can be made explicit and finitely compared
without identifying equivalence with equality of representation.

This work applies those commitments to interpretation.
It studies two independently variable sources of equivalence:
the regime under which occurrences are compared
and the regime under which interpretive frames are compared.
A finite probe family connects the two.
The resulting construction determines which occurrence distinctions
remain observable under a fixed family of probes and a fixed frame regime.

The earlier papers do not supply the occurrence regimes,
frame regimes,
probe families,
interpretation functions,
or interpretive kernels defined here.
Those objects are introduced in Sections~\ref{ik.sec.setting}
and~\ref{ik.sec.kernel}.

\subsection{Imported Commitments}

The present construction assumes the following commitments from the prior papers.

\begin{enumerate}[label=(I\arabic*)]
  \item \emph{Declared comparison.}
        Equivalence is evaluated relative to an explicitly declared regime.

  \item \emph{Separation of representation and interpretation.}
        A shared carrier may support multiple interpretations
        without incorporating any one interpretation into the carrier itself.

  \item \emph{Finite extensional comparison.}
        On the examined artifact,
        each regime induces a decidable equivalence relation
        and therefore a finite partition of its declared carrier.

  \item \emph{Typed interpretation.}
        The inputs,
        environmental parameters,
        and outputs of an interpretation are fixed by a declared signature.

  \item \emph{Declared output comparison.}
        Outputs are compared only through an explicitly declared
        canonicalization or comparison lifting.
\end{enumerate}

\subsection{Preconditions on the Comparison Artifact}

The constructions of this paper are relative to a finite comparison artifact.
The artifact supplies:

\begin{enumerate}[label=(P\arabic*)]
  \item a finite occurrence carrier;
  \item a finite frame-environment carrier;
  \item decidable occurrence and frame regimes on those carriers;
  \item a finite probe family;
  \item a total interpretation function on the examined domain;
  \item a declared output canonicalization or comparison lifting; and
  \item any finite horizon required to generate the examined occurrences.
\end{enumerate}

The finite horizon is part of the probe declaration.
An unbounded recurrence specification may denote an indefinitely extensible family,
but a finite probe evaluates only the occurrences generated within its declared horizon.

Failure to supply one of these components does not refute an interpretive equivalence.
It leaves the corresponding kernel or threshold undefined for the examined artifact.


% =============================================================================
\section{Mathematical Setting}
\label{ik.sec.setting}
% =============================================================================

This section fixes the finite typed setting used throughout the paper.
The construction distinguishes occurrence environments,
frame environments,
and the probes under which one is interpreted relative to the other.
It then represents occurrence and frame regimes
as equivalence relations on their respective carriers.

\subsection{Finite Carriers}

\begin{definition}[Finite Comparison Artifact]
  \label{ik.def.artifact}
  A \emph{finite comparison artifact} consists of:
  \begin{enumerate}[label=(\alph*)]
    \item a finite set \(\mathcal{A}\) of occurrence environments;
    \item a finite set \(\mathcal{E}\) of frame environments;
    \item a finite probe family \(P\); and
    \item a total interpretation function
          \[
            I:
            P\times\mathcal{E}
            \longrightarrow
            \mathcal{A}.
          \]
  \end{enumerate}
\end{definition}

The finiteness requirement is relative to the examined artifact.
It does not assert that the interpreted domain,
the set of possible frame environments,
or the family of possible probes is globally finite.
It asserts that the carriers and probes admitted to one comparison
are finite and explicitly delimited.

For domains with indefinitely extensible behavior,
finiteness is obtained through the probe declaration.
A recurrence rule may admit occurrences beyond every fixed date,
but a probe evaluates that rule only within its declared finite horizon.
The resulting occurrence environment is therefore finite
even when the underlying specification is unbounded.

\begin{remark}[Artifact Relativity]
  \label{ik.rem.artifact-relative}
  Every partition,
  kernel,
  threshold,
  witness,
  and comparison in this paper
  is relative to a fixed finite comparison artifact.

  Enlarging \(\mathcal{A}\),
  enlarging \(\mathcal{E}\),
  adding probes to \(P\),
  extending a probe horizon,
  or revising the interpretation function
  produces a revised artifact
  on which the constructions must be recomputed.
\end{remark}

\subsection{Partitions and Refinement}

Let \(X\) be a finite set.
Write \(\operatorname{Eq}(X)\)
for the set of equivalence relations on \(X\).

Each \(\alpha\in\operatorname{Eq}(X)\)
induces the partition
\[
  X/\alpha
  =
  \bigl\{
  [x]_{\alpha}
  \bigm|
  x\in X
  \bigr\}.
\]
The relation and its induced partition
will be used interchangeably when no ambiguity results.

\begin{definition}[Refinement Order]
  \label{ik.def.refinement}
  For \(\alpha,\beta\in\operatorname{Eq}(X)\),
  write
  \[
    \alpha\sqsubseteq\beta
  \]
  when
  \[
    \alpha\subseteq\beta
  \]
  as subsets of \(X\times X\).

  Equivalently,
  \(\alpha\sqsubseteq\beta\)
  when every block of \(X/\alpha\)
  is contained in a block of \(X/\beta\).
  Thus \(\alpha\) is finer than \(\beta\),
  and \(\beta\) is coarser than \(\alpha\).
\end{definition}

Under \(\sqsubseteq\),
\(\operatorname{Eq}(X)\) is a finite lattice.
Its least element is the diagonal relation
\[
  \bot_X
  =
  \{
  (x,x)
  \mid
  x\in X
  \},
\]
and its greatest element is
\[
  \top_X
  =
  X\times X.
\]

For \(\alpha,\beta\in\operatorname{Eq}(X)\),
their meet is intersection:
\[
  \alpha\wedge\beta
  =
  \alpha\cap\beta.
\]
Their join is the equivalence closure of their union:
\[
  \alpha\vee\beta
  =
  \operatorname{EqCl}_X(\alpha\cup\beta).
\]

More generally,
for a relation \(R\subseteq X\times X\),
write
\[
  \operatorname{EqCl}_X(R)
\]
for the least equivalence relation on \(X\)
containing \(R\).

\begin{definition}[Generating Relation]
  \label{ik.def.generating}
  Let \(\alpha\in\operatorname{Eq}(X)\).
  A relation \(G\subseteq X\times X\)
  \emph{generates} \(\alpha\) when
  \[
    \operatorname{EqCl}_X(G)
    =
    \alpha.
  \]
\end{definition}

A generating relation need not contain every pair in \(\alpha\).
For a finite equivalence class,
a connected set of generating pairs is sufficient.
Section~\ref{ik.sec.computational-procedure}
uses this fact to replace all within-class pairs
by a spanning forest.

\subsection{Occurrence and Frame Environments}

The two carriers represent different kinds of variation.

\begin{definition}[Occurrence Environment]
  \label{ik.def.occurrence-environment}
  An \emph{occurrence environment}
  \(A\in\mathcal{A}\)
  is a finite canonical representation
  of the occurrences observed by a probe.

  The representation contains every component
  declared relevant to occurrence comparison,
  including any ordering,
  multiplicity,
  boundary,
  or termination information
  required by the domain signature.
\end{definition}

An occurrence environment is not necessarily a single occurrence.
It may be a finite sequence,
a finite indexed family,
a bounded trace,
or another finite structured output.
Its internal form is fixed by the declared interpretation signature.

When occurrence environments contain ordered outputs,
the order is part of their canonical representation.
For the schedule interpretation of Section~\ref{ik.sec.schedule-interpretation},
that order is chronological.
The comparison does not minimize over arbitrary bijections
between generated occurrences.

A terminal or end tag is also part of the canonical representation
when the domain distinguishes completion,
truncation,
or exhaustion of generated output.
Differences in length or termination are then ordinary represented differences
rather than exceptional cases outside the comparison.

\begin{definition}[Frame Environment]
  \label{ik.def.frame-environment}
  A \emph{frame environment}
  \(F\in\mathcal{E}\)
  is a complete assignment
  to the frame variables of the declared interpretation signature.

  Frame variables are parameters
  under which the same probe may be interpreted differently
  without changing the probe declaration itself.
\end{definition}

A frame environment may include,
for example,
a resolution policy,
a calendar convention,
a time-zone database release,
a locale,
or another declared interpretive context.
The abstract construction does not prescribe
which variables belong to the frame.
That choice is part of the domain interpretation.

\begin{definition}[Occurrence and Frame Regimes]
  \label{ik.def.regimes}
  An \emph{occurrence regime} is an equivalence relation
  \[
    \rho\in\operatorname{Eq}(\mathcal{A}).
  \]
  A \emph{frame regime} is an equivalence relation
  \[
    \varphi\in\operatorname{Eq}(\mathcal{E}).
  \]
\end{definition}

An occurrence regime declares
which occurrence environments count as equivalent.
A frame regime declares
which frame environments count as equivalent.
Neither equivalence is inferred from raw representation equality
unless equality is itself the declared regime.

The two regime lattices are
\[
  \operatorname{Eq}(\mathcal{A})
  \qquad\text{and}\qquad
  \operatorname{Eq}(\mathcal{E}),
\]
each ordered by refinement as in Definition~\ref{ik.def.refinement}.

\subsection{Interpretation Functions}

\begin{definition}[Interpretation Function]
  \label{ik.def.interpretation}
  For a finite probe family \(P\),
  an \emph{interpretation function} is a total function
  \[
    I:
    P\times\mathcal{E}
    \longrightarrow
    \mathcal{A}.
  \]

  For \(p\in P\) and \(F\in\mathcal{E}\),
  the value
  \[
    I(p,F)
  \]
  is the occurrence environment obtained
  by interpreting probe \(p\)
  under frame environment \(F\).
\end{definition}

Totality is required on the examined artifact.
A missing,
invalid,
or inapplicable interpretation result
must therefore be represented by a declared value in \(\mathcal{A}\)
or excluded by a precondition on the artifact.
It is not treated as an implicit absence of output.

The interpretation function may be factored through raw output.
Let \(\mathcal{O}\) be a finite raw-output carrier,
let
\[
  J:
  P\times\mathcal{E}
  \longrightarrow
  \mathcal{O}
\]
be a raw interpretation,
and let
\[
  c:
  \mathcal{O}
  \longrightarrow
  \mathcal{A}
\]
be a declared canonicalization.
Then
\[
  I
  =
  c\circ J.
\]
The abstract results depend on \(I\),
not on the internal procedure used to compute it.

For a fixed frame environment \(F\), the section
\(I({\cdot},F)\colon P \to \mathcal{A}\)
plays the role of a valuation function familiar from denotational semantics
\citep[\S4.1]{schmidt1986}:
it assigns to each probe its meaning as an occurrence environment.
What distinguishes the present setting is that the frame is a parameter.
The same probe is valued against many frames,
and the object of study is the equivalence induced on frames
when two of them value every selected probe alike.

A domain realization must nevertheless declare
how raw output is transformed into the occurrence environment
on which its occurrence regimes are defined.

\begin{remark}[Canonicalization Is Part of the Declaration]
  \label{ik.rem.canonicalization}
  Canonicalization is not an optimization applied after comparison.
  It determines the carrier on which comparison is performed.

  Two outputs that differ only by an undeclared presentation order
  may be mapped to the same canonical occurrence environment.
  Two outputs whose order is semantically declared
  remain distinguishable through that order.

  Replacing declared canonicalization
  by minimization over all bijections
  defines a different comparison
  and may produce different kernels.
\end{remark}

\begin{remark}[Comparison Liftings]
  \label{ik.rem.lifting}
  A domain may define an occurrence regime directly on \(\mathcal{A}\),
  or obtain it by lifting an equivalence relation
  on simpler output values
  to sequences,
  finite families,
  or other structured occurrence environments.

  Such a lifting must be declared before the interpretive kernel is computed.
  The kernel construction receives the resulting occurrence regime \(\rho\);
  it does not choose the lifting
  or infer one from the observed outputs.
\end{remark}

\subsection{Probes}

\begin{definition}[Probe]
  \label{ik.def.probe}
  A \emph{probe} \(p\)
  is a finite observation declaration
  sufficient to determine \(I(p,F)\)
  for every frame environment \(F\in\mathcal{E}\).

  A probe fixes:
  \begin{enumerate}[label=(\alph*)]
    \item the interpreted input;
    \item the finite observation horizon;
    \item the requested occurrence output; and
    \item any additional observation parameters
          not supplied by the frame environment.
  \end{enumerate}
\end{definition}

The horizon is constitutive of the probe.
Two declarations that differ only in horizon
are distinct probes
because they may produce different occurrence environments.

An unbounded recurrence specification
does not produce an infinite occurrence environment under a probe.
The probe horizon determines the finite generated prefix,
and the canonical occurrence environment records
the declared boundary or end condition.

\begin{definition}[Probe Family]
  \label{ik.def.probe-family}
  A \emph{probe family}
  \(P\)
  is a finite set of probes
  evaluated against the same frame carrier \(\mathcal{E}\)
  and occurrence carrier \(\mathcal{A}\)
  through one interpretation function
  \[
    I:
    P\times\mathcal{E}
    \longrightarrow
    \mathcal{A}.
  \]
\end{definition}

The probe family is the observational basis of the comparison.
A frame distinction invisible to one probe family
may become visible when another probe is added.
The interpretive kernels of Section~\ref{ik.sec.kernel}
are therefore indexed by \(P\).

\begin{remark}[Both Parameters Are Required]
  \label{ik.rem.two-parameters}
  A probe family alone does not determine an interpretive kernel.
  The same probes may identify different frame environments
  under different occurrence regimes.

  An occurrence regime alone does not determine an interpretive kernel.
  The same regime may identify different frame environments
  under different probe families.

  Every kernel reported in this paper
  is therefore indexed by both
  a probe family \(P\)
  and an occurrence regime \(\rho\).
\end{remark}

The standing mathematical data for the next section are thus
\[
  P,
  \qquad
  \mathcal{A},
  \qquad
  \mathcal{E},
  \qquad
  I:P\times\mathcal{E}\longrightarrow\mathcal{A},
\]
together with an occurrence regime
\[
  \rho\in\operatorname{Eq}(\mathcal{A})
\]
or a frame regime
\[
  \varphi\in\operatorname{Eq}(\mathcal{E}),
\]
according to the direction of the construction.


% =============================================================================
\section{Interpretive Kernels}
\label{ik.sec.kernel}
% =============================================================================

The interpretation function assigns an occurrence environment
to each probe and frame environment.
An occurrence regime determines which of those outputs count as equivalent.
The resulting pattern of probe-relative equivalence
induces an equivalence relation on the frame environments.

\subsection{Kernel Induced by Probe Family and Occurrence Regime}

Fix a finite comparison artifact
\[
  \bigl(
  P,
  \mathcal{A},
  \mathcal{E},
  I
  \bigr)
\]
as in Definition~\ref{ik.def.artifact},
and let
\[
  \rho\in\operatorname{Eq}(\mathcal{A})
\]
be an occurrence regime.

\begin{definition}[Probe Signature]
  \label{ik.def.probe-signature}
  The \emph{probe signature} induced by \(P\) and \(\rho\) is the function
  \[
    \Sigma_{P,\rho}:
    \mathcal{E}
    \longrightarrow
    \bigl(\mathcal{A}/\rho\bigr)^{P}
  \]
  defined by
  \[
    \Sigma_{P,\rho}(F)(p)
    =
    [I(p,F)]_{\rho}
  \]
  for every \(F\in\mathcal{E}\) and \(p\in P\).
\end{definition}

The probe signature records,
for each frame environment,
the occurrence-equivalence class produced by every probe.
Its codomain is a function space indexed by \(P\),
so the definition requires no arbitrary ordering of the probe family.

\begin{definition}[Interpretive Kernel]
  \label{ik.def.kernel}
  The \emph{interpretive kernel} induced by the probe family \(P\)
  and occurrence regime \(\rho\) is
  \[
    K_P(\rho)
    =
    \ker(\Sigma_{P,\rho}).
  \]
  Equivalently,
  for \(F,F'\in\mathcal{E}\),
  \[
    F\mathrel{K_P(\rho)}F'
  \]
  when
  \[
    I(p,F)\mathrel{\rho}I(p,F')
    \qquad
    \text{for every }p\in P.
  \]
\end{definition}

Two frame environments lie in the same block of \(K_P(\rho)\)
when no probe in \(P\) distinguishes them
under the declared occurrence regime \(\rho\).

\begin{proposition}[The Kernel Is a Frame Regime]
  \label{ik.prop.kernel-equivalence}
  For every finite probe family \(P\)
  and occurrence regime
  \(\rho\in\operatorname{Eq}(\mathcal{A})\),
  the relation
  \[
    K_P(\rho)
  \]
  is an equivalence relation on \(\mathcal{E}\).
  Hence
  \[
    K_P(\rho)
    \in
    \operatorname{Eq}(\mathcal{E}).
  \]
\end{proposition}

\begin{proof}
  The relation \(K_P(\rho)\)
  is the kernel of the function \(\Sigma_{P,\rho}\),
  so it is an equivalence relation.

  Explicitly,
  reflexivity follows because
  \(I(p,F)\mathrel{\rho}I(p,F)\)
  for every \(p\in P\).

  Symmetry follows because
  \(I(p,F)\mathrel{\rho}I(p,F')\)
  implies
  \(I(p,F')\mathrel{\rho}I(p,F)\).

  Transitivity follows because
  \[
    I(p,F)\mathrel{\rho}I(p,F')
    \qquad\text{and}\qquad
    I(p,F')\mathrel{\rho}I(p,F'')
  \]
  imply
  \[
    I(p,F)\mathrel{\rho}I(p,F'')
  \]
  for every \(p\in P\).
\end{proof}

\begin{definition}[Interpretive Partition]
  \label{ik.def.interpretive-partition}
  The quotient
  \[
    \pi_{P,\rho}
    =
    \mathcal{E}/K_P(\rho)
  \]
  is the \emph{interpretive partition}
  induced by \(P\) and \(\rho\).
\end{definition}

Each block of \(\pi_{P,\rho}\)
contains the frame environments
that produce the same probe signature
up to the occurrence distinctions preserved by \(\rho\).

The kernel reports observational equivalence
relative to the declared comparison artifact.
It does not assert that two frame environments
are identical in representation,
interchangeable for every purpose,
or equivalent under probes outside \(P\).

\begin{remark}[Kernel Terminology]
  \label{ik.rem.kernel-terminology}
  The term \emph{kernel} is literal.
  The relation \(K_P(\rho)\)
  is the kernel equivalence of the quotient-valued signature map
  \(\Sigma_{P,\rho}\).

  It is also interpretation-relative.
  Replacing \(I\),
  \(P\),
  or \(\rho\)
  may change the signature map
  and therefore its kernel.
\end{remark}

\subsection{Both Parameters Are Constitutive}

The notation \(K_P(\rho)\)
retains both parameters because neither can be recovered from the other.

The probe family determines
which interpretive differences are exercised.
The occurrence regime determines
which differences in the resulting occurrence environments
count as observable.

\begin{proposition}[Probe Dependence]
  \label{ik.prop.probe-dependence}
  There exist a finite interpretation function,
  an occurrence regime \(\rho\),
  and finite probe families
  \[
    P\subsetneq Q
  \]
  such that
  \[
    K_Q(\rho)
    \subsetneq
    K_P(\rho).
  \]
\end{proposition}

\begin{proof}
  Let
  \[
    \mathcal{E}
    =
    \{F_0,F_1\},
    \qquad
    \mathcal{A}
    =
    \{a,b\},
  \]
  and let \(\rho=\bot_{\mathcal{A}}\)
  be equality on \(\mathcal{A}\).

  Take distinct probes \(p\) and \(q\),
  with
  \[
    I(p,F_0)=I(p,F_1)=a,
  \]
  and
  \[
    I(q,F_0)=a,
    \qquad
    I(q,F_1)=b.
  \]

  Let
  \[
    P=\{p\},
    \qquad
    Q=\{p,q\}.
  \]

  Probe \(p\) does not distinguish the two frame environments,
  so
  \[
    K_P(\rho)
    =
    \top_{\mathcal{E}}.
  \]

  Probe \(q\) distinguishes them under equality,
  so
  \[
    K_Q(\rho)
    =
    \bot_{\mathcal{E}}.
  \]

  Therefore
  \[
    K_Q(\rho)
    \subsetneq
    K_P(\rho).
  \]
\end{proof}

\begin{proposition}[Occurrence-Regime Dependence]
  \label{ik.prop.regime-dependence}
  There exist a finite probe family \(P\)
  and occurrence regimes
  \[
    \rho\subsetneq\sigma
  \]
  such that
  \[
    K_P(\rho)
    \subsetneq
    K_P(\sigma).
  \]
\end{proposition}

\begin{proof}
  Let
  \[
    \mathcal{E}
    =
    \{F_0,F_1\},
    \qquad
    \mathcal{A}
    =
    \{a,b\},
  \]
  and let \(P=\{p\}\).

  Define
  \[
    I(p,F_0)=a,
    \qquad
    I(p,F_1)=b.
  \]

  Let
  \[
    \rho=\bot_{\mathcal{A}}
  \]
  and
  \[
    \sigma=\top_{\mathcal{A}}.
  \]

  Under \(\rho\),
  the outputs \(a\) and \(b\) are distinct,
  so
  \[
    K_P(\rho)
    =
    \bot_{\mathcal{E}}.
  \]

  Under \(\sigma\),
  every occurrence environment is equivalent,
  so
  \[
    K_P(\sigma)
    =
    \top_{\mathcal{E}}.
  \]

  Hence
  \[
    K_P(\rho)
    \subsetneq
    K_P(\sigma).
  \]
\end{proof}

Propositions~\ref{ik.prop.probe-dependence}
and~\ref{ik.prop.regime-dependence}
show that neither parameter is decorative.

Reporting only \(P\)
omits the rule under which outputs were compared.
Reporting only \(\rho\)
omits the observations under which frame equivalence was tested.
A reported interpretive kernel must identify both.

\begin{remark}[No Unqualified Frame Equivalence]
  \label{ik.rem.no-unqualified-equivalence}
  The construction does not define an unqualified statement
  that two frame environments are interpretively equivalent.

  It defines the indexed statement
  \[
    F\mathrel{K_P(\rho)}F',
  \]
  meaning that the frame environments agree
  under the probes in \(P\)
  up to the occurrence distinctions retained by \(\rho\).
\end{remark}

\subsection{Canonicalization and Comparison Liftings}

The kernel compares elements of the occurrence carrier \(\mathcal{A}\).
The construction therefore begins only after raw interpretive output
has been converted into the canonical occurrence representation
on which the occurrence regime is defined.

Suppose the domain realization supplies
a raw-output carrier \(\mathcal{O}\),
a raw interpretation
\[
  J:
  P\times\mathcal{E}
  \longrightarrow
  \mathcal{O},
\]
and a canonicalization
\[
  c:
  \mathcal{O}
  \longrightarrow
  \mathcal{A}.
\]
Then
\[
  I
  =
  c\circ J,
\]
and Definition~\ref{ik.def.kernel}
becomes
\[
  F\mathrel{K_P(\rho)}F'
  \quad\Longleftrightarrow\quad
  c(J(p,F))
  \mathrel{\rho}
  c(J(p,F'))
  \text{ for every }p\in P.
\]

Canonicalization fixes the represented structure
before comparison begins.
It may normalize distinctions
that the domain declares representational,
but it must preserve distinctions
that the domain declares interpretively relevant.

\begin{definition}[Comparison Lifting]
  \label{ik.def.comparison-lifting}
  Let \(V\) be a finite component-value carrier
  and let
  \[
    \delta\in\operatorname{Eq}(V)
  \]
  be a component equivalence.

  A \emph{comparison lifting} for the occurrence carrier \(\mathcal{A}\)
  assigns to \(\delta\) an occurrence regime
  \[
    \widehat{\delta}
    \in
    \operatorname{Eq}(\mathcal{A}).
  \]
\end{definition}

The lifting states how component equivalence
extends to complete occurrence environments.
For sequence-valued environments,
the lifting may require equal length,
index-wise component equivalence,
and agreement of declared terminal information.
For other structured outputs,
a different lifting may be appropriate.

The kernel receives the lifted occurrence regime
\[
  \widehat{\delta}
\]
and computes
\[
  K_P(\widehat{\delta}).
\]
The component relation \(\delta\) alone
does not determine the comparison
until its lifting to \(\mathcal{A}\) has been declared.

\begin{definition}[Monotone Comparison Lifting]
  \label{ik.def.monotone-lifting}
  A comparison lifting is \emph{monotone}
  when
  \[
    \delta\sqsubseteq\delta'
  \]
  implies
  \[
    \widehat{\delta}
    \sqsubseteq
    \widehat{\delta'}.
  \]
\end{definition}

Monotonicity says that weakening component comparison
cannot make the lifted comparison stricter.
Any stronger structural laws required of a domain lifting
must be stated as part of that domain's declared comparison protocol.

\begin{remark}[Canonical Order Is Not Matching]
  \label{ik.rem.canonical-order}
  When the occurrence environment carries a declared canonical order,
  comparison is performed in that order.

  The construction does not minimize over bijections
  between the members of two occurrence environments.
  Such minimization would quotient away the declared order
  and define a different occurrence regime.

  In the schedule interpretation,
  generated occurrences are placed in chronological order.
  Comparison is then index-wise under the declared lifting,
  with length and terminal information retained
  as components of the canonical occurrence environment.
\end{remark}

\begin{remark}[Canonicalization Precedes the Kernel]
  \label{ik.rem.canonicalization-precedes}
  Canonicalization and comparison lifting
  are not operations performed by \(K_P\).

  They determine the occurrence carrier
  and the occurrence regime supplied to \(K_P\).
  The kernel then transports that declared output comparison
  to a frame regime through the interpretation function.
\end{remark}

\subsection{Monotonicity}

The interpretive-kernel construction is monotone
in the occurrence regime.

\begin{proposition}[Monotonicity in the Occurrence Regime]
  \label{ik.prop.kernel-monotone}
  Let
  \[
    \rho,\sigma
    \in
    \operatorname{Eq}(\mathcal{A}).
  \]
  If
  \[
    \rho\sqsubseteq\sigma,
  \]
  then
  \[
    K_P(\rho)
    \sqsubseteq
    K_P(\sigma).
  \]
\end{proposition}

\begin{proof}
  Suppose
  \[
    F\mathrel{K_P(\rho)}F'.
  \]
  Then
  \[
    I(p,F)\mathrel{\rho}I(p,F')
  \]
  for every \(p\in P\).

  Because
  \[
    \rho\subseteq\sigma,
  \]
  it follows that
  \[
    I(p,F)\mathrel{\sigma}I(p,F')
  \]
  for every \(p\in P\).

  Hence
  \[
    F\mathrel{K_P(\sigma)}F',
  \]
  so
  \[
    K_P(\rho)
    \subseteq
    K_P(\sigma).
  \]
\end{proof}

A coarser occurrence regime
identifies at least as many occurrence environments.
The induced interpretive kernel therefore identifies
at least as many frame environments.

The dependence on the probe family has the opposite direction
under ordinary set inclusion.

\begin{proposition}[Probe-Family Extension]
  \label{ik.prop.probe-extension}
  Let \(P\) and \(Q\) be finite probe families
  for the same interpretation function,
  with
  \[
    P\subseteq Q.
  \]
  Then for every occurrence regime \(\rho\),
  \[
    K_Q(\rho)
    \sqsubseteq
    K_P(\rho).
  \]
\end{proposition}

\begin{proof}
  Suppose
  \[
    F\mathrel{K_Q(\rho)}F'.
  \]
  Then
  \[
    I(q,F)\mathrel{\rho}I(q,F')
  \]
  for every \(q\in Q\).

  Because \(P\subseteq Q\),
  the same condition holds for every \(p\in P\).
  Therefore
  \[
    F\mathrel{K_P(\rho)}F'.
  \]
\end{proof}

Adding probes can preserve a frame equivalence
or split an existing equivalence class.
It cannot merge frame environments
that the smaller probe family distinguished.

\begin{proposition}[Probe-Union Meet Law]
  \label{ik.prop.probe-union}
  For finite probe families \(P\) and \(Q\)
  and every occurrence regime \(\rho\),
  \[
    K_{P\cup Q}(\rho)
    =
    K_P(\rho)
    \wedge
    K_Q(\rho).
  \]
\end{proposition}

\begin{proof}
  For \(F,F'\in\mathcal{E}\),
  \[
    F\mathrel{K_{P\cup Q}(\rho)}F'
  \]
  holds when
  \[
    I(p,F)\mathrel{\rho}I(p,F')
  \]
  for every \(p\in P\cup Q\).

  This is equivalent to the conjunction of
  \[
    F\mathrel{K_P(\rho)}F'
  \]
  and
  \[
    F\mathrel{K_Q(\rho)}F'.
  \]

  Therefore
  \[
    K_{P\cup Q}(\rho)
    =
    K_P(\rho)\cap K_Q(\rho).
  \]
  Meet in \(\operatorname{Eq}(\mathcal{E})\)
  is intersection,
  so the result follows.
\end{proof}

\begin{corollary}[Kernel as a Meet of Single-Probe Kernels]
  \label{ik.cor.single-probe-meet}
  For every finite probe family \(P\)
  and occurrence regime \(\rho\),
  \[
    K_P(\rho)
    =
    \bigwedge_{p\in P}
    K_{\{p\}}(\rho).
  \]
\end{corollary}

\begin{corollary}[Boundary Kernels]
  \label{ik.cor.boundary-kernels}
  For every finite probe family \(P\),
  \[
    K_P(\top_{\mathcal{A}})
    =
    \top_{\mathcal{E}}.
  \]

  For the empty probe family,
  \[
    K_{\varnothing}(\rho)
    =
    \top_{\mathcal{E}}
  \]
  for every occurrence regime \(\rho\).
\end{corollary}

The first equality holds because
the indiscrete occurrence regime
identifies every pair of occurrence environments.
The second holds because
the universal condition over an empty probe family
is vacuously satisfied.

For equality on the occurrence carrier,
\[
  K_P(\bot_{\mathcal{A}})
\]
identifies those frame environments
that produce equal canonical occurrence environments
for every probe in \(P\).

The map
\[
  K_P:
  \operatorname{Eq}(\mathcal{A})
  \longrightarrow
  \operatorname{Eq}(\mathcal{E})
\]
is therefore a monotone map between finite partition lattices.
The next section constructs the map in the opposite direction
from a declared frame regime
and establishes the relation between the two constructions.


% =============================================================================
\section{The Adjoint Construction}
\label{ik.sec.adjoint}
% =============================================================================

The interpretive-kernel construction sends an occurrence regime
to the frame regime that it induces under a fixed probe family.
This section constructs the map in the opposite direction.

A declared frame regime requires certain pairs of frame environments
to remain indistinguishable.
For each probe,
that requirement determines pairs of occurrence environments
that must be treated as equivalent.
Taking the equivalence closure of those required pairs
yields the least occurrence regime
capable of supporting the declared frame regime.

\subsection{Generated Occurrence Regime}

Fix a finite comparison artifact
\[
  \bigl(
  P,
  \mathcal{A},
  \mathcal{E},
  I
  \bigr)
\]
and let
\[
  \varphi
  \in
  \operatorname{Eq}(\mathcal{E})
\]
be a frame regime.

\begin{definition}[Frame-Generated Output Relation]
  \label{ik.def.generated-relation}
  The \emph{frame-generated output relation}
  induced by \(P\) and \(\varphi\) is
  \[
    R_{P,\varphi}
    =
    \bigl\{
    \bigl(
    I(p,F),
    I(p,F')
    \bigr)
    \bigm|
    p\in P
    \text{ and }
    F\mathrel{\varphi}F'
    \bigr\}.
  \]
\end{definition}

The relation \(R_{P,\varphi}\)
contains every pair of occurrence environments
that the interpretation function produces
from frame environments
the declared frame regime requires to remain equivalent.

The relation need not itself be an equivalence relation.
Different frame pairs may produce overlapping output pairs,
and the required output equivalence may arise only
after symmetry and transitivity are imposed.

\begin{definition}[Generated Occurrence Regime]
  \label{ik.def.generated-regime}
  The \emph{generated occurrence regime}
  induced by the probe family \(P\)
  and frame regime \(\varphi\) is
  \[
    L_P(\varphi)
    =
    \operatorname{EqCl}_{\mathcal{A}}
    \bigl(
    R_{P,\varphi}
    \bigr).
  \]
\end{definition}

By construction,
\[
  L_P(\varphi)
  \in
  \operatorname{Eq}(\mathcal{A}).
\]

The map
\[
  L_P:
  \operatorname{Eq}(\mathcal{E})
  \longrightarrow
  \operatorname{Eq}(\mathcal{A})
\]
therefore sends a frame regime
to an occurrence regime.

\begin{proposition}[Least Supporting Occurrence Regime]
  \label{ik.prop.least-supporting}
  Let
  \[
    \varphi
    \in
    \operatorname{Eq}(\mathcal{E}).
  \]
  Then \(L_P(\varphi)\) has the following properties:
  \begin{enumerate}[label=(\alph*)]
    \item for every \(p\in P\) and \(F,F'\in\mathcal{E}\),
          \[
            F\mathrel{\varphi}F'
            \quad\Longrightarrow\quad
            I(p,F)
            \mathrel{L_P(\varphi)}
            I(p,F');
          \]
    \item if
          \[
            \rho
            \in
            \operatorname{Eq}(\mathcal{A})
          \]
          also satisfies
          \[
            F\mathrel{\varphi}F'
            \quad\Longrightarrow\quad
            I(p,F)\mathrel{\rho}I(p,F')
          \]
          for every \(p\in P\),
          then
          \[
            L_P(\varphi)
            \sqsubseteq
            \rho.
          \]
  \end{enumerate}
\end{proposition}

\begin{proof}
  Every pair
  \[
    \bigl(
    I(p,F),
    I(p,F')
    \bigr)
  \]
  with \(F\mathrel{\varphi}F'\)
  lies in \(R_{P,\varphi}\),
  and therefore lies in its equivalence closure
  \(L_P(\varphi)\).
  This proves the first statement.

  For the second,
  suppose \(\rho\) contains every such pair.
  Then
  \[
    R_{P,\varphi}
    \subseteq
    \rho.
  \]
  Because \(\rho\) is an equivalence relation,
  it contains the least equivalence relation
  generated by \(R_{P,\varphi}\).
  Hence
  \[
    L_P(\varphi)
    =
    \operatorname{EqCl}_{\mathcal{A}}
    \bigl(
    R_{P,\varphi}
    \bigr)
    \subseteq
    \rho.
  \]
\end{proof}

Thus \(L_P(\varphi)\)
does not introduce an arbitrary occurrence comparison.
It contains the equivalences forced
by the frame regime under the probes in \(P\),
together with those required by equivalence closure.

\begin{proposition}[Monotonicity of the Generated Regime]
  \label{ik.prop.generated-monotone}
  Let
  \[
    \varphi,\psi
    \in
    \operatorname{Eq}(\mathcal{E}).
  \]
  If
  \[
    \varphi
    \sqsubseteq
    \psi,
  \]
  then
  \[
    L_P(\varphi)
    \sqsubseteq
    L_P(\psi).
  \]
\end{proposition}

\begin{proof}
  If
  \[
    \varphi
    \subseteq
    \psi,
  \]
  then every pair used to construct
  \(R_{P,\varphi}\)
  is also used to construct
  \(R_{P,\psi}\).
  Therefore
  \[
    R_{P,\varphi}
    \subseteq
    R_{P,\psi}.
  \]
  Equivalence closure is monotone,
  so
  \[
    L_P(\varphi)
    \subseteq
    L_P(\psi).
  \]
\end{proof}

\subsection{Generating-Set Reduction}

Definition~\ref{ik.def.generated-relation}
ranges over every pair in the frame equivalence relation.
That full pair set is not required.

A generating relation for the frame regime
is sufficient to generate the same occurrence regime.

\begin{definition}[Probe Image of a Frame Relation]
  \label{ik.def.relation-image}
  For a relation
  \[
    G
    \subseteq
    \mathcal{E}\times\mathcal{E},
  \]
  define
  \[
    R_{P,G}
    =
    \bigl\{
    \bigl(
    I(p,F),
    I(p,F')
    \bigr)
    \bigm|
    p\in P
    \text{ and }
    (F,F')\in G
    \bigr\}.
  \]
\end{definition}

\begin{theorem}[Generating-Set Reduction]
  \label{ik.thm.generating-reduction}
  Let
  \[
    \varphi
    \in
    \operatorname{Eq}(\mathcal{E}),
  \]
  and let
  \[
    G
    \subseteq
    \mathcal{E}\times\mathcal{E}
  \]
  generate \(\varphi\), so that
  \[
    \operatorname{EqCl}_{\mathcal{E}}(G)
    =
    \varphi.
  \]
  Then
  \[
    \operatorname{EqCl}_{\mathcal{A}}
    \bigl(
    R_{P,G}
    \bigr)
    =
    L_P(\varphi).
  \]
\end{theorem}

\begin{proof}
  Because
  \[
    G
    \subseteq
    \varphi,
  \]
  Definition~\ref{ik.def.relation-image}
  gives
  \[
    R_{P,G}
    \subseteq
    R_{P,\varphi}.
  \]
  Therefore
  \[
    \operatorname{EqCl}_{\mathcal{A}}
    \bigl(
    R_{P,G}
    \bigr)
    \subseteq
    L_P(\varphi).
  \]

  For the reverse inclusion,
  take a generating pair
  \[
    \bigl(
    I(p,F),
    I(p,F')
    \bigr)
    \in
    R_{P,\varphi}.
  \]
  Then
  \[
    F\mathrel{\varphi}F'.
  \]

  Since
  \[
    \operatorname{EqCl}_{\mathcal{E}}(G)
    =
    \varphi,
  \]
  there exists a finite sequence
  \[
    F
    =
    F_0,
    F_1,
    \ldots,
    F_k
    =
    F'
  \]
  such that each consecutive pair
  is contained in \(G\)
  in one direction or the other.

  Fixing the same probe \(p\),
  each consecutive output pair
  \[
    \bigl(
    I(p,F_i),
    I(p,F_{i+1})
    \bigr)
  \]
  therefore lies in \(R_{P,G}\)
  in one direction or the other.

  Symmetry and transitivity of
  \[
    \operatorname{EqCl}_{\mathcal{A}}
    \bigl(
    R_{P,G}
    \bigr)
  \]
  give
  \[
    I(p,F)
    \mathrel{
      \operatorname{EqCl}_{\mathcal{A}}
      \left(
      R_{P,G}
      \right)
    }
    I(p,F').
  \]

  Every generator in \(R_{P,\varphi}\)
  therefore lies in
  \[
    \operatorname{EqCl}_{\mathcal{A}}
    \bigl(
    R_{P,G}
    \bigr).
  \]
  Taking equivalence closures yields
  \[
    L_P(\varphi)
    \subseteq
    \operatorname{EqCl}_{\mathcal{A}}
    \bigl(
    R_{P,G}
    \bigr).
  \]
\end{proof}

\begin{corollary}[Spanning-Forest Reduction]
  \label{ik.cor.spanning-forest}
  Let
  \[
    \varphi
    \in
    \operatorname{Eq}(\mathcal{E}).
  \]
  For each block of
  \[
    \mathcal{E}/\varphi,
  \]
  choose a spanning tree on the members of that block,
  and let \(G_{\varphi}\)
  be the union of the selected edge sets.

  Then
  \[
    \operatorname{EqCl}_{\mathcal{E}}
    \bigl(
    G_{\varphi}
    \bigr)
    =
    \varphi,
  \]
  and therefore
  \[
    L_P(\varphi)
    =
    \operatorname{EqCl}_{\mathcal{A}}
    \bigl(
    R_{P,G_{\varphi}}
    \bigr).
  \]

  The number of selected frame pairs is
  \[
    |\mathcal{E}|
    -
    |\mathcal{E}/\varphi|.
  \]
\end{corollary}

\begin{proof}
  A spanning tree connects every pair of vertices
  in its selected equivalence class,
  while no selected edge connects different classes.
  The equivalence closure of the forest
  is therefore \(\varphi\).

  A block containing \(n\) frame environments
  contributes \(n-1\) edges.
  Summing over the blocks gives
  \[
    \sum_{C\in\mathcal{E}/\varphi}
    \bigl(
    |C|-1
    \bigr)
    =
    |\mathcal{E}|
    -
    |\mathcal{E}/\varphi|.
  \]
  The result then follows from
  Theorem~\ref{ik.thm.generating-reduction}.
\end{proof}

The reduction replaces a possible quadratic enumeration
of all pairs within frame-equivalence classes
by a linear-size generating relation.
For each selected frame pair,
the interpretation need be compared only under the probes in \(P\).

\subsection{The Adjunction}

The maps
\[
  L_P:
  \operatorname{Eq}(\mathcal{E})
  \longrightarrow
  \operatorname{Eq}(\mathcal{A})
\]
and
\[
  K_P:
  \operatorname{Eq}(\mathcal{A})
  \longrightarrow
  \operatorname{Eq}(\mathcal{E})
\]
form an adjoint pair.

\begin{theorem}[Interpretive Adjunction]
  \label{ik.thm.adjunction}
  For every frame regime
  \[
    \varphi
    \in
    \operatorname{Eq}(\mathcal{E})
  \]
  and occurrence regime
  \[
    \rho
    \in
    \operatorname{Eq}(\mathcal{A}),
  \]
  \[
    L_P(\varphi)
    \sqsubseteq
    \rho
    \quad\Longleftrightarrow\quad
    \varphi
    \sqsubseteq
    K_P(\rho).
  \]
  Hence
  \[
    L_P
    \dashv
    K_P.
  \]
\end{theorem}

\begin{proof}
  Suppose first that
  \[
    L_P(\varphi)
    \sqsubseteq
    \rho.
  \]

  Let
  \[
    F\mathrel{\varphi}F'.
  \]
  By Proposition~\ref{ik.prop.least-supporting},
  \[
    I(p,F)
    \mathrel{L_P(\varphi)}
    I(p,F')
  \]
  for every \(p\in P\).

  Since
  \[
    L_P(\varphi)
    \subseteq
    \rho,
  \]
  it follows that
  \[
    I(p,F)
    \mathrel{\rho}
    I(p,F')
  \]
  for every \(p\in P\).

  Therefore
  \[
    F\mathrel{K_P(\rho)}F',
  \]
  and hence
  \[
    \varphi
    \sqsubseteq
    K_P(\rho).
  \]

  Conversely,
  suppose
  \[
    \varphi
    \sqsubseteq
    K_P(\rho).
  \]

  If
  \[
    F\mathrel{\varphi}F',
  \]
  then
  \[
    F\mathrel{K_P(\rho)}F'.
  \]
  By Definition~\ref{ik.def.kernel},
  \[
    I(p,F)
    \mathrel{\rho}
    I(p,F')
  \]
  for every \(p\in P\).

  Thus
  \[
    R_{P,\varphi}
    \subseteq
    \rho.
  \]
  Since \(\rho\) is an equivalence relation,
  it contains the equivalence closure of
  \(R_{P,\varphi}\).
  Therefore
  \[
    L_P(\varphi)
    \sqsubseteq
    \rho.
  \]
\end{proof}

The theorem states the same compatibility condition
in two equivalent ways.

The left-hand statement asks whether
the occurrence regime \(\rho\)
contains every output equivalence
required by the frame regime \(\varphi\).

The right-hand statement asks whether
the interpretive kernel induced by \(\rho\)
contains every frame equivalence
declared by \(\varphi\).

\begin{corollary}[Unit and Counit]
  \label{ik.cor.unit-counit}
  For every frame regime
  \[
    \varphi
    \in
    \operatorname{Eq}(\mathcal{E}),
  \]
  \[
    \varphi
    \sqsubseteq
    K_P
    \bigl(
    L_P(\varphi)
    \bigr).
  \]

  For every occurrence regime
  \[
    \rho
    \in
    \operatorname{Eq}(\mathcal{A}),
  \]
  \[
    L_P
    \bigl(
    K_P(\rho)
    \bigr)
    \sqsubseteq
    \rho.
  \]
\end{corollary}

\begin{proof}
  Apply Theorem~\ref{ik.thm.adjunction}
  first with
  \[
    \rho=L_P(\varphi),
  \]
  using reflexivity of the refinement order,
  and then with
  \[
    \varphi=K_P(\rho).
  \]
\end{proof}

The first inequality says that
the least occurrence regime generated by a frame regime
supports at least that frame regime
when transported back through the kernel.

The second says that
regenerating an occurrence regime
from its induced frame kernel
introduces no occurrence equivalence
that was absent from the original regime.

\begin{remark}[Equality Need Not Hold]
  \label{ik.rem.unit-counit-strict}
  Neither inequality in
  Corollary~\ref{ik.cor.unit-counit}
  is asserted to be an equality.

  The probe family may be unable to observe
  every distinction carried by a frame regime,
  so
  \[
    \varphi
    \subsetneq
    K_P
    \bigl(
    L_P(\varphi)
    \bigr)
  \]
  may hold.

  An occurrence regime may also contain distinctions
  unrelated to the outputs generated by \(P\),
  so
  \[
    L_P
    \bigl(
    K_P(\rho)
    \bigr)
    \subsetneq
    \rho
  \]
  may hold.
\end{remark}

\subsection{Preservation Laws}

As a right adjoint,
the interpretive-kernel map preserves meets.
As a left adjoint,
the generated-regime map preserves joins.

\begin{theorem}[Meet Preservation]
  \label{ik.thm.kernel-meets}
  Let
  \[
    \{
    \rho_i
    \}_{i\in J}
    \subseteq
    \operatorname{Eq}(\mathcal{A})
  \]
  be a family of occurrence regimes.
  Then
  \[
    K_P
    \left(
    \bigwedge_{i\in J}
    \rho_i
    \right)
    =
    \bigwedge_{i\in J}
    K_P(\rho_i).
  \]
\end{theorem}

\begin{proof}
  For \(F,F'\in\mathcal{E}\),
  \[
    F
    \mathrel{
      K_P
      \left(
      \bigwedge_{i\in J}
      \rho_i
      \right)
    }
    F'
  \]
  holds when
  \[
    I(p,F)
    \mathrel{
      \bigwedge_{i\in J}
      \rho_i
    }
    I(p,F')
  \]
  for every \(p\in P\).

  Since meet in
  \(\operatorname{Eq}(\mathcal{A})\)
  is intersection,
  this holds when
  \[
    I(p,F)
    \mathrel{\rho_i}
    I(p,F')
  \]
  for every \(p\in P\) and every \(i\in J\).

  That condition is equivalent to
  \[
    F
    \mathrel{K_P(\rho_i)}
    F'
  \]
  for every \(i\in J\),
  which is
  \[
    F
    \mathrel{
      \bigwedge_{i\in J}
      K_P(\rho_i)
    }
    F'.
  \]
\end{proof}

\begin{theorem}[Join Preservation]
  \label{ik.thm.generated-joins}
  Let
  \[
    \{
    \varphi_i
    \}_{i\in J}
    \subseteq
    \operatorname{Eq}(\mathcal{E})
  \]
  be a family of frame regimes.
  Then
  \[
    L_P
    \left(
    \bigvee_{i\in J}
    \varphi_i
    \right)
    =
    \bigvee_{i\in J}
    L_P(\varphi_i).
  \]
\end{theorem}

\begin{proof}
  Let
  \[
    \rho
    \in
    \operatorname{Eq}(\mathcal{A}).
  \]
  By Theorem~\ref{ik.thm.adjunction},
  \[
    L_P
    \left(
    \bigvee_{i\in J}
    \varphi_i
    \right)
    \sqsubseteq
    \rho
  \]
  holds when
  \[
    \bigvee_{i\in J}
    \varphi_i
    \sqsubseteq
    K_P(\rho).
  \]

  By the defining property of join,
  this holds when
  \[
    \varphi_i
    \sqsubseteq
    K_P(\rho)
  \]
  for every \(i\in J\).

  Applying the adjunction to each \(i\),
  this is equivalent to
  \[
    L_P(\varphi_i)
    \sqsubseteq
    \rho
  \]
  for every \(i\in J\).

  By the defining property of join,
  the latter holds when
  \[
    \bigvee_{i\in J}
    L_P(\varphi_i)
    \sqsubseteq
    \rho.
  \]

  The two left-hand elements therefore have
  the same upper bounds,
  and hence are equal.
\end{proof}

\begin{corollary}[Boundary Preservation]
  \label{ik.cor.boundary-preservation}
  \[
    K_P
    \bigl(
    \top_{\mathcal{A}}
    \bigr)
    =
    \top_{\mathcal{E}},
  \]
  and
  \[
    L_P
    \bigl(
    \bot_{\mathcal{E}}
    \bigr)
    =
    \bot_{\mathcal{A}}.
  \]
\end{corollary}

The preservation laws are directional.
The adjunction guarantees meet preservation by \(K_P\)
and join preservation by \(L_P\).
It does not by itself establish
join preservation by \(K_P\)
or meet preservation by \(L_P\).

\subsection{Probe-Family Duality}

Changing the probe family
acts dually on the two constructions.

Adding probes makes the interpretive kernel finer,
because more observations must agree.
The same addition makes the generated occurrence regime coarser,
because more output-equivalence requirements are imposed.

\begin{theorem}[Probe-Union Duality]
  \label{ik.thm.probe-union-duality}
  Let \(P\) and \(Q\)
  be finite probe families
  for the same occurrence carrier,
  frame carrier,
  and interpretation function.

  For every occurrence regime
  \[
    \rho
    \in
    \operatorname{Eq}(\mathcal{A}),
  \]
  \[
    K_{P\cup Q}(\rho)
    =
    K_P(\rho)
    \wedge
    K_Q(\rho).
  \]

  For every frame regime
  \[
    \varphi
    \in
    \operatorname{Eq}(\mathcal{E}),
  \]
  \[
    L_{P\cup Q}(\varphi)
    =
    L_P(\varphi)
    \vee
    L_Q(\varphi).
  \]
\end{theorem}

\begin{proof}
  The first equality is
  Proposition~\ref{ik.prop.probe-union}.

  For the second,
  Definition~\ref{ik.def.generated-relation}
  gives
  \[
    R_{P\cup Q,\varphi}
    =
    R_{P,\varphi}
    \cup
    R_{Q,\varphi}.
  \]

  Therefore
  \[
    L_{P\cup Q}(\varphi)
    =
    \operatorname{EqCl}_{\mathcal{A}}
    \bigl(
    R_{P,\varphi}
    \cup
    R_{Q,\varphi}
    \bigr).
  \]

  Join in
  \(\operatorname{Eq}(\mathcal{A})\)
  is the equivalence closure of union,
  so
  \[
    L_{P\cup Q}(\varphi)
    =
    L_P(\varphi)
    \vee
    L_Q(\varphi).
  \]
\end{proof}

\begin{corollary}[Finite Probe-Family Decomposition]
  \label{ik.cor.probe-decomposition}
  Let
  \[
    \{
    P_i
    \}_{i\in J}
  \]
  be a finite family of finite probe families.

  For every occurrence regime \(\rho\),
  \[
    K_{\bigcup_{i\in J}P_i}(\rho)
    =
    \bigwedge_{i\in J}
    K_{P_i}(\rho).
  \]

  For every frame regime \(\varphi\),
  \[
    L_{\bigcup_{i\in J}P_i}(\varphi)
    =
    \bigvee_{i\in J}
    L_{P_i}(\varphi).
  \]
\end{corollary}

\begin{corollary}[Empty Probe Family]
  \label{ik.cor.empty-probe-adjoint}
  For every occurrence regime \(\rho\),
  \[
    K_{\varnothing}(\rho)
    =
    \top_{\mathcal{E}}.
  \]

  For every frame regime \(\varphi\),
  \[
    L_{\varnothing}(\varphi)
    =
    \bot_{\mathcal{A}}.
  \]
\end{corollary}

The empty probe family distinguishes no frame environments
and imposes no nontrivial output-equivalence requirements.

The adjunction and its preservation laws
provide the algebraic basis for the threshold construction.
A declared frame regime \(\varphi\)
is supported by an occurrence regime \(\rho\) when
\[
  L_P(\varphi)
  \sqsubseteq
  \rho,
\]
equivalently when
\[
  \varphi
  \sqsubseteq
  K_P(\rho).
\]
The next section studies the boundary
between occurrence regimes that satisfy this condition
and those that do not.

% =============================================================================
\section{Thresholds}
\label{ik.sec.thresholds}
% =============================================================================

The adjunction characterizes when an occurrence regime
is sufficiently coarse
to support a declared frame regime.

That characterization has the form of a threshold.
For each probe family and frame regime,
there is a least occurrence regime
that identifies every output pair
whose equivalence is required by the frame regime.
Every coarser occurrence regime is also sufficient,
while every regime not lying above that least element is insufficient.

The threshold is therefore lattice-valued.
It is an equivalence relation on the occurrence carrier,
not a scalar score or a position in a total order.

\subsection{Admissible Occurrence Regimes}

Fix a finite comparison artifact,
a probe family \(P\),
and a frame regime
\[
  \varphi
  \in
  \operatorname{Eq}(\mathcal{E}).
\]

\begin{definition}[Admissible Occurrence Regime]
  \label{ik.def.admissible}
  An occurrence regime
  \[
    \rho
    \in
    \operatorname{Eq}(\mathcal{A})
  \]
  is \emph{admissible for \(\varphi\) under \(P\)}
  when
  \[
    \varphi
    \sqsubseteq
    K_P(\rho).
  \]

  Write
  \[
    \operatorname{Adm}_P(\varphi)
    =
    \bigl\{
    \rho
    \in
    \operatorname{Eq}(\mathcal{A})
    \bigm|
    \varphi
    \sqsubseteq
    K_P(\rho)
    \bigr\}
  \]
  for the set of admissible occurrence regimes.
\end{definition}

Admissibility means that every pair of frame environments
identified by \(\varphi\)
also lies in the interpretive kernel induced by \(\rho\).
Equivalently,
for every \(F,F'\in\mathcal{E}\),
\[
  F\mathrel{\varphi}F'
\]
implies
\[
  I(p,F)\mathrel{\rho}I(p,F')
\]
for every \(p\in P\).

\begin{proposition}[Adjoint Characterization of Admissibility]
  \label{ik.prop.admissible-adjoint}
  For every occurrence regime
  \[
    \rho
    \in
    \operatorname{Eq}(\mathcal{A}),
  \]
  \[
    \rho
    \in
    \operatorname{Adm}_P(\varphi)
  \]
  if and only if
  \[
    L_P(\varphi)
    \sqsubseteq
    \rho.
  \]
\end{proposition}

\begin{proof}
  By Definition~\ref{ik.def.admissible},
  \(\rho\) is admissible when
  \[
    \varphi
    \sqsubseteq
    K_P(\rho).
  \]

  By Theorem~\ref{ik.thm.adjunction},
  this holds when
  \[
    L_P(\varphi)
    \sqsubseteq
    \rho.
  \]
\end{proof}

Thus
\[
  \operatorname{Adm}_P(\varphi)
  =
  \bigl\{
  \rho
  \in
  \operatorname{Eq}(\mathcal{A})
  \bigm|
  L_P(\varphi)
  \sqsubseteq
  \rho
  \bigr\}.
\]

The admissible occurrence regimes form
the principal upper set generated by \(L_P(\varphi)\).
Admissibility is therefore not an arbitrary predicate
over the occurrence-regime lattice.

\begin{proposition}[Upward Closure of Admissibility]
  \label{ik.prop.admissible-upward}
  Let
  \[
    \rho,\sigma
    \in
    \operatorname{Eq}(\mathcal{A}).
  \]
  If
  \[
    \rho
    \in
    \operatorname{Adm}_P(\varphi)
    \qquad\text{and}\qquad
    \rho
    \sqsubseteq
    \sigma,
  \]
  then
  \[
    \sigma
    \in
    \operatorname{Adm}_P(\varphi).
  \]
\end{proposition}

\begin{proof}
  Proposition~\ref{ik.prop.admissible-adjoint}
  gives
  \[
    L_P(\varphi)
    \sqsubseteq
    \rho.
  \]
  Combining this with
  \[
    \rho
    \sqsubseteq
    \sigma
  \]
  and transitivity yields
  \[
    L_P(\varphi)
    \sqsubseteq
    \sigma.
  \]
  The same proposition then gives
  \[
    \sigma
    \in
    \operatorname{Adm}_P(\varphi).
  \]
\end{proof}

Coarsening an admissible occurrence regime
cannot invalidate the frame equivalences it already supports.
It may identify additional occurrence environments
and thereby induce additional frame equivalences,
but it cannot remove an equivalence required by \(\varphi\).

\subsection{Lattice-Valued Thresholds}

\begin{definition}[Interpretive Threshold]
  \label{ik.def.threshold}
  The \emph{interpretive threshold}
  of a frame regime \(\varphi\)
  under probe family \(P\) is
  \[
    \Theta_P(\varphi)
    =
    L_P(\varphi).
  \]
\end{definition}

The threshold is an element of
\[
  \operatorname{Eq}(\mathcal{A}).
\]
It is therefore a partition of the occurrence carrier
rather than a numerical value.

\begin{theorem}[Exact Threshold Characterization]
  \label{ik.thm.threshold-characterization}
  For every frame regime \(\varphi\)
  and occurrence regime \(\rho\),
  \[
    \rho
    \in
    \operatorname{Adm}_P(\varphi)
    \quad\Longleftrightarrow\quad
    \Theta_P(\varphi)
    \sqsubseteq
    \rho.
  \]
\end{theorem}

\begin{proof}
  Substitute
  \[
    \Theta_P(\varphi)
    =
    L_P(\varphi)
  \]
  into Proposition~\ref{ik.prop.admissible-adjoint}.
\end{proof}

The theorem gives the threshold its exact meaning.

An occurrence regime is sufficient
precisely when it contains
every occurrence equivalence generated
by the frame regime under the probes in \(P\).

A regime that is not above the threshold
is not admissible.
It either omits a required threshold pair
or is incomparable with the threshold
and therefore still omits at least one required equivalence.

\begin{corollary}[Threshold Is Admissible]
  \label{ik.cor.threshold-admissible}
  For every frame regime \(\varphi\),
  \[
    \Theta_P(\varphi)
    \in
    \operatorname{Adm}_P(\varphi).
  \]
\end{corollary}

\begin{proof}
  Reflexivity gives
  \[
    \Theta_P(\varphi)
    \sqsubseteq
    \Theta_P(\varphi).
  \]
  Apply Theorem~\ref{ik.thm.threshold-characterization}.
\end{proof}

\begin{corollary}[Threshold Is the Finest Admissible Regime]
  \label{ik.cor.threshold-finest}
  For every
  \[
    \rho
    \in
    \operatorname{Adm}_P(\varphi),
  \]
  \[
    \Theta_P(\varphi)
    \sqsubseteq
    \rho.
  \]
  Hence \(\Theta_P(\varphi)\)
  is the least element of
  \(\operatorname{Adm}_P(\varphi)\)
  under the refinement order
  and the finest admissible occurrence regime.
\end{corollary}

The threshold introduces no occurrence equivalence
beyond those forced by the declared frame regime
and the selected probes.
Any strictly coarser admissible regime
contains additional occurrence equivalences
not required by that demand.

\begin{proposition}[Monotonicity of Thresholds]
  \label{ik.prop.threshold-monotone}
  Let
  \[
    \varphi,\psi
    \in
    \operatorname{Eq}(\mathcal{E}).
  \]
  If
  \[
    \varphi
    \sqsubseteq
    \psi,
  \]
  then
  \[
    \Theta_P(\varphi)
    \sqsubseteq
    \Theta_P(\psi).
  \]
\end{proposition}

\begin{proof}
  By Definition~\ref{ik.def.threshold},
  \[
    \Theta_P(\varphi)
    =
    L_P(\varphi)
  \]
  and
  \[
    \Theta_P(\psi)
    =
    L_P(\psi).
  \]
  The result follows from
  Proposition~\ref{ik.prop.generated-monotone}.
\end{proof}

A coarser frame regime identifies more frame pairs
and therefore makes a stronger invariance demand.
Its threshold contains at least
the occurrence equivalences required
by every finer frame regime below it.

\begin{proposition}[Probe Extension Raises the Threshold]
  \label{ik.prop.threshold-probes}
  Let
  \[
    P
    \subseteq
    Q
  \]
  be finite probe families.
  Then for every frame regime \(\varphi\),
  \[
    \Theta_P(\varphi)
    \sqsubseteq
    \Theta_Q(\varphi).
  \]
\end{proposition}

\begin{proof}
  Since
  \[
    Q
    =
    P\cup Q,
  \]
  Theorem~\ref{ik.thm.probe-union-duality}
  gives
  \[
    \Theta_Q(\varphi)
    =
    \Theta_P(\varphi)
    \vee
    \Theta_Q(\varphi).
  \]
  Hence
  \[
    \Theta_P(\varphi)
    \sqsubseteq
    \Theta_Q(\varphi).
  \]
\end{proof}

Adding probes may reveal additional occurrence equivalences
that are required to preserve the same frame regime.
It can leave the threshold unchanged
or move it upward to a coarser partition.
It cannot make the threshold finer.

\begin{remark}[Why the Threshold Is Not Scalar]
  \label{ik.rem.threshold-not-scalar}
  The occurrence-regime lattice is generally only partially ordered.

  Two thresholds may require different occurrence equivalences
  without either containing the other.
  Assigning them scalar levels would impose
  a linear comparison not supplied by the mathematics.

  Threshold reporting must therefore retain
  the partition itself,
  or an equivalent finite representation of it,
  together with the probe family
  and frame regime that generated it.
\end{remark}

\subsection{Incomparable Thresholds}

Distinct frame demands may generate incomparable thresholds.

\begin{proposition}[Incomparable Thresholds Exist]
  \label{ik.prop.incomparable-thresholds}
  There exist a finite comparison artifact
  and frame regimes
  \[
    \varphi,\psi
    \in
    \operatorname{Eq}(\mathcal{E})
  \]
  such that
  \[
    \Theta_P(\varphi)
  \]
  and
  \[
    \Theta_P(\psi)
  \]
  are incomparable under refinement.
\end{proposition}

\begin{proof}
  Let
  \[
    \mathcal{E}
    =
    \{F_0,F_1,F_2\},
    \qquad
    \mathcal{A}
    =
    \{a,b,c\},
    \qquad
    P
    =
    \{p\}.
  \]

  Define the interpretation by
  \[
    I(p,F_0)=a,
    \qquad
    I(p,F_1)=b,
    \qquad
    I(p,F_2)=c.
  \]

  Let \(\varphi\) induce the frame partition
  \[
    \mathcal{E}/\varphi
    =
    \bigl\{
    \{F_0,F_1\},
    \{F_2\}
    \bigr\},
  \]
  and let \(\psi\) induce
  \[
    \mathcal{E}/\psi
    =
    \bigl\{
    \{F_0,F_2\},
    \{F_1\}
    \bigr\}.
  \]

  The threshold generated by \(\varphi\)
  has occurrence partition
  \[
    \mathcal{A}/\Theta_P(\varphi)
    =
    \bigl\{
    \{a,b\},
    \{c\}
    \bigr\}.
  \]

  The threshold generated by \(\psi\)
  has occurrence partition
  \[
    \mathcal{A}/\Theta_P(\psi)
    =
    \bigl\{
    \{a,c\},
    \{b\}
    \bigr\}.
  \]

  The first threshold relates \(a\) and \(b\)
  but not \(a\) and \(c\).
  The second relates \(a\) and \(c\)
  but not \(a\) and \(b\).

  Neither relation contains the other,
  so the thresholds are incomparable.
\end{proof}

The example exhibits two different occurrence demands.
One requires the outputs \(a\) and \(b\) to be identified.
The other requires \(a\) and \(c\) to be identified.
Neither demand is uniformly stronger than the other.

When both frame demands must be supported,
their thresholds combine by join.

\begin{theorem}[Threshold of Combined Frame Demands]
  \label{ik.thm.combined-threshold}
  Let
  \[
    \{
    \varphi_i
    \}_{i\in J}
  \]
  be a finite family of frame regimes.
  Then
  \[
    \Theta_P
    \left(
    \bigvee_{i\in J}
    \varphi_i
    \right)
    =
    \bigvee_{i\in J}
    \Theta_P(\varphi_i).
  \]
\end{theorem}

\begin{proof}
  By Definition~\ref{ik.def.threshold},
  the threshold map is \(L_P\).
  The result is therefore
  Theorem~\ref{ik.thm.generated-joins}.
\end{proof}

\begin{corollary}[Simultaneous Admissibility]
  \label{ik.cor.simultaneous-admissibility}
  An occurrence regime \(\rho\)
  is admissible for every
  \(\varphi_i\)
  if and only if
  \[
    \bigvee_{i\in J}
    \Theta_P(\varphi_i)
    \sqsubseteq
    \rho.
  \]
\end{corollary}

\begin{proof}
  The regime \(\rho\)
  is admissible for every \(\varphi_i\)
  when
  \[
    \Theta_P(\varphi_i)
    \sqsubseteq
    \rho
  \]
  for every \(i\in J\).

  By the defining property of join,
  this holds when
  \[
    \bigvee_{i\in J}
    \Theta_P(\varphi_i)
    \sqsubseteq
    \rho.
  \]
\end{proof}

In the construction of
Proposition~\ref{ik.prop.incomparable-thresholds},
the join of the two thresholds is
\[
  \top_{\mathcal{A}},
\]
because the required pairs
connect \(a\) to both \(b\) and \(c\).
Supporting both frame demands
therefore requires all three occurrence environments
to be equivalent.

The individual demands are not linearly ordered,
but their joint threshold remains well defined.

\subsection{Extremal Admissible Regimes}

The refinement order reverses
a possible informal reading of the words
\emph{least} and \emph{greatest}.

The least element of an admissible set
is its finest regime.
The greatest element is its coarsest regime.

\begin{theorem}[Extremal Admissible Regimes]
  \label{ik.thm.extremal-admissible}
  For every frame regime \(\varphi\),
  the admissible set
  \[
    \operatorname{Adm}_P(\varphi)
  \]
  is nonempty and has:
  \begin{enumerate}[label=(\alph*)]
    \item the least element
          \[
            \Theta_P(\varphi);
          \]
    \item the greatest element
          \[
            \top_{\mathcal{A}}.
          \]
  \end{enumerate}
\end{theorem}

\begin{proof}
  Corollary~\ref{ik.cor.threshold-admissible}
  shows that
  \[
    \Theta_P(\varphi)
  \]
  is admissible.

  Corollary~\ref{ik.cor.threshold-finest}
  shows that it refines every admissible regime,
  so it is the least element.

  By Corollary~\ref{ik.cor.boundary-kernels},
  \[
    K_P(\top_{\mathcal{A}})
    =
    \top_{\mathcal{E}}.
  \]
  Every frame regime satisfies
  \[
    \varphi
    \sqsubseteq
    \top_{\mathcal{E}},
  \]
  so
  \[
    \top_{\mathcal{A}}
    \in
    \operatorname{Adm}_P(\varphi).
  \]

  Since
  \[
    \top_{\mathcal{A}}
  \]
  is the greatest occurrence regime,
  it is the greatest admissible regime.
\end{proof}

The greatest admissible regime
is therefore always available,
but carries little diagnostic information.
It identifies every pair of occurrence environments
and consequently supports every frame regime.

The mathematically informative boundary
is the least admissible regime,
because it is the finest occurrence comparison
that still supports the declared frame demand.

\begin{proposition}[Meet Closure of Admissibility]
  \label{ik.prop.admissible-meets}
  Let
  \[
    \{
    \rho_i
    \}_{i\in J}
    \subseteq
    \operatorname{Adm}_P(\varphi)
  \]
  be a nonempty family of admissible occurrence regimes.
  Then
  \[
    \bigwedge_{i\in J}
    \rho_i
    \in
    \operatorname{Adm}_P(\varphi).
  \]
\end{proposition}

\begin{proof}
  Since each \(\rho_i\) is admissible,
  Theorem~\ref{ik.thm.threshold-characterization}
  gives
  \[
    \Theta_P(\varphi)
    \sqsubseteq
    \rho_i
  \]
  for every \(i\in J\).

  Thus \(\Theta_P(\varphi)\)
  is a lower bound of the family.
  By the defining property of meet,
  \[
    \Theta_P(\varphi)
    \sqsubseteq
    \bigwedge_{i\in J}
    \rho_i.
  \]

  Applying the threshold characterization again
  proves that the meet is admissible.
\end{proof}

\begin{corollary}[Threshold as the Meet of All Admissible Regimes]
  \label{ik.cor.threshold-meet}
  For every frame regime \(\varphi\),
  \[
    \Theta_P(\varphi)
    =
    \bigwedge
    \operatorname{Adm}_P(\varphi).
  \]
\end{corollary}

\begin{proof}
  The threshold belongs to the admissible set
  and refines every member of it.
  It is therefore the greatest lower bound
  of the admissible regimes.
\end{proof}

\begin{remark}[The Existence Question Is Resolved]
  \label{ik.rem.threshold-existence}
  Before the adjoint construction is available,
  one may ask whether the admissible occurrence regimes
  are closed under meet
  and therefore possess a finest member.

  The adjunction gives the stronger result.
  The admissible set is not merely meet-closed.
  It is the principal upper set
  generated by
  \[
    L_P(\varphi).
  \]

  No additional closure hypothesis is required.
  The finest admissible occurrence regime
  exists for every frame regime
  in the finite comparison artifact
  and is
  \[
    \Theta_P(\varphi)
    =
    L_P(\varphi).
  \]
\end{remark}

Thresholds therefore summarize
the exact occurrence-equivalence demand
made by a frame regime under a fixed probe family.

Their values may be equal,
strictly ordered,
or incomparable.
Multiple demands combine by join,
and every admissible occurrence regime
lies in the principal upper set
above their combined threshold.



% =============================================================================
\section{Representation}
\label{ik.sec.representation}
% =============================================================================

The threshold construction assigns an occurrence regime
to an entire frame regime.
The same construction can be applied to the least frame regime
that identifies one selected pair of frame environments.

The resulting occurrence regime measures
the exact output equivalences required
to make that pair interpretively indistinguishable.
These pairwise values determine the complete interpretive-kernel map.

\subsection{Pairwise Interpretive Demand}

\begin{definition}[Principal Pair Regime]
  \label{ik.def.principal-pair}
  For frame environments
  \(F,F'\in\mathcal{E}\),
  define
  \[
    \gamma_{F,F'}
    =
    \operatorname{EqCl}_{\mathcal{E}}
    \bigl(
    \{
    (F,F')
    \}
    \bigr).
  \]

  The relation \(\gamma_{F,F'}\)
  is the least frame regime
  that identifies \(F\) and \(F'\).
\end{definition}

When \(F\neq F'\),
the partition induced by \(\gamma_{F,F'}\)
has one block \(\{F,F'\}\)
and a singleton block for every other frame environment.
When \(F=F'\),
\[
  \gamma_{F,F}
  =
  \bot_{\mathcal{E}}.
\]

\begin{definition}[Pairwise Interpretive Demand]
  \label{ik.def.pairwise-demand}
  The \emph{pairwise interpretive demand}
  induced by the probe family \(P\) is the function
  \[
    \delta_P:
    \mathcal{E}\times\mathcal{E}
    \longrightarrow
    \operatorname{Eq}(\mathcal{A})
  \]
  defined by
  \[
    \delta_P(F,F')
    =
    L_P
    \bigl(
    \gamma_{F,F'}
    \bigr).
  \]
\end{definition}

The value \(\delta_P(F,F')\)
is the least occurrence regime
under which the probes in \(P\)
make \(F\) and \(F'\) interpretively equivalent.

\begin{proposition}[Explicit Pairwise Construction]
  \label{ik.prop.pairwise-explicit}
  For every \(F,F'\in\mathcal{E}\),
  \[
    \delta_P(F,F')
    =
    \operatorname{EqCl}_{\mathcal{A}}
    \Bigl(
    \bigl\{
    \bigl(
    I(p,F),
    I(p,F')
    \bigr)
    \bigm|
    p\in P
    \bigr\}
    \Bigr).
  \]
\end{proposition}

\begin{proof}
  By Definition~\ref{ik.def.pairwise-demand},
  \[
    \delta_P(F,F')
    =
    L_P
    \bigl(
    \gamma_{F,F'}
    \bigr).
  \]

  The principal pair regime
  is generated by the single pair
  \((F,F')\).
  Applying Theorem~\ref{ik.thm.generating-reduction}
  to that generating relation gives
  \[
    L_P
    \bigl(
    \gamma_{F,F'}
    \bigr)
    =
    \operatorname{EqCl}_{\mathcal{A}}
    \Bigl(
    \bigl\{
    \bigl(
    I(p,F),
    I(p,F')
    \bigr)
    \bigm|
    p\in P
    \bigr\}
    \Bigr).
  \]
\end{proof}

Thus the pairwise demand may be computed directly.
For each probe,
the two interpreted outputs contribute one generating pair,
and equivalence closure produces
the least occurrence regime containing all of them.

\begin{theorem}[Pairwise Threshold Characterization]
  \label{ik.thm.pairwise-threshold}
  For every occurrence regime
  \[
    \rho
    \in
    \operatorname{Eq}(\mathcal{A})
  \]
  and frame environments
  \(F,F'\in\mathcal{E}\),
  \[
    F\mathrel{K_P(\rho)}F'
    \quad\Longleftrightarrow\quad
    \delta_P(F,F')
    \sqsubseteq
    \rho.
  \]
\end{theorem}

\begin{proof}
  Because \(K_P(\rho)\)
  is an equivalence relation,
  \[
    F\mathrel{K_P(\rho)}F'
  \]
  holds when
  \[
    \gamma_{F,F'}
    \sqsubseteq
    K_P(\rho).
  \]

  By Theorem~\ref{ik.thm.adjunction},
  this is equivalent to
  \[
    L_P
    \bigl(
    \gamma_{F,F'}
    \bigr)
    \sqsubseteq
    \rho.
  \]

  Definition~\ref{ik.def.pairwise-demand}
  identifies the left-hand side as
  \(\delta_P(F,F')\).
\end{proof}

The complete kernel can therefore be recovered
from the pairwise demands:
\[
  K_P(\rho)
  =
  \bigl\{
  (F,F')
  \in
  \mathcal{E}\times\mathcal{E}
  \bigm|
  \delta_P(F,F')
  \sqsubseteq
  \rho
  \bigr\}.
\]

\subsection{Generalized Ultrametric Laws}

The codomain of \(\delta_P\)
is the partition lattice
\[
  \operatorname{Eq}(\mathcal{A}).
\]
The lattice join plays the role
taken by maximum in an ordinary ultrametric.

\begin{definition}[Equivalence-Valued Generalized Ultrametric]
  \label{ik.def.generalized-ultrametric}
  Let \(X\) and \(A\) be sets.
  An
  \emph{\(\operatorname{Eq}(A)\)-valued generalized ultrametric}
  on \(X\)
  is a function
  \[
    \delta:
    X\times X
    \longrightarrow
    \operatorname{Eq}(A)
  \]
  satisfying:
  \begin{enumerate}[label=(U\arabic*)]
    \item \emph{reflexive zero:}
          \[
            \delta(x,x)
            =
            \bot_A;
          \]
    \item \emph{symmetry:}
          \[
            \delta(x,y)
            =
            \delta(y,x);
          \]
    \item \emph{strong triangle inequality:}
          \[
            \delta(x,z)
            \sqsubseteq
            \delta(x,y)
            \vee
            \delta(y,z).
          \]
  \end{enumerate}
\end{definition}

The term \emph{zero}
refers to the least equivalence relation,
not to a numerical value.

The strong triangle inequality says that
every occurrence equivalence required
to identify \(x\) with \(z\)
is contained in the regime generated jointly
by the demands for \(x\) with \(y\)
and \(y\) with \(z\).

\begin{theorem}[Interpretive Demand Is a Generalized Ultrametric]
  \label{ik.thm.demand-ultrametric}
  The pairwise interpretive demand
  \[
    \delta_P:
    \mathcal{E}\times\mathcal{E}
    \longrightarrow
    \operatorname{Eq}(\mathcal{A})
  \]
  is an
  \(\operatorname{Eq}(\mathcal{A})\)-valued
  generalized ultrametric.
\end{theorem}

\begin{proof}
  For reflexive zero,
  \[
    \gamma_{F,F}
    =
    \bot_{\mathcal{E}}.
  \]
  By Corollary~\ref{ik.cor.boundary-preservation},
  \[
    \delta_P(F,F)
    =
    L_P
    \bigl(
    \bot_{\mathcal{E}}
    \bigr)
    =
    \bot_{\mathcal{A}}.
  \]

  For symmetry,
  \[
    \gamma_{F,F'}
    =
    \gamma_{F',F},
  \]
  so
  \[
    \delta_P(F,F')
    =
    \delta_P(F',F).
  \]

  For the strong triangle inequality,
  let
  \[
    \rho
    =
    \delta_P(F,F')
    \vee
    \delta_P(F',F'').
  \]

  Since
  \[
    \delta_P(F,F')
    \sqsubseteq
    \rho,
  \]
  Theorem~\ref{ik.thm.pairwise-threshold}
  gives
  \[
    F\mathrel{K_P(\rho)}F'.
  \]

  Similarly,
  \[
    F'\mathrel{K_P(\rho)}F''.
  \]

  The relation \(K_P(\rho)\)
  is transitive by
  Proposition~\ref{ik.prop.kernel-equivalence},
  so
  \[
    F\mathrel{K_P(\rho)}F''.
  \]

  Applying
  Theorem~\ref{ik.thm.pairwise-threshold}
  again gives
  \[
    \delta_P(F,F'')
    \sqsubseteq
    \rho.
  \]

  Substituting the definition of \(\rho\)
  yields
  \[
    \delta_P(F,F'')
    \sqsubseteq
    \delta_P(F,F')
    \vee
    \delta_P(F',F'').
  \]
\end{proof}

\begin{remark}[Distinct Frames May Have Zero Demand]
  \label{ik.rem.zero-demand}
  The generalized ultrametric need not separate
  distinct frame environments.

  For \(F\neq F'\),
  \[
    \delta_P(F,F')
    =
    \bot_{\mathcal{A}}
  \]
  when
  \[
    I(p,F)
    =
    I(p,F')
  \]
  for every \(p\in P\).

  Distinct frame representations may therefore
  have zero pairwise demand
  when the selected probes produce
  identical canonical occurrence environments.
\end{remark}

The construction is consequently
a generalized pseudoultrametric
when the word \emph{ultrametric}
is reserved for separating distances.
No separation hypothesis is needed
for the representation result below.

\subsection{Representation Theorem}

The preceding construction is not specific
to interpretation functions.
It characterizes all complete meet-preserving maps
between finite partition lattices.

\begin{theorem}[Representation of Complete Meet Homomorphisms]
  \label{ik.thm.representation}
  Let \(A\) and \(X\) be finite sets,
  and let
  \[
    K:
    \operatorname{Eq}(A)
    \longrightarrow
    \operatorname{Eq}(X)
  \]
  be a map.

  The following are equivalent:
  \begin{enumerate}[label=(\alph*)]
    \item \(K\) preserves all meets,
          including the empty meet;
    \item there exists an
          \(\operatorname{Eq}(A)\)-valued
          generalized ultrametric
          \[
            \delta:
            X\times X
            \longrightarrow
            \operatorname{Eq}(A)
          \]
          such that for every
          \(\rho\in\operatorname{Eq}(A)\)
          and \(x,y\in X\),
          \[
            x\mathrel{K(\rho)}y
            \quad\Longleftrightarrow\quad
            \delta(x,y)
            \sqsubseteq
            \rho.
          \]
  \end{enumerate}

  When it exists,
  the representing generalized ultrametric is unique and is given by
  \[
    \delta_K(x,y)
    =
    \bigwedge
    \bigl\{
    \rho
    \in
    \operatorname{Eq}(A)
    \bigm|
    x\mathrel{K(\rho)}y
    \bigr\}.
  \]
\end{theorem}

\begin{proof}
  Suppose first that \(K\)
  preserves all meets,
  including the empty meet.

  Empty-meet preservation gives
  \[
    K(\top_A)
    =
    \top_X.
  \]
  Therefore, for every \(x,y\in X\),
  the set
  \[
    S_{x,y}
    =
    \bigl\{
    \rho
    \in
    \operatorname{Eq}(A)
    \bigm|
    x\mathrel{K(\rho)}y
    \bigr\}
  \]
  is nonempty.

  Define
  \[
    \delta_K(x,y)
    =
    \bigwedge S_{x,y}.
  \]

  Since \(K\) preserves meets,
  \[
    K
    \bigl(
    \delta_K(x,y)
    \bigr)
    =
    \bigwedge_{\rho\in S_{x,y}}
    K(\rho).
  \]
  The pair \((x,y)\)
  belongs to every relation in the meet,
  so
  \[
    x
    \mathrel{
      K
      \left(
      \delta_K(x,y)
      \right)
    }
    y.
  \]

  If
  \[
    \delta_K(x,y)
    \sqsubseteq
    \rho,
  \]
  monotonicity of \(K\) gives
  \[
    K
    \bigl(
    \delta_K(x,y)
    \bigr)
    \sqsubseteq
    K(\rho).
  \]
  Hence
  \[
    x\mathrel{K(\rho)}y.
  \]

  Conversely,
  if
  \[
    x\mathrel{K(\rho)}y,
  \]
  then
  \[
    \rho\in S_{x,y},
  \]
  and therefore
  \[
    \delta_K(x,y)
    \sqsubseteq
    \rho.
  \]

  It remains to verify the generalized ultrametric laws.

  Reflexivity of every \(K(\rho)\)
  gives
  \[
    S_{x,x}
    =
    \operatorname{Eq}(A),
  \]
  so
  \[
    \delta_K(x,x)
    =
    \bot_A.
  \]

  Symmetry of every \(K(\rho)\)
  gives
  \[
    S_{x,y}
    =
    S_{y,x},
  \]
  and hence
  \[
    \delta_K(x,y)
    =
    \delta_K(y,x).
  \]

  Let
  \[
    \rho
    =
    \delta_K(x,y)
    \vee
    \delta_K(y,z).
  \]
  The threshold characterization already proved gives
  \[
    x\mathrel{K(\rho)}y
    \qquad\text{and}\qquad
    y\mathrel{K(\rho)}z.
  \]
  Transitivity of \(K(\rho)\)
  gives
  \[
    x\mathrel{K(\rho)}z.
  \]
  Therefore
  \[
    \delta_K(x,z)
    \sqsubseteq
    \rho,
  \]
  which is the strong triangle inequality.

  Conversely,
  suppose an
  \(\operatorname{Eq}(A)\)-valued
  generalized ultrametric \(\delta\)
  is given.
  Define
  \[
    K_{\delta}(\rho)
    =
    \bigl\{
    (x,y)
    \in
    X\times X
    \bigm|
    \delta(x,y)
    \sqsubseteq
    \rho
    \bigr\}.
  \]

  Reflexive zero gives reflexivity,
  and symmetry of \(\delta\)
  gives symmetry of \(K_{\delta}(\rho)\).

  If
  \[
    \delta(x,y)
    \sqsubseteq
    \rho
    \qquad\text{and}\qquad
    \delta(y,z)
    \sqsubseteq
    \rho,
  \]
  then
  \[
    \delta(x,y)
    \vee
    \delta(y,z)
    \sqsubseteq
    \rho.
  \]
  The strong triangle inequality gives
  \[
    \delta(x,z)
    \sqsubseteq
    \rho.
  \]
  Thus \(K_{\delta}(\rho)\)
  is transitive
  and therefore an equivalence relation.

  Finally, for any family
  \[
    \{
    \rho_i
    \}_{i\in J},
  \]
  \[
    \delta(x,y)
    \sqsubseteq
    \bigwedge_{i\in J}\rho_i
  \]
  holds when
  \[
    \delta(x,y)
    \sqsubseteq
    \rho_i
  \]
  for every \(i\in J\).
  Hence
  \[
    K_{\delta}
    \left(
    \bigwedge_{i\in J}\rho_i
    \right)
    =
    \bigwedge_{i\in J}
    K_{\delta}(\rho_i).
  \]

  For the empty family,
  the right-hand condition is vacuous,
  so
  \[
    K_{\delta}(\top_A)
    =
    \top_X.
  \]

  Uniqueness follows because any representing
  generalized ultrametric must satisfy
  \[
    \delta(x,y)
    =
    \bigwedge
    \bigl\{
    \rho
    \in
    \operatorname{Eq}(A)
    \bigm|
    x\mathrel{K(\rho)}y
    \bigr\}.
  \]
\end{proof}

\begin{corollary}[Representation of Interpretive Kernels]
  \label{ik.cor.interpretive-representation}
  For every finite probe family \(P\),
  the complete meet homomorphism
  \[
    K_P:
    \operatorname{Eq}(\mathcal{A})
    \longrightarrow
    \operatorname{Eq}(\mathcal{E})
  \]
  is represented by the unique generalized ultrametric
  \[
    \delta_P(F,F')
    =
    L_P
    \bigl(
    \gamma_{F,F'}
    \bigr).
  \]

  In particular,
  \[
    K_P(\rho)
    =
    \bigl\{
    (F,F')
    \in
    \mathcal{E}\times\mathcal{E}
    \bigm|
    \delta_P(F,F')
    \sqsubseteq
    \rho
    \bigr\}.
  \]
\end{corollary}

The representation converts a lattice homomorphism
into a pairwise threshold function.
Conversely,
the pairwise threshold function determines
the entire family of kernel equivalence relations.

\subsection{Exact Hypotheses and Limits}

The representation theorem depends on
complete meet preservation,
not on monotonicity alone.

A monotone map
\[
  K:
  \operatorname{Eq}(A)
  \longrightarrow
  \operatorname{Eq}(X)
\]
need not be recoverable
from least pairwise thresholds.
For a pairwise threshold to represent \(K\),
the pair must belong to \(K\)
at the meet of every regime
under which it belongs.
That is the closure property
supplied by meet preservation.

\begin{remark}[Empty-Meet Preservation]
  \label{ik.rem.empty-meet}
  Preservation of the empty meet is the condition
  \[
    K(\top_A)
    =
    \top_X.
  \]

  It ensures that every pair \(x,y\in X\)
  is equivalent under at least one occurrence regime,
  so the defining set for
  \(\delta_K(x,y)\)
  is nonempty.

  If only nonempty meets are assumed,
  top preservation must be imposed separately.
\end{remark}

\begin{remark}[No Join-Preservation Claim]
  \label{ik.rem.no-join-preservation}
  The representation theorem does not imply
  \[
    K(\rho\vee\sigma)
    =
    K(\rho)\vee K(\sigma).
  \]

  A right adjoint preserves meets.
  It need not preserve joins.

  The strong triangle inequality uses join
  inside the value lattice,
  but does not convert \(K\)
  into a join homomorphism.
\end{remark}

\begin{remark}[Canonicalization Remains Prior]
  \label{ik.rem.representation-canonicalization}
  The representation theorem begins with a map
  whose domain is already
  \(\operatorname{Eq}(\mathcal{A})\).

  It does not choose a canonical occurrence representation,
  construct a comparison lifting,
  or justify quotienting over alternative matchings.
  Those choices determine the carrier
  and occurrence regimes to which the theorem applies.

  A proposed comparison that fails meet preservation
  is not repaired by the representation theorem.
  It fails the theorem's hypothesis.
\end{remark}

\begin{remark}[No Scalar Ranking]
  \label{ik.rem.no-scalar-ranking}
  The values of \(\delta_P\)
  are equivalence relations
  and may be incomparable under refinement.

  The construction therefore supplies
  no canonical scalar distance,
  total ranking,
  or single numerical severity measure.
  Any scalarization would add structure
  not contained in the partition lattice.
\end{remark}

\begin{remark}[Finite and Infinite Forms]
  \label{ik.rem.finite-representation}
  The representation is not intrinsically finite.
  For any set \(A\), the equivalence relations
  \(\operatorname{Eq}(A)\) form a complete lattice,
  so the meets used in the proof of
  Theorem~\ref{ik.thm.representation}---the pairwise
  value \(\bigwedge S_{x,y}\) and the arbitrary meets
  \(\bigwedge_{i}\rho_i\)---exist without any finiteness
  assumption.
  The representation and its uniqueness
  therefore hold for arbitrary \(A\) and \(X\).

  Finiteness is required only by the computational claims
  of Section~\ref{ik.sec.computational-procedure}:
  interned occurrence identifiers, hashed signature grouping,
  generating-forest thresholds, finite witness extraction,
  and the executable schedule realization.
\end{remark}

\begin{remark}[Observational Rather Than Ontological Distance]
  \label{ik.rem.observational-distance}
  The value
  \[
    \delta_P(F,F')
  \]
  does not measure how different
  the frame environments are in themselves.

  It records the least occurrence regime
  that makes them indistinguishable
  under the selected probes
  and interpretation function.

  Changing \(P\),
  changing \(I\),
  changing the canonical occurrence carrier,
  or changing the comparison lifting
  may change the value
  without changing either frame representation.
\end{remark}

The representation result therefore characterizes
the algebra carried by interpretive kernels
without claiming that every proposed output comparison
produces such a kernel.

The next section supplies finite counterexamples
that separate the valid construction
from comparison rules that lose canonical structure
or fail the required preservation laws.

% The strongest result is not that our particular Kp
% has a distance-like description. It is the equivalence:
% complete meet-preserving partition map == equivalence-valued generalized ultrametric.
% That may be a standalone theorem rather than schedule-specific machinery.


% =============================================================================
\section{Finite Counterexamples}
\label{ik.sec.finite-counterexamples}
% =============================================================================

The preceding results depend on three structural choices.

First,
occurrence environments are placed in a declared canonical representation
before they are compared.
Second,
occurrence and frame regimes are required to be equivalence relations.
Third,
the interpretive-kernel map preserves meets
because comparison is performed pointwise
over a fixed probe family.

These requirements are not presentational conveniences.
This section gives finite counterexamples
to several natural alternatives.

The first counterexample shows that index-wise comparison is not well defined
on unordered output families without canonicalization.

The second rejects a proposed repair
that minimizes over all bijections.
That repair is permutation-invariant,
but it does not preserve meets
and its associated lattice-valued comparison
fails the strong triangle inequality.

The third separates resolution policies
from frame identity regimes.
A deterministic fallback procedure
need not induce an equivalence relation at all.

\subsection{Presentation Dependence Without Canonicalization}

Suppose raw interpretation produces
an unordered family of two occurrence values
drawn from a finite carrier
\[
  O
  =
  \{a,b\}.
\]

The unordered family
\[
  X
  =
  \{a,b\}
\]
has both ordered presentations
\[
  (a,b)
  \qquad\text{and}\qquad
  (b,a).
\]

Let equality be the component comparison.
Consider the attempted index-wise rule
\[
  (x_1,x_2)
  \mathrel{\widehat{=}}
  (y_1,y_2)
\]
when
\[
  x_1=y_1
  \qquad\text{and}\qquad
  x_2=y_2.
\]

Using the same presentation on both sides gives
\[
  (a,b)
  \mathrel{\widehat{=}}
  (a,b).
\]

Using two different presentations
of the same unordered family gives
\[
  (a,b)
  \mathrel{\not\widehat{=}}
  (b,a).
\]

The attempted relation therefore assigns
different comparison results
to different representatives
of the same unordered object.

\begin{proposition}[Index-Wise Comparison Does Not Descend]
  \label{ik.prop.no-canonicalization}
  Index-wise component comparison
  does not define a relation
  on unordered finite families
  unless a representative is selected
  by a declared canonicalization
  or the comparison is otherwise proved
  independent of representative choice.
\end{proposition}

\begin{proof}
  The family
  \[
    X=\{a,b\}
  \]
  is equal to itself as an unordered family.

  The ordered representatives
  \[
    (a,b)
    \qquad\text{and}\qquad
    (a,b)
  \]
  satisfy index-wise equality,
  while the ordered representatives
  \[
    (a,b)
    \qquad\text{and}\qquad
    (b,a)
  \]
  do not.

  The comparison therefore depends
  on the selected representatives
  and does not descend to the quotient
  by permutation.
\end{proof}

A declared canonicalization removes the ambiguity.
If a declared order places \(a\) before \(b\),
then both presentations are mapped to
\[
  (a,b),
\]
and comparison is performed
on that canonical indexed form.

The construction does not require
that every domain use ordering.
It requires that the domain declare
how an occurrence environment is represented
before index-wise comparison is applied.

\subsection{Failure of the Matching Comparison}

A natural response to
Proposition~\ref{ik.prop.no-canonicalization}
is to avoid selecting a canonical order
and instead compare unordered families
through possible bijections.

Let \(O\) be a finite component carrier,
and let
\[
  \rho
  \in
  \operatorname{Eq}(O).
\]

For equal-cardinality finite families
\(X\) and \(Y\),
define the proposed matching lift
\[
  X
  \mathrel{\widehat{\rho}_{\mathrm{m}}}
  Y
\]
when there exists a bijection
\[
  \sigma:
  X
  \longrightarrow
  Y
\]
such that
\[
  x\mathrel{\rho}\sigma(x)
  \qquad
  \text{for every }x\in X.
\]

The definition is invariant under permutation.
It nevertheless fails to preserve meets.

For compactness,
write a partition by listing its blocks.
Thus
\[
  ac|bd
\]
denotes
\[
  \bigl\{
  \{a,c\},
  \{b,d\}
  \bigr\}.
\]

\begin{proposition}[Matching Lift Does Not Preserve Meets]
  \label{ik.prop.matching-no-meets}
  There exist component equivalences
  \[
    \rho,\sigma
    \in
    \operatorname{Eq}(O)
  \]
  and unordered finite families \(X,Y\)
  such that
  \[
    X
    \mathrel{\widehat{\rho}_{\mathrm{m}}}
    Y
    \qquad\text{and}\qquad
    X
    \mathrel{\widehat{\sigma}_{\mathrm{m}}}
    Y,
  \]
  but
  \[
    X
    \mathrel{\not\widehat{
        \rho\wedge\sigma
      }_{\mathrm{m}}}
    Y.
  \]

  Consequently,
  \[
    \widehat{
      \rho\wedge\sigma
    }_{\mathrm{m}}
    \neq
    \widehat{\rho}_{\mathrm{m}}
    \wedge
    \widehat{\sigma}_{\mathrm{m}}.
  \]
\end{proposition}

\begin{proof}
  Let
  \[
    O
    =
    \{a,b,c,d\},
  \]
  and define
  \[
    \rho
    =
    ac|bd,
    \qquad
    \sigma
    =
    ad|bc.
  \]

  Let
  \[
    X
    =
    \{a,b\},
    \qquad
    Y
    =
    \{c,d\}.
  \]

  Under \(\rho\),
  the bijection
  \[
    a\longmapsto c,
    \qquad
    b\longmapsto d
  \]
  witnesses
  \[
    X
    \mathrel{\widehat{\rho}_{\mathrm{m}}}
    Y.
  \]

  Under \(\sigma\),
  the different bijection
  \[
    a\longmapsto d,
    \qquad
    b\longmapsto c
  \]
  witnesses
  \[
    X
    \mathrel{\widehat{\sigma}_{\mathrm{m}}}
    Y.
  \]

  The meet of the component equivalences is
  \[
    \rho\wedge\sigma
    =
    a|b|c|d,
  \]
  the equality relation on \(O\).

  The families \(X\) and \(Y\) are disjoint,
  so no bijection between them
  pairs equal component values.
  Therefore
  \[
    X
    \mathrel{\not\widehat{
        \rho\wedge\sigma
      }_{\mathrm{m}}}
    Y.
  \]
\end{proof}

The failure occurs because
the bijection witnessing comparison under \(\rho\)
need not be the bijection witnessing comparison under \(\sigma\).
Membership in the meet of the lifted relations
permits separate witnesses,
while lifting the meet requires one witness
that satisfies both component regimes simultaneously.

Permutation invariance therefore does not imply
meet preservation.

The same obstruction appears
in the lattice-valued version of the proposal.

Let
\[
  d:
  O\times O
  \longrightarrow
  \operatorname{Eq}(O)
\]
be defined by
\[
  d(u,v)
  =
  \operatorname{EqCl}_{O}
  \bigl(
  \{
  (u,v)
  \}
  \bigr).
\]

Thus \(d(u,v)\)
is the least component equivalence
that identifies \(u\) and \(v\).

\begin{lemma}[The Component Comparison Is a Generalized Ultrametric]
  \label{ik.lem.component-ultrametric}
  The function \(d\)
  is an
  \(\operatorname{Eq}(O)\)-valued
  generalized ultrametric.
\end{lemma}

\begin{proof}
  For every \(u\in O\),
  \[
    d(u,u)
    =
    \bot_O.
  \]

  Symmetry follows because
  the least equivalence relation containing
  \((u,v)\)
  is also the least equivalence relation containing
  \((v,u)\).

  For \(u,v,w\in O\),
  the join
  \[
    d(u,v)
    \vee
    d(v,w)
  \]
  contains both pairs
  \((u,v)\) and \((v,w)\).
  By transitivity,
  it contains \((u,w)\).
  Therefore
  \[
    d(u,w)
    \sqsubseteq
    d(u,v)
    \vee
    d(v,w).
  \]
\end{proof}

For equal-cardinality unordered families,
define the proposed best-matching comparison
\[
  d_{\mathrm{m}}(X,Y)
  =
  \bigwedge_{\sigma:X\cong Y}
  \;
  \bigvee_{x\in X}
  d
  \bigl(
  x,
  \sigma(x)
  \bigr).
\]

The inner join is the occurrence regime
required by one selected matching.
The outer meet attempts to retain
only the equivalences common
to all possible matching costs.

\begin{proposition}[Best-Matching Comparison Fails the Strong Triangle Law]
  \label{ik.prop.matching-no-triangle}
  The comparison \(d_{\mathrm{m}}\)
  need not satisfy
  \[
    d_{\mathrm{m}}(X,Z)
    \sqsubseteq
    d_{\mathrm{m}}(X,Y)
    \vee
    d_{\mathrm{m}}(Y,Z).
  \]
\end{proposition}

\begin{proof}
  Let
  \[
    O
    =
    \{a,b,c,d\},
  \]
  and take
  \[
    X
    =
    \{a,b\},
    \qquad
    Y
    =
    \{a,c\},
    \qquad
    Z
    =
    \{b,d\}.
  \]

  There are two bijections
  between each pair of families.

  For \(X\) and \(Y\),
  the two matching costs are
  \[
    a|bc|d
    \qquad\text{and}\qquad
    abc|d.
  \]
  Their meet is
  \[
    d_{\mathrm{m}}(X,Y)
    =
    a|bc|d.
  \]

  For \(Y\) and \(Z\),
  the two matching costs are
  \[
    ab|cd
    \qquad\text{and}\qquad
    ad|bc.
  \]
  Their meet is
  \[
    d_{\mathrm{m}}(Y,Z)
    =
    a|b|c|d
    =
    \bot_O.
  \]

  For \(X\) and \(Z\),
  the two matching costs are
  \[
    abd|c
    \qquad\text{and}\qquad
    ad|b|c.
  \]
  Their meet is
  \[
    d_{\mathrm{m}}(X,Z)
    =
    ad|b|c.
  \]

  Therefore
  \[
    d_{\mathrm{m}}(X,Y)
    \vee
    d_{\mathrm{m}}(Y,Z)
    =
    a|bc|d.
  \]

  But
  \[
    ad|b|c
    \not\sqsubseteq
    a|bc|d,
  \]
  because the left-hand relation identifies
  \(a\) with \(d\),
  while the right-hand relation does not.

  Hence
  \[
    d_{\mathrm{m}}(X,Z)
    \not\sqsubseteq
    d_{\mathrm{m}}(X,Y)
    \vee
    d_{\mathrm{m}}(Y,Z).
  \]
\end{proof}

The proposition gives an immediate failure
at the level of threshold relations.

Let
\[
  \rho
  =
  a|bc|d,
\]
and define
\[
  X
  \mathrel{R_{\rho}}
  Y
  \quad\Longleftrightarrow\quad
  d_{\mathrm{m}}(X,Y)
  \sqsubseteq
  \rho.
\]

Then
\[
  X\mathrel{R_{\rho}}Y
\]
and
\[
  Y\mathrel{R_{\rho}}Z,
\]
but
\[
  X\mathrel{\not R_{\rho}}Z.
\]

Thus \(R_{\rho}\)
is not transitive.
The best-matching comparison
does not merely fail
to produce the desired adjunction.
Its threshold relations need not be
frame regimes at all.

\begin{remark}[Source of the Failure]
  \label{ik.rem.matching-failure}
  For any fixed pair of bijections,
  composition and the component strong triangle law
  provide the expected bound.

  The invalid step is moving
  the meets over alternative bijections
  through the lattice join.

  That step would require a distributive law
  not available in
  \(\operatorname{Eq}(O)\),
  whose partition lattice is generally nondistributive.
\end{remark}

The canonicalized construction avoids this step.
After canonicalization,
one indexed comparison is made.
There is no meet over alternative matchings,
and arbitrary meets of component regimes
commute with the universal index-wise condition.

\subsection{Resolution Policies Are Not Frame Regimes}

A frame regime is an equivalence relation on the frame carrier.
A resolution policy instead selects which frame will be used
when a declaration is evaluated.

The two objects may interact, but are not interchangeable.

Consider three frame releases
\[
  r_0,
  r_1,
  r_2,
\]
ordered from oldest to newest.

Suppose a pin-with-fallback policy behaves as follows:

\begin{enumerate}[label=(\alph*)]
  \item a declaration pinned to \(r_i\)
        resolves to \(r_i\)
        when that release is available;
  \item when \(r_2\) is unavailable,
        the declaration may resolve to \(r_1\);
  \item when \(r_1\) is unavailable,
        the declaration may resolve to \(r_0\).
\end{enumerate}

If this policy is misread as a frame-sameness relation,
the resulting directed relation is
\[
  R_{\mathrm{pf}}
  =
  \bot_{\mathcal{E}}
  \cup
  \{
  (r_2,r_1),
  (r_1,r_0)
  \}.
\]

\begin{proposition}[Pin-With-Fallback Need Not Be a Frame Regime]
  \label{ik.prop.pin-fallback}
  The relation \(R_{\mathrm{pf}}\)
  is not an equivalence relation.
\end{proposition}

\begin{proof}
  The relation contains
  \[
    (r_2,r_1),
  \]
  but does not contain
  \[
    (r_1,r_2).
  \]

  It is therefore not symmetric.

  In the displayed finite policy,
  it also contains
  \[
    (r_2,r_1)
    \qquad\text{and}\qquad
    (r_1,r_0),
  \]
  but not
  \[
    (r_2,r_0),
  \]
  so it is not transitive either.
\end{proof}

The failure is not repaired
by requiring the policy to be deterministic.
A total deterministic function
\[
  q:
  D\times C
  \longrightarrow
  \mathcal{E}
\]
from declarations and availability contexts
to selected frames
is still a selection function,
not an equivalence relation
on \(\mathcal{E}\).

The same distinction applies to
always-latest,
latest-compatible,
and similar deployment practices.
They answer
\emph{which frame should be selected?}
rather than
\emph{which frames are the same?}

\begin{corollary}[Partition Validation Is Required]
  \label{ik.cor.partition-validation}
  A declared frame regime
  must be validated as an equivalence relation
  before it is supplied to
  \(L_P\),
  \(K_P\),
  or the threshold construction.

  A resolution policy that fails this validation
  may be represented as part of the frame environment
  or interpretation procedure,
  but not as a frame regime.
\end{corollary}

\subsection{Further Non-Implications}

The preceding examples delimit
several additional claims
that do not follow from the valid construction.

\begin{proposition}[Monotonicity Does Not Imply Meet Preservation]
  \label{ik.prop.monotone-not-meet}
  A monotone map between finite partition lattices
  need not preserve meets,
  even when it preserves the greatest element.
\end{proposition}

\begin{proof}
  Let
  \[
    A
    =
    \{a,b,c\},
    \qquad
    X
    =
    \{x,y\}.
  \]

  Define
  \[
    \alpha
    =
    ab|c,
    \qquad
    \beta
    =
    ac|b.
  \]

  Define
  \[
    M:
    \operatorname{Eq}(A)
    \longrightarrow
    \operatorname{Eq}(X)
  \]
  by
  \[
    M(\rho)
    =
    \begin{cases}
      \top_X,
       &
      \text{when }
      \alpha\sqsubseteq\rho
      \text{ or }
      \beta\sqsubseteq\rho;
      \\
      \bot_X,
       &
      \text{otherwise.}
    \end{cases}
  \]

  The map is monotone,
  and
  \[
    M(\top_A)
    =
    \top_X.
  \]

  However,
  \[
    \alpha\wedge\beta
    =
    \bot_A,
  \]
  so
  \[
    M
    \bigl(
    \alpha\wedge\beta
    \bigr)
    =
    \bot_X.
  \]

  On the other hand,
  \[
    M(\alpha)
    =
    M(\beta)
    =
    \top_X,
  \]
  and therefore
  \[
    M(\alpha)
    \wedge
    M(\beta)
    =
    \top_X.
  \]

  Hence
  \[
    M
    \bigl(
    \alpha\wedge\beta
    \bigr)
    \neq
    M(\alpha)
    \wedge
    M(\beta).
  \]
\end{proof}

Monotonicity alone therefore does not supply
the adjunction,
least pairwise thresholds,
or generalized-ultrametric representation.
The stronger meet-preservation law is essential.

\begin{proposition}[Meet Preservation Does Not Imply Join Preservation]
  \label{ik.prop.meet-not-join}
  An interpretive-kernel map
  may preserve every meet
  while failing to preserve joins.
\end{proposition}

\begin{proof}
  Let
  \[
    \mathcal{A}
    =
    \{a,b,c\},
    \qquad
    \mathcal{E}
    =
    \{F_0,F_1\},
    \qquad
    P
    =
    \{p\}.
  \]

  Define
  \[
    I(p,F_0)=a,
    \qquad
    I(p,F_1)=c.
  \]

  Let
  \[
    \rho
    =
    ab|c,
    \qquad
    \sigma
    =
    a|bc.
  \]

  Neither occurrence regime identifies
  \(a\) with \(c\).
  Therefore
  \[
    K_P(\rho)
    =
    K_P(\sigma)
    =
    \bot_{\mathcal{E}}.
  \]

  Their join is
  \[
    \rho\vee\sigma
    =
    \top_{\mathcal{A}},
  \]
  because the pairs \(a\sim b\)
  and \(b\sim c\)
  generate \(a\sim c\).

  Consequently,
  \[
    K_P
    \bigl(
    \rho\vee\sigma
    \bigr)
    =
    \top_{\mathcal{E}},
  \]
  while
  \[
    K_P(\rho)
    \vee
    K_P(\sigma)
    =
    \bot_{\mathcal{E}}.
  \]

  Thus
  \[
    K_P
    \bigl(
    \rho\vee\sigma
    \bigr)
    \neq
    K_P(\rho)
    \vee
    K_P(\sigma).
  \]
\end{proof}

\begin{remark}[Zero Demand Does Not Imply Frame Equality]
  \label{ik.rem.zero-not-equality}
  Distinct frame environments may satisfy
  \[
    \delta_P(F,F')
    =
    \bot_{\mathcal{A}}
  \]
  when every probe produces
  the same canonical occurrence environment
  under both.

  Zero pairwise demand establishes
  observational equivalence under \(P\),
  not equality of frame representation.
\end{remark}

\begin{remark}[Admissibility Does Not Imply Minimality]
  \label{ik.rem.admissible-not-minimal}
  An admissible occurrence regime
  need not equal the threshold.

  Admissibility gives
  \[
    \Theta_P(\varphi)
    \sqsubseteq
    \rho.
  \]

  Equality holds only when \(\rho\)
  introduces no occurrence equivalence
  beyond those required
  by the frame demand and probe family.
\end{remark}

\begin{remark}[Permutation Invariance Is Insufficient]
  \label{ik.rem.permutation-insufficient}
  The matching comparison is permutation-invariant,
  but fails both meet preservation
  and the generalized ultrametric law.

  Permutation invariance therefore does not by itself
  make an unordered comparison suitable
  for interpretive-kernel construction.
\end{remark}

The counterexamples identify the boundary
of the valid construction.

A domain may use unordered raw output,
resolution policies,
or comparison procedures other than index-wise equality.
It must nevertheless supply:

\begin{enumerate}[label=(\alph*)]
  \item a declared canonical occurrence carrier;
  \item occurrence and frame comparisons
        that are genuine equivalence relations;
  \item a comparison lifting
        satisfying the required preservation laws; and
  \item both the probe family
        and occurrence regime
        indexing every reported kernel.
\end{enumerate}

The schedule realization of the next section
meets these requirements
through chronological canonicalization,
finite probe horizons,
declared frame-environment tuples,
and index-wise occurrence comparison.

% The matching example is strong because the same six values do three jobs:
% it refutes the strong triangle law,
% produces a nontransitive threshold relation, and
% pinpoints nondistributivity as the failed proof step.

% =============================================================================
\section{Schedule Interpretation}
\label{ik.sec.schedule-interpretation}
% =============================================================================

This section instantiates the preceding construction
for finite schedule interpretation.

The schedule realization preserves the distinction
between semantic policy and external interpretive facts.
Denotation-affecting policy belongs to the schedule declaration
and has no implicit default.
The frame environment supplies the cited external artifacts
under which that declaration is interpreted.

A probe combines a schedule declaration
with a finite observation horizon
and a requested output channel.
Interpretation produces a finite occurrence environment,
which is placed in chronological canonical form
before an occurrence regime is applied.

\subsection{Declared Schedule Signature}

Let
\[
  \mathcal{S}
\]
be a finite carrier of well-formed schedule declarations.

A declaration contains the semantic choices
that determine its meaning,
including any applicable recurrence,
exception,
resolution,
ambiguity,
month-length,
observance,
or combination policy.
A denotation-affecting policy is represented explicitly.
It is not supplied by an undeclared implementation default.

Let
\[
  \mathcal{C}
\]
be a finite carrier of civil occurrence values,
and let
\[
  \mathcal{T}
\]
be a finite carrier of resolved instants.

The two carriers represent distinct output dimensions.

A civil value records a position
in a declared civil-time domain,
including its date,
local time,
and zone-bearing type where applicable.

A resolved value records an instant
on the resolved time line.

The tagged schedule-output carrier is
\[
  \mathcal{V}
  =
  \bigl(
  \{\mathsf{civil}\}
  \times
  \mathcal{C}
  \bigr)
  \;\cup\;
  \bigl(
  \{\mathsf{resolved}\}
  \times
  \mathcal{T}
  \bigr).
\]

The tags prevent a civil value
and a resolved instant
from being treated as values
of the same output dimension.

\begin{definition}[Schedule Probe]
  \label{ik.def.schedule-probe}
  A \emph{schedule probe} is a tuple
  \[
    p
    =
    \langle
    s,
    W,
    \chi
    \rangle,
  \]
  where:
  \begin{enumerate}[label=(\alph*)]
    \item \(s\in\mathcal{S}\)
          is a well-formed schedule declaration;
    \item \(W\) is a finite observation horizon; and
    \item
          \[
            \chi
            \in
            \{
            \mathsf{civil},
            \mathsf{resolved}
            \}
          \]
          is the requested output channel.
  \end{enumerate}
\end{definition}

A probe requests one typed view
of the schedule's occurrences.
Two probes may use the same declaration and horizon
while requesting different output channels.

Let
\[
  \mathsf{End}
  =
  \{
  \mathsf{exhausted},
  \mathsf{bounded}
  \}.
\]

The value \(\mathsf{exhausted}\)
records that the interpreted declaration
produced no further occurrence
because its denotation ended within the horizon.

The value \(\mathsf{bounded}\)
records that interpretation stopped
at the declared probe boundary.

\begin{definition}[Schedule Occurrence Environment]
  \label{ik.def.schedule-occurrence-environment}
  A \emph{schedule occurrence environment} is a pair
  \[
    A
    =
    \langle
    \mathbf{v},
    e
    \rangle,
  \]
  where
  \[
    \mathbf{v}
    =
    \langle
    v_1,
    \ldots,
    v_n
    \rangle
  \]
  is a finite canonical sequence
  of values from \(\mathcal{V}\),
  and
  \[
    e\in\mathsf{End}
  \]
  is its terminal tag.
\end{definition}

Let
\[
  \mathcal{A}_{\mathrm{sch}}
\]
be the finite carrier
of schedule occurrence environments
admitted to the comparison artifact.

The schedule interpretation function has signature
\[
  I_{\mathrm{sch}}:
  P
  \times
  \mathcal{E}_{\mathrm{sch}}
  \longrightarrow
  \mathcal{A}_{\mathrm{sch}}.
\]

For
\[
  p
  =
  \langle
  s,
  W,
  \chi
  \rangle
\]
and frame environment \(F\),
the value
\[
  I_{\mathrm{sch}}(p,F)
\]
is the canonical finite sequence
obtained by interpreting \(s\)
under \(F\) over \(W\),
projecting the result to channel \(\chi\),
and attaching the appropriate terminal tag.

\subsection{Chronological Canonicalization}

A schedule may generate occurrences
through unions,
intersections,
exceptions,
observance rules,
recurrence expansion,
or other compositional operators.

The order in which an implementation discovers those occurrences
is not part of their meaning
unless the schedule declaration explicitly makes it so.
Comparison therefore does not use
generation order.

\begin{definition}[Schedule Canonicalization]
  \label{ik.def.schedule-canonicalization}
  Schedule canonicalization maps
  a finite raw family of typed occurrence values
  to a finite sequence ordered:
  \begin{enumerate}[label=(\alph*)]
    \item primarily by chronological position
          in the requested output channel; and
    \item secondarily by a declared canonical ordering
          of the complete typed occurrence representation.
  \end{enumerate}

  Multiplicity is retained.
  The terminal tag is appended after ordering.
\end{definition}

For the resolved channel,
chronological order is order by instant.

For the civil channel,
chronological order is the declared order
of the civil-time type,
including its date,
local time,
and zone-bearing information.

The secondary order makes canonicalization deterministic
when distinct occurrence representations
share the same primary chronological position.

\begin{remark}[Canonicalization and Normalization]
  \label{ik.rem.schedule-normalization}
  Canonicalization orders
  a schedule's represented occurrences.
  It does not silently change their denotation.

  A domain normalization that alters,
  removes,
  combines,
  shifts,
  or resolves occurrences
  is a semantic operation
  and must be declared in the schedule expression
  or its explicit policy.

  Such an operation occurs before canonicalization.
\end{remark}

\begin{remark}[No Occurrence Matching]
  \label{ik.rem.schedule-no-matching}
  Two schedule occurrence environments
  are not compared by searching
  for a favorable bijection
  between their generated values.

  Each is first placed in chronological canonical form.
  The declared occurrence regime
  is then applied to the resulting indexed sequences,
  including their lengths, multiplicities, and terminal tags.
\end{remark}

\subsection{Finite Recurrence Horizons}

A recurring schedule may describe
an indefinitely extensible denotation.
That does not make a finite probe
or its occurrence environment infinite.

\begin{definition}[Observation Horizon]
  \label{ik.def.schedule-horizon}
  An \emph{observation horizon}
  is a finite declared interval
  \[
    W
    =
    [t_0,t_1)
  \]
  in the input domain
  appropriate to the schedule declaration.

  Interpretation includes the occurrences selected by the declaration
  whose relevant positions lie within \(W\),
  subject to its explicit boundary policy.
\end{definition}

The horizon belongs to the probe,
not to the schedule declaration's identity.
Changing the horizon changes the observation
and therefore produces a different probe.

\begin{proposition}[Finite Schedule Output]
  \label{ik.prop.finite-schedule-output}
  Suppose:
  \begin{enumerate}[label=(\alph*)]
    \item the schedule declaration is well formed;
    \item its operators are finite over every finite horizon;
    \item the frame environment resolves every cited external artifact; and
    \item the observation horizon is finite.
  \end{enumerate}

  Then
  \[
    I_{\mathrm{sch}}(p,F)
  \]
  is a finite schedule occurrence environment.
\end{proposition}

The proposition does not require
the schedule declaration to denote
a globally finite recurrence.

The sequence length and terminal tag
remain part of the comparison.
Two frames may produce the same canonical prefix
while differing in whether the schedule exhausted
or was merely cut off at the horizon.

\begin{remark}[Horizon Relativity]
  \label{ik.rem.schedule-horizon-relative}
  A kernel computed under one horizon
  is not a claim about all future occurrences.

  Extending the horizon may expose
  a frame distinction that was absent
  from the shorter probe family.
  In the abstract construction,
  this is ordinary probe-family extension.
\end{remark}

\subsection{Frame Environments}

A schedule frame environment
is a tuple of cited external artifacts
used during interpretation.

\begin{definition}[Schedule Frame Environment]
  \label{ik.def.schedule-frame}
  A \emph{schedule frame environment} is a tuple
  \[
    F
    =
    \langle
    z,
    h,
    b,
    v
    \rangle,
  \]
  where:
  \begin{enumerate}[label=(\alph*)]
    \item \(z\) identifies a time-zone rule artifact
          and its release;
    \item \(h\) identifies a holiday-calendar artifact
          and its release;
    \item \(b\) identifies a business-calendar artifact
          and its release; and
    \item \(v\) identifies any provider or authority artifact
          whose cited facts participate in interpretation.
  \end{enumerate}
\end{definition}

A concrete comparison artifact may fix
some components across every frame
and vary only the components under examination.

The frame tuple contains external facts.
It does not absorb semantic policy
that belongs in the schedule declaration.

For example,
a holiday artifact may state
which dates are holidays.
Whether a schedule skips,
shifts,
retains,
or otherwise treats those dates
is schedule policy.

Likewise,
a time-zone artifact supplies
the applicable offset rules.
A policy for resolving an ambiguous
or nonexistent civil time
belongs to the schedule declaration.

\begin{remark}[Frame Regimes and Resolution Policies]
  \label{ik.rem.schedule-frame-policy}
  A frame regime is an equivalence relation
  on schedule frame environments.

  Exact pinning,
  pin-with-fallback,
  latest-compatible,
  and always-latest are resolution policies
  for selecting a frame environment.
  They are not frame regimes
  unless a separate declared equivalence relation
  is supplied and validated.
\end{remark}

Frame regimes may compare complete tuples
or may be induced by declared projections.

For example,
the projection
\[
  q_z(F)
  =
  z
\]
induces the frame regime
\[
  F
  \mathrel{\varphi_z}
  F'
  \quad\Longleftrightarrow\quad
  q_z(F)=q_z(F').
\]

Similarly,
\[
  q_h(F)
  =
  h
\]
induces
\[
  F
  \mathrel{\varphi_h}
  F'
  \quad\Longleftrightarrow\quad
  q_h(F)=q_h(F').
\]

These regimes answer different questions.
Two frame environments may share
the same time-zone artifact
while differing in holiday data,
or share the same holiday artifact
while differing in time-zone rules.

\subsection{A Nondegenerate Finite Example}

The following finite construction
exhibits two independently varying frame components
and two typed occurrence dimensions.

It is an abstract schedule example.
A later worked application may replace
the symbolic releases and outputs
with specific cited artifacts,
dates,
and generated occurrences
without changing the construction.

Let the frame carrier be
\[
  \mathcal{E}_{\mathrm{sch}}
  =
  \{
  F_0,
  F_1,
  F_2
  \},
\]
where
\[
  F_0
  =
  \langle
  z_0,
  h_0,
  b_0,
  v_0
  \rangle,
\]
\[
  F_1
  =
  \langle
  z_1,
  h_0,
  b_0,
  v_0
  \rangle,
\]
and
\[
  F_2
  =
  \langle
  z_0,
  h_1,
  b_0,
  v_0
  \rangle.
\]

Thus \(F_1\) differs from \(F_0\)
only in its time-zone artifact,
while \(F_2\) differs from \(F_0\)
only in its holiday artifact.

Define the time-zone frame regime
by equality of the time-zone projection:
\[
  \mathcal{E}_{\mathrm{sch}}/\varphi_z
  =
  \bigl\{
  \{F_0,F_2\},
  \{F_1\}
  \bigr\}.
\]

Define the holiday frame regime
by equality of the holiday projection:
\[
  \mathcal{E}_{\mathrm{sch}}/\varphi_h
  =
  \bigl\{
  \{F_0,F_1\},
  \{F_2\}
  \bigr\}.
\]

The two frame regimes are non-discrete
and incomparable.

Let the probe family be
\[
  P
  =
  \{
  p_z,
  p_h
  \}.
\]

The probe \(p_z\)
requests resolved output
for a local-time schedule
whose instant differs between
the time-zone artifacts \(z_0\) and \(z_1\)
within the selected horizon.

The holiday artifacts do not affect this probe.

The probe \(p_h\)
requests civil output
for a holiday-relative schedule
whose civil occurrence differs between
the holiday artifacts \(h_0\) and \(h_1\)
within the selected horizon.

The time-zone artifacts do not affect this probe.

Let
\[
  u_0,u_1\in\mathcal{T},
  \qquad
  u_0\neq u_1,
\]
be the two resolved outputs,
and let
\[
  c_0,c_1\in\mathcal{C},
  \qquad
  c_0\neq c_1,
\]
be the two civil outputs.

For compactness,
write
\[
  A_{u_i}
  =
  \left\langle
  \left\langle
  (\mathsf{resolved},u_i)
  \right\rangle,
  \mathsf{bounded}
  \right\rangle
\]
and
\[
  A_{c_i}
  =
  \left\langle
  \left\langle
  (\mathsf{civil},c_i)
  \right\rangle,
  \mathsf{bounded}
  \right\rangle.
\]

The interpretation table is:

\[
  \begin{array}{c|cc}
        & p_z     & p_h     \\
    \hline
    F_0 & A_{u_0} & A_{c_0} \\
    F_1 & A_{u_1} & A_{c_0} \\
    F_2 & A_{u_0} & A_{c_1}
  \end{array}
\]

Let the examined occurrence carrier be
\[
  \mathcal{A}_0
  =
  \{
  A_{u_0},
  A_{u_1},
  A_{c_0},
  A_{c_1}
  \}.
\]

Under strict equality on \(\mathcal{A}_0\),
the three frame environments have distinct probe signatures:
\[
  \Sigma_{P,\bot}(F_0)
  =
  \langle
  A_{u_0},
  A_{c_0}
  \rangle,
\]
\[
  \Sigma_{P,\bot}(F_1)
  =
  \langle
  A_{u_1},
  A_{c_0}
  \rangle,
\]
and
\[
  \Sigma_{P,\bot}(F_2)
  =
  \langle
  A_{u_0},
  A_{c_1}
  \rangle.
\]

Therefore
\[
  K_P
  \bigl(
  \bot_{\mathcal{A}_0}
  \bigr)
  =
  \bot_{\mathcal{E}_{\mathrm{sch}}}.
\]

The threshold of the time-zone frame regime is
\[
  \Theta_P(\varphi_z)
  =
  A_{c_0}A_{c_1}
  \mid
  A_{u_0}
  \mid
  A_{u_1}.
\]

Indeed,
\(\varphi_z\) requires \(F_0\) and \(F_2\)
to remain equivalent.
Their resolved outputs already agree,
while their civil outputs differ.
The least supporting occurrence regime
therefore identifies
\(A_{c_0}\) with \(A_{c_1}\)
and introduces no other nontrivial equivalence.

Similarly,
the threshold of the holiday frame regime is
\[
  \Theta_P(\varphi_h)
  =
  A_{u_0}A_{u_1}
  \mid
  A_{c_0}
  \mid
  A_{c_1}.
\]

The regime \(\varphi_h\)
requires \(F_0\) and \(F_1\)
to remain equivalent.
Their civil outputs already agree,
while their resolved outputs differ.
The least supporting occurrence regime
therefore identifies
\(A_{u_0}\) with \(A_{u_1}\).

\begin{proposition}[Incomparable Schedule Thresholds]
  \label{ik.prop.schedule-incomparable-thresholds}
  In the finite schedule construction above,
  \[
    \Theta_P(\varphi_z)
  \]
  and
  \[
    \Theta_P(\varphi_h)
  \]
  are non-discrete and incomparable.

  Moreover,
  \[
    K_P
    \bigl(
    \Theta_P(\varphi_z)
    \bigr)
    =
    \varphi_z,
  \]
  and
  \[
    K_P
    \bigl(
    \Theta_P(\varphi_h)
    \bigr)
    =
    \varphi_h.
  \]
\end{proposition}

\begin{proof}
  The threshold
  \(\Theta_P(\varphi_z)\)
  identifies
  \(A_{c_0}\) with \(A_{c_1}\)
  but does not identify
  \(A_{u_0}\) with \(A_{u_1}\).

  The threshold
  \(\Theta_P(\varphi_h)\)
  identifies
  \(A_{u_0}\) with \(A_{u_1}\)
  but does not identify
  \(A_{c_0}\) with \(A_{c_1}\).

  Neither relation contains the other,
  so they are incomparable.

  Under \(\Theta_P(\varphi_z)\),
  the signatures of \(F_0\) and \(F_2\)
  agree in both probe positions,
  while the signature of \(F_1\)
  remains distinct through \(p_z\).
  Hence
  \[
    K_P
    \bigl(
    \Theta_P(\varphi_z)
    \bigr)
    =
    \varphi_z.
  \]

  The argument for \(\varphi_h\)
  is symmetric.
\end{proof}

The example is nondegenerate in four respects.

First,
neither frame regime is equality
or the indiscrete relation.

Second,
neither threshold is equality
or the indiscrete occurrence regime.

Third,
the frame regimes are incomparable.

Fourth,
their thresholds are incomparable
because they require equivalence
in different typed output dimensions.

Their combined demand satisfies
\[
  \Theta_P
  \bigl(
  \varphi_z
  \vee
  \varphi_h
  \bigr)
  =
  \Theta_P(\varphi_z)
  \vee
  \Theta_P(\varphi_h),
\]
which has partition
\[
  A_{c_0}A_{c_1}
  \mid
  A_{u_0}A_{u_1}.
\]

The civil and resolved channels remain distinct
because the output carrier is typed.
Combining the two frame demands
does not force a civil occurrence value
to be identified with a resolved instant.

This finite construction supplies
the schedule-level shape required by the theory.
The empirical application must additionally identify:

\begin{enumerate}[label=(\alph*)]
  \item the actual schedule declarations;
  \item the cited frame artifacts and releases;
  \item the finite horizons;
  \item the generated civil and resolved outputs;
  \item the declared frame regimes;
  \item the occurrence regimes admitted by the vocabulary; and
  \item the finite witnesses for every failed admissibility claim.
\end{enumerate}



% =============================================================================
\section{Computational Procedure}
\label{ik.sec.computational-procedure}
% =============================================================================

The constructions of the preceding sections are finite and executable.

A single table of canonical interpretation results
supports computation of interpretive kernels,
generated occurrence regimes,
thresholds,
pairwise interpretive demands,
admissibility verdicts,
and finite witnesses.

The algebraic operations are inexpensive
relative to schedule interpretation itself.
The potentially costly step is generating
and canonicalizing the occurrence environments.
The procedure therefore separates
interpretation cost from partition-processing cost
and reuses each computed interpretation result.

\subsection{Partition Construction}

Let \(X\) be a finite carrier.

A partition of \(X\) is represented computationally
by assigning a block label to every element.
Write
\[
  \lambda_{\alpha}:
  X
  \longrightarrow
  B_{\alpha}
\]
for a block-label function representing
an equivalence relation
\[
  \alpha
  \in
  \operatorname{Eq}(X),
\]
so that
\[
  x\mathrel{\alpha}y
  \quad\Longleftrightarrow\quad
  \lambda_{\alpha}(x)
  =
  \lambda_{\alpha}(y).
\]

The particular block identifiers are not significant.
Only equality of labels is used.

\begin{definition}[Normalized Partition Representation]
  \label{ik.def.normalized-partition}
  A \emph{normalized partition representation}
  of an equivalence relation \(\alpha\)
  over a finite carrier \(X\) consists of:
  \begin{enumerate}[label=(\alph*)]
    \item a stable identifier for every element of \(X\);
    \item one block identifier assigned to every element; and
    \item a block inventory containing the nonempty fibers of the block-label function.
  \end{enumerate}
\end{definition}

Supplying disjoint nonempty blocks
whose union is \(X\)
establishes partition-hood directly.

When a proposed regime is instead supplied
as an arbitrary binary relation,
the implementation must validate reflexivity,
symmetry,
and transitivity
before admitting it as an occurrence or frame regime.

\begin{definition}[Interpretation Table]
  \label{ik.def.interpretation-table}
  For a finite comparison artifact,
  the \emph{interpretation table} is
  \[
    T_I:
    P\times\mathcal{E}
    \longrightarrow
    \mathcal{A}
  \]
  with
  \[
    T_I(p,F)
    =
    I(p,F).
  \]

  Each table value is computed once,
  placed in canonical form,
  interned in the occurrence carrier,
  and thereafter referenced by its stable identifier.
\end{definition}

The table contains
\[
  |P|\,|\mathcal{E}|
\]
entries.

For schedule interpretation,
constructing one entry includes recurrence generation,
application of declared schedule policy,
projection to the requested output channel,
chronological canonicalization,
and attachment of the terminal tag.

\begin{proposition}[Kernel Construction by Signatures]
  \label{ik.prop.kernel-construction}
  Let
  \[
    \rho
    \in
    \operatorname{Eq}(\mathcal{A})
  \]
  be represented by the block-label function
  \(\lambda_{\rho}\).

  For every frame environment \(F\),
  construct the finite signature
  \[
    S_{\rho}(F)
    =
    \bigl(
    \lambda_{\rho}
    \bigl(
    T_I(p,F)
    \bigr)
    \bigr)_{p\in P}.
  \]

  Grouping frame environments
  by equality of these signatures
  produces the partition
  \[
    \mathcal{E}/K_P(\rho).
  \]
\end{proposition}

\begin{proof}
  Two frame environments \(F\) and \(F'\)
  receive equal signatures when
  \[
    \lambda_{\rho}
    \bigl(
    T_I(p,F)
    \bigr)
    =
    \lambda_{\rho}
    \bigl(
    T_I(p,F')
    \bigr)
  \]
  for every \(p\in P\).

  Equality of the block labels holds when
  \[
    I(p,F)
    \mathrel{\rho}
    I(p,F')
  \]
  for every \(p\in P\).

  By Definition~\ref{ik.def.kernel},
  this condition is exactly
  \[
    F\mathrel{K_P(\rho)}F'.
  \]
\end{proof}

A stable ordering of the probe identifiers
may be used to serialize the signatures.
The mathematical definition remains indexed by \(P\)
and does not depend on that implementation order.

Standard partition operations
are also computable from block labels.

\begin{proposition}[Finite Partition Operations]
  \label{ik.prop.partition-operations}
  Let
  \[
    \alpha,\beta
    \in
    \operatorname{Eq}(X)
  \]
  be represented by block labels.

  Then:
  \begin{enumerate}[label=(\alph*)]
    \item refinement
          \[
            \alpha
            \sqsubseteq
            \beta
          \]
          is decidable by checking that
          every \(\alpha\)-block lies within one \(\beta\)-block;
    \item the meet
          \[
            \alpha\wedge\beta
          \]
          is obtained by grouping each \(x\in X\)
          under the paired label
          \[
            \bigl(
            \lambda_{\alpha}(x),
            \lambda_{\beta}(x)
            \bigr);
          \]
    \item the join
          \[
            \alpha\vee\beta
          \]
          is obtained by initializing a union-find structure on \(X\)
          and joining a spanning set of members
          within every block of \(\alpha\)
          and every block of \(\beta\).
  \end{enumerate}
\end{proposition}

The join procedure does not enumerate
all pairs within each block.
A block with \(k\) members
requires only \(k-1\) union operations.

\subsection{Generating Forests}

Let
\[
  \varphi
  \in
  \operatorname{Eq}(\mathcal{E})
\]
be a frame regime.

For each block
\[
  C
  \in
  \mathcal{E}/\varphi,
\]
choose one member as its root
and connect that root
to every other member of \(C\).

The union of those blockwise stars
is a spanning forest
\[
  G_{\varphi}
  \subseteq
  \mathcal{E}\times\mathcal{E}
\]
that generates \(\varphi\).

Any other spanning tree within each block
produces the same equivalence closure.

\begin{definition}[Forest-Based Generated-Regime Procedure]
  \label{ik.def.forest-procedure}
  To compute
  \[
    L_P(\varphi),
  \]
  perform the following steps:
  \begin{enumerate}[label=(\arabic*)]
    \item initialize one union-find node
          for every occurrence environment
          in \(\mathcal{A}\);
    \item construct a spanning forest
          \(G_{\varphi}\)
          for the blocks of \(\varphi\);
    \item for every edge
          \[
            (F,F')
            \in
            G_{\varphi}
          \]
          and every probe \(p\in P\),
          union the occurrence identifiers
          \[
            T_I(p,F)
            \qquad\text{and}\qquad
            T_I(p,F');
          \]
    \item return the equivalence relation
          represented by the resulting
          union-find components.
  \end{enumerate}
\end{definition}

\begin{theorem}[Correctness of the Forest Procedure]
  \label{ik.thm.forest-procedure}
  The procedure of
  Definition~\ref{ik.def.forest-procedure}
  returns
  \[
    L_P(\varphi).
  \]
\end{theorem}

\begin{proof}
  The selected forest generates \(\varphi\), so
  \[
    \operatorname{EqCl}_{\mathcal{E}}
    \bigl(
    G_{\varphi}
    \bigr)
    =
    \varphi.
  \]

  The union operations construct
  \[
    \operatorname{EqCl}_{\mathcal{A}}
    \bigl(
    R_{P,G_{\varphi}}
    \bigr).
  \]

  By Theorem~\ref{ik.thm.generating-reduction},
  \[
    \operatorname{EqCl}_{\mathcal{A}}
    \bigl(
    R_{P,G_{\varphi}}
    \bigr)
    =
    L_P(\varphi).
  \]
\end{proof}

Because
\[
  \Theta_P(\varphi)
  =
  L_P(\varphi),
\]
the same procedure computes
the interpretive threshold.

For a principal pair regime
\(\gamma_{F,F'}\),
the generating forest consists of
the single nontrivial edge
\[
  (F,F').
\]
The procedure then computes
\[
  \delta_P(F,F')
\]
using one union operation per probe.

\begin{remark}[Release-Family Reduction]
  \label{ik.rem.release-family-reduction}
  When a frame regime partitions
  a finite release family,
  its spanning forest requires
  one fewer edge than the number of releases
  in each nonempty block.

  The procedure therefore does not compare
  every release with every other release.
  It interprets the selected probe family once per release
  and performs union operations only along
  the generating forest.
\end{remark}

\subsection{Complexity}

Let
\[
  n_{\mathcal{E}}
  =
  |\mathcal{E}|,
  \qquad
  n_{\mathcal{A}}
  =
  |\mathcal{A}|,
  \qquad
  n_P
  =
  |P|.
\]

For a frame regime \(\varphi\),
let
\[
  b_{\varphi}
  =
  |\mathcal{E}/\varphi|
\]
be its number of blocks.

Its spanning forest has
\[
  n_{\mathcal{E}}
  -
  b_{\varphi}
\]
edges.

Let
\[
  C_I
\]
denote the total cost
of computing and canonicalizing
all entries of the interpretation table:
\[
  C_I
  =
  \sum_{F\in\mathcal{E}}
  \sum_{p\in P}
  C_I(p,F).
\]

This quantity includes domain-specific work
such as recurrence generation
and chronological sorting.

Let
\[
  \alpha
\]
denote the inverse-Ackermann function
arising from union-find operations.

\begin{theorem}[Computational Complexity]
  \label{ik.thm.computational-complexity}
  Under the finite comparison assumptions:
  \begin{enumerate}[label=(\alph*)]
    \item the interpretation table requires
          \[
            C_I
          \]
          time and
          \[
            O
            \bigl(
            n_P n_{\mathcal{E}}
            \bigr)
          \]
          stored occurrence identifiers,
          in addition to the canonical outputs themselves;

    \item after the table and occurrence-regime labels are available,
          constructing
          \[
            K_P(\rho)
          \]
          by hashed signature grouping requires expected time
          \[
            O
            \bigl(
            n_P n_{\mathcal{E}}
            \bigr);
          \]

    \item deterministic sorting of the signatures
          gives the alternative bound
          \[
            O
            \bigl(
            n_P n_{\mathcal{E}}
            \log n_{\mathcal{E}}
            \bigr);
          \]

    \item constructing a spanning forest
          for \(\varphi\)
          requires
          \[
            O
            \bigl(
            n_{\mathcal{E}}
            \bigr)
          \]
          time when its blocks are already known;

    \item computing
          \[
            L_P(\varphi)
            =
            \Theta_P(\varphi)
          \]
          from that forest requires
          \[
            n_P
            \bigl(
            n_{\mathcal{E}}
            -
            b_{\varphi}
            \bigr)
          \]
          union operations and therefore
          \[
            O
            \left(
            n_P
            \bigl(
              n_{\mathcal{E}}
              -
              b_{\varphi}
              \bigr)
            \alpha
            \bigl(
              n_{\mathcal{A}}
              \bigr)
            \right)
          \]
          additional time;

    \item computing one pairwise interpretive demand
          \[
            \delta_P(F,F')
          \]
          requires
          \[
            n_P
          \]
          union operations and therefore
          \[
            O
            \left(
            n_P
            \alpha
            \bigl(
              n_{\mathcal{A}}
              \bigr)
            \right)
          \]
          time after interpretation-table construction.
  \end{enumerate}
\end{theorem}

The algebraic overhead is therefore near-linear
in the size of the finite carriers
and selected probe family.

The procedure avoids:

\begin{enumerate}[label=(\alph*)]
  \item factorial enumeration of occurrence matchings;
  \item quadratic enumeration of all pairs
        within frame-regime blocks; and
  \item repeated interpretation of the same probe
        under the same frame environment.
\end{enumerate}

The interpretation cost \(C_I\)
may dominate the partition operations.
That is expected.
The computational burden lies
in producing the domain interpretation,
not in the finite lattice algebra
used to compare its results.

\begin{remark}[Cost of Canonicalization]
  \label{ik.rem.canonicalization-cost}
  Suppose one schedule interpretation result
  contains \(k\) generated occurrences.

  Chronological canonicalization
  may require
  \[
    O(k\log k)
  \]
  sorting time.

  Once canonicalized,
  index-wise processing is linear in \(k\).

  The rejected matching comparison
  would instead range over
  as many as
  \[
    k!
  \]
  bijections,
  with lattice operations performed
  for each matching.
\end{remark}

\begin{remark}[Incremental Probe Addition]
  \label{ik.rem.incremental-probes}
  Suppose a kernel has already been computed
  for a probe family \(P\),
  and a new probe \(q\) is added.

  The new interpretation values
  \[
    T_I(q,F)
  \]
  are computed once for every
  \(F\in\mathcal{E}\).

  The probe-union law gives
  \[
    K_{P\cup\{q\}}(\rho)
    =
    K_P(\rho)
    \wedge
    K_{\{q\}}(\rho).
  \]

  Threshold construction is likewise incremental:
  \[
    \Theta_{P\cup\{q\}}(\varphi)
    =
    \Theta_P(\varphi)
    \vee
    \Theta_{\{q\}}(\varphi).
  \]
\end{remark}

The interpretation table may therefore be extended
one probe at a time,
and the corresponding kernel or threshold
updated by one partition meet or join.

\subsection{Witness Extraction}

Every failed admissibility claim
has a finite witness.

\begin{definition}[Admissibility-Failure Witness]
  \label{ik.def.admissibility-witness}
  Let
  \[
    \varphi
    \in
    \operatorname{Eq}(\mathcal{E})
  \]
  be a frame regime,
  and let
  \[
    \rho
    \in
    \operatorname{Eq}(\mathcal{A})
  \]
  be a proposed occurrence regime.

  An \emph{admissibility-failure witness}
  is a tuple
  \[
    \langle
    F,
    F',
    p,
    A,
    A'
    \rangle
  \]
  such that:
  \begin{enumerate}[label=(\alph*)]
    \item
          \[
            F\mathrel{\varphi}F';
          \]
    \item
          \[
            A
            =
            I(p,F);
          \]
    \item
          \[
            A'
            =
            I(p,F');
          \]
    \item
          \[
            A
            \not\mathrel{\rho}
            A'.
          \]
  \end{enumerate}
\end{definition}

The witness identifies:

\begin{enumerate}[label=(\alph*)]
  \item the frame pair
        the declaration requires to remain equivalent;
  \item the probe exposing the disagreement; and
  \item the two canonical occurrence environments
        the proposed occurrence regime fails to identify.
\end{enumerate}

\begin{theorem}[Finite Admissibility Witness]
  \label{ik.thm.finite-admissibility-witness}
  The occurrence regime \(\rho\)
  is not admissible for \(\varphi\) under \(P\)
  if and only if
  an admissibility-failure witness exists.
\end{theorem}

\begin{proof}
  By Definition~\ref{ik.def.admissible},
  admissibility fails when
  \[
    \varphi
    \not\sqsubseteq
    K_P(\rho).
  \]

  This holds when
  there exist \(F,F'\in\mathcal{E}\)
  such that
  \[
    F\mathrel{\varphi}F'
  \]
  but
  \[
    F
    \not\mathrel{K_P(\rho)}
    F'.
  \]

  By Definition~\ref{ik.def.kernel},
  the latter condition holds when
  some \(p\in P\) satisfies
  \[
    I(p,F)
    \not\mathrel{\rho}
    I(p,F').
  \]

  These finite objects form
  the required witness.
\end{proof}

A witness may be extracted
without enumerating every pair in \(\varphi\).

\begin{theorem}[Generating-Forest Witness Extraction]
  \label{ik.thm.forest-witness}
  Let \(G_{\varphi}\)
  be any spanning forest generating \(\varphi\).

  If \(\rho\)
  is not admissible for \(\varphi\),
  then there exist
  \[
    (F,F')
    \in
    G_{\varphi}
  \]
  and
  \[
    p\in P
  \]
  such that
  \[
    I(p,F)
    \not\mathrel{\rho}
    I(p,F').
  \]
\end{theorem}

\begin{proof}
  Suppose every forest edge
  satisfies
  \[
    I(p,F)
    \mathrel{\rho}
    I(p,F')
  \]
  for every \(p\in P\).

  For any two frame environments
  in the same \(\varphi\)-block,
  the forest supplies a finite path
  between them.

  Applying the same probe
  along that path
  and using transitivity of \(\rho\)
  shows that their interpreted outputs
  are \(\rho\)-equivalent.

  Therefore every pair in \(\varphi\)
  lies in \(K_P(\rho)\),
  so
  \[
    \varphi
    \sqsubseteq
    K_P(\rho).
  \]

  This contradicts failed admissibility.
  Hence some forest edge and probe
  provide the required witness.
\end{proof}

After construction of the interpretation table,
the forest-based admissibility test examines at most
\[
  n_P
  \bigl(
  n_{\mathcal{E}}
  -
  b_{\varphi}
  \bigr)
\]
output-label pairs.

It stops immediately
when the first mismatch is found.

\begin{definition}[Kernel-Separation Witness]
  \label{ik.def.kernel-separation-witness}
  For frame environments \(F,F'\)
  lying in different blocks of
  \(K_P(\rho)\),
  a \emph{kernel-separation witness}
  is a probe \(p\in P\)
  satisfying
  \[
    I(p,F)
    \not\mathrel{\rho}
    I(p,F').
  \]
\end{definition}

Such a witness is obtained
by scanning the two stored signatures
until their first unequal
\(\rho\)-block label is found.

\begin{definition}[Refinement-Failure Witness]
  \label{ik.def.refinement-witness}
  For equivalence relations
  \[
    \alpha,\beta
    \in
    \operatorname{Eq}(X),
  \]
  a \emph{refinement-failure witness}
  for
  \[
    \alpha
    \not\sqsubseteq
    \beta
  \]
  is a pair
  \[
    (x,y)
    \in
    X\times X
  \]
  such that
  \[
    x\mathrel{\alpha}y
    \qquad\text{and}\qquad
    x\not\mathrel{\beta}y.
  \]
\end{definition}

A refinement witness can be extracted
by scanning each \(\alpha\)-block
and checking whether all of its members
share one \(\beta\)-block label.

If two partitions are incomparable, the procedure returns two witnesses:

\[
  x\mathrel{\alpha}y
  \qquad\text{and}\qquad
  x\not\mathrel{\beta}y,
\]
together with
\[
  u\mathrel{\beta}v
  \qquad\text{and}\qquad
  u\not\mathrel{\alpha}v.
\]

For threshold comparison, these pairs state
which occurrence equivalence one demand requires and the other omits.

\begin{theorem}[Finite Decidability]
  \label{ik.thm.finite-decidability}
  Under the finite comparison assumptions,
  and assuming that the interpretation function
  and declared canonicalization are computable,
  the following are computable:
  \begin{enumerate}[label=(\alph*)]
    \item the interpretive kernel
          \(K_P(\rho)\);
    \item the generated occurrence regime
          \(L_P(\varphi)\);
    \item the threshold
          \(\Theta_P(\varphi)\);
    \item the pairwise interpretive demand
          \(\delta_P(F,F')\);
    \item admissibility of \(\rho\)
          for \(\varphi\);
    \item refinement,
          equality,
          strict refinement,
          and incomparability
          of finite regimes; and
    \item finite witnesses
          refuting every failed admissibility
          or refinement claim.
  \end{enumerate}
\end{theorem}

The procedure does not discover
the correct probe family,
frame carrier,
occurrence carrier,
interpretation signature,
canonicalization,
or declared regimes.

Those are semantic declarations supplied by the domain realization.

Given those declarations,
the procedure computes their consequences and returns finite witnesses
when a proposed occurrence regime does not support a declared frame demand.


% =============================================================================
\section{Evidence Plan}
\label{ik.sec.evidence-plan}
% =============================================================================

The mathematical results of this paper
hold for every finite comparison artifact
satisfying the stated hypotheses.

The schedule interpretation is intended to establish a different claim:
that the construction identifies nontrivial and practically relevant structure
in a real domain.

That claim is not established by the algebra alone.

The evidence plan therefore separates:

\begin{enumerate}[label=(\alph*)]
  \item the proved properties
        of interpretive kernels;
  \item the prospective classification
        of a previously assembled schedule-question corpus;
  \item the executable computation of kernels and thresholds
        over pinned schedule artifacts; and
  \item the residuals that would limit, revise,
        or disconfirm the proposed schedule realization.
\end{enumerate}

No corpus classification, kernel result, or residual count is reported in this section.
Those results must be generated from the frozen evidence artifacts
described below.

\subsection{Question Corpus}

The initial evidence base is
a corpus of fifty-one schedule-sameness questions
assembled before the present classification.

The questions concern recurring rules,
exceptions,
time-zone changes,
holiday and business calendars,
copied and overlaid calendars,
identifiers,
version histories,
cancellation and restoration,
declared and executed schedules,
boundaries,
granularity,
and frame-dependent interpretation.

The corpus was assembled as a challenge set,
not as a random sample
from a defined population of schedule questions.
It therefore supports
finite coverage,
separation,
and residual claims
about the examined questions.
It does not support
statistical prevalence claims
about scheduling practice as a whole.

\begin{definition}[Question-Corpus Entry]
  \label{ik.def.corpus-entry}
  A \emph{question-corpus entry}
  contains:
  \begin{enumerate}[label=(\alph*)]
    \item a stable question identifier;
    \item the question text as originally recorded;
    \item its source or construction provenance;
    \item the date and version at which it entered the corpus;
    \item any examples originally attached to it; and
    \item an amendment history
          that preserves rather than replaces
          the original wording.
  \end{enumerate}
\end{definition}

The original question text is immutable
after the classification protocol is frozen.
A correction,
clarification,
or scope restriction
is recorded as a new attributed amendment.

The corpus preserves the distinction
between observed phenomena
and explanatory classifications.

Apparent groupings such as
rule versus extension,
frame dependence,
per-observer copies,
identifier versus content,
persistence through history,
and declared versus executed behavior
were recorded as observations.

They are not treated as established regimes,
kernel classes,
or proof that a new primitive is required.

\begin{definition}[Kernel Analysis Record]
  \label{ik.def.kernel-analysis-record}
  For each corpus question admitted
  to kernel analysis,
  the evidence record declares:
  \begin{enumerate}[label=(\alph*)]
    \item the schedule declaration or declarations;
    \item the finite probe family;
    \item the observation horizon of each probe;
    \item the schedule frame carrier;
    \item the interpretation function
          and canonicalization procedure;
    \item the candidate frame regime;
    \item the candidate occurrence regime;
    \item the resulting interpretive kernel;
    \item the generated threshold;
    \item the admissibility verdict; and
    \item any finite witness
          supporting a failed refinement
          or admissibility claim.
  \end{enumerate}
\end{definition}

A corpus question need not supply
all of these declarations in its original wording.

Missing information is not silently inferred.
It is recorded as part of the classification result.

\begin{definition}[Question Disposition]
  \label{ik.def.question-disposition}
  Each corpus question receives
  one primary disposition:
  \begin{enumerate}[label=(D\arabic*)]
    \item \emph{kernel-expressible}:
          the question determines
          a finite comparison artifact
          and a declared kernel comparison;

    \item \emph{object boundary undeclared}:
          the question does not determine
          which objects form the comparison carrier;

    \item \emph{frame regime undeclared}:
          the question does not determine
          which frame environments
          are required to count as equivalent;

    \item \emph{occurrence regime undeclared}:
          the question does not determine
          which generated occurrence environments
          count as equivalent;

    \item \emph{probe domain undeclared}:
          the question omits
          the horizon,
          output channel,
          or other observation declaration
          required to make interpretation finite;

    \item \emph{resolution policy presented as regime}:
          the stated practice selects a frame
          but does not define
          an equivalence relation on frames;

    \item \emph{interpretation incomplete}:
          the declared inputs do not determine
          a total interpretation function
          on the examined artifact;

    \item \emph{outside kernel scope}:
          the question concerns another relation,
          such as causal responsibility,
          normative authority,
          execution compliance,
          or object identity
          without an interpretive-equivalence claim; or

    \item \emph{residual}:
          the question is well scoped
          but cannot be represented faithfully
          by the declared kernel construction.
  \end{enumerate}
\end{definition}

These dispositions are mutually exclusive
at the primary level.
An evidence record may carry
additional secondary diagnoses.

In particular,
\emph{object boundary undeclared}
is distinct from
\emph{frame regime undeclared}
and
\emph{occurrence regime undeclared}.

The first concerns
what the compared objects are.
The latter two concern
which equivalence relations
are declared over known carriers.

For a kernel-expressible question,
the computational result is recorded separately.

\begin{definition}[Kernel Result Class]
  \label{ik.def.kernel-result-class}
  A kernel-expressible question receives
  one of the following result classes:
  \begin{enumerate}[label=(K\arabic*)]
    \item \emph{admissible at threshold}:
          the examined occurrence regime
          equals the generated threshold;

    \item \emph{admissible above threshold}:
          the examined occurrence regime
          is strictly coarser than the threshold;

    \item \emph{inadmissible below threshold}:
          the examined occurrence regime
          is strictly finer than the threshold;

    \item \emph{inadmissible and incomparable}:
          the examined occurrence regime
          and threshold are incomparable; or

    \item \emph{observational collapse}:
          distinct declared frame environments
          produce zero pairwise demand
          under the selected probe family.
  \end{enumerate}
\end{definition}

The result classes report
the relation between
a proposed occurrence regime
and the generated threshold.

They do not rank the substantive importance
of the question
or determine which regime
an institution ought to adopt.

\subsection{Predicted Classifications}

The evidence record preserves
two stages of prediction.

The first stage consists
of the competing predictions
recorded before the corpus was classified
and before the present kernel algebra
was established.

Those predictions concern:

\begin{enumerate}[label=(\alph*)]
  \item how completely
        the referential regimes
        of the earlier Structural Explainability work
        would classify the schedule questions;
  \item whether frame identity
        would require a new regime,
        operate as an external context parameter,
        or reuse the existing regime machinery
        over frame objects;
  \item whether apparent schedule phenomena
        would collapse into fewer explanatory distinctions;
  \item whether some established regimes
        would receive no support from the corpus; and
  \item whether any question
        would remain well scoped
        but unclassifiable.
\end{enumerate}

Those predictions remain frozen
in their original preregistration artifact and
are not rewritten to agree with the results of this paper.

The present work adds
a second stage of predictions
specific to interpretive kernels.

\begin{prediction}[Nondegenerate Schedule Kernels]
  \label{ik.prediction.nondegenerate}
  At least one pinned schedule comparison
  will produce an interpretive kernel
  that is neither
  \[
    \bot_{\mathcal{E}}
  \]
  nor
  \[
    \top_{\mathcal{E}}.
  \]
\end{prediction}

This prediction tests whether
the construction has discriminating power
in the schedule domain.

If every examined kernel is discrete,
the probes distinguish every frame environment
under the selected occurrence regimes.

If every examined kernel is indiscrete,
the probes distinguish none.

Either extreme is mathematically valid,
but a corpus concentrated at one extreme
would weaken the claim
that the kernel supplies useful schedule structure.

\begin{prediction}[Occurrence-Regime Sensitivity]
  \label{ik.prediction.regime-sensitivity}
  There exists at least one fixed probe family \(P\),
  frame carrier,
  and pair of occurrence regimes
  \[
    \rho
    \sqsubset
    \sigma
  \]
  for which
  \[
    K_P(\rho)
    \subsetneq
    K_P(\sigma).
  \]
\end{prediction}

The theorem permits strict dependence
but does not guarantee
that a selected schedule corpus exhibits it.

The evidence must therefore include
at least one concrete frame pair
whose interpretive equivalence changes
when the occurrence comparison changes.

\begin{prediction}[Nontrivial Thresholds]
  \label{ik.prediction.nontrivial-thresholds}
  At least one declared schedule frame regime
  has a threshold
  that is neither
  the equality relation
  nor the indiscrete occurrence regime.
\end{prediction}

A nontrivial threshold identifies
some generated occurrence environments
while preserving distinctions among others.

Such a result would show
that the schedule demand is not merely
``all outputs must be equal''
or
``all outputs may be collapsed.''

\begin{prediction}[Incomparable Schedule Demands]
  \label{ik.prediction.incomparable-demands}
  The schedule evidence contains
  two declared frame demands
  whose generated occurrence thresholds
  are incomparable under refinement.
\end{prediction}

The candidate construction
separates time-zone-sensitive
and calendar-sensitive demands.

The empirical result requires
actual pinned schedule declarations,
frame artifacts,
and generated occurrence environments
exhibiting the two incomparable requirements.

\begin{prediction}[Policy-Regime Separation]
  \label{ik.prediction.policy-regime}
  At least one production-style practice
  initially described as frame sameness
  will fail equivalence-relation validation
  and will instead be classified
  as a resolution policy.
\end{prediction}

Pin-with-fallback,
latest-compatible,
and always-latest
are candidate practices.

The evidence must test
the actual declared relation.
The classification is not assigned
from the policy's name alone.

\begin{prediction}[Frame-Component Interaction]
  \label{ik.prediction.frame-interaction}
  At least one schedule interpretation
  depends jointly on two components
  of the frame tuple
  in a way not recoverable
  by independently computing
  and combining per-component kernels.
\end{prediction}

Business-day interpretation
under both a holiday calendar
and time-zone rules
is a candidate interaction.

This prediction is supported only if
a finite witness shows that
the complete frame kernel differs
from the proposed componentwise construction.

\begin{prediction}[Parameter Necessity]
  \label{ik.prediction.parameter-necessity}
  The corpus contains questions
  for which reporting the probe family
  without the occurrence regime,
  or the occurrence regime
  without the probe family,
  leaves the interpretive-equivalence claim
  underdetermined.
\end{prediction}

This is an evidentiary prediction
about the questions and declarations,
not a restatement of
Propositions~\ref{ik.prop.probe-dependence}
and~\ref{ik.prop.regime-dependence}.

The finite evidence must show
the same schedule material
receiving different classifications
under the omitted parameter.

The following mathematical laws
are not registered as empirical predictions:

\begin{enumerate}[label=(\alph*)]
  \item monotonicity of \(K_P\);
  \item meet preservation by \(K_P\);
  \item join preservation by \(L_P\);
  \item the adjunction;
  \item probe-union duality; and
  \item existence of finite witnesses under the stated finite hypotheses.
\end{enumerate}

They are proved results.

Executable checks of those laws
are conformance tests
for the implementation.

A failed executable law
indicates an implementation defect,
an undeclared change of carrier,
or violation of a theorem hypothesis.

\subsection{Provenance and Preregistration}

The evidence package is maintained
separately from the mathematical manuscript
and the schedule implementation.

The manuscript states the protocol and reports the results.

The schedule implementation computes
interpretations,
canonical forms,
partitions,
thresholds,
and witnesses.

The evidence repository preserves
the frozen inputs,
preregistered predictions,
execution manifests,
and generated result bundle.

\begin{definition}[Evidence Release]
  \label{ik.def.evidence-release}
  An \emph{evidence release}
  contains:
  \begin{enumerate}[label=(\alph*)]
    \item the immutable question corpus;
    \item the original preregistration;
    \item the kernel-specific preregistration;
    \item the schedule declarations;
    \item the probe declarations
          and finite horizons;
    \item the frame-environment inventory;
    \item the candidate frame regimes;
    \item the candidate occurrence regimes;
    \item the canonicalization declaration;
    \item the pinned interpreter
          and dependency versions;
    \item the cited external frame artifacts;
    \item the complete interpretation table;
    \item the computed kernels,
          thresholds,
          and pairwise demands;
    \item every finite failure witness;
    \item the classification results; and
    \item a manifest identifying
          the hashes and versions
          of all included artifacts.
  \end{enumerate}
\end{definition}

For schedule evidence, the pinned external artifacts include, where applicable:

\begin{enumerate}[label=(\alph*)]
  \item time-zone database releases;
  \item holiday-calendar releases;
  \item business-calendar releases;
  \item provider or authority artifacts;
  \item schedule vocabulary releases; and
  \item the schedule interpretation implementation.
\end{enumerate}

A human-readable release name
is not sufficient provenance.
The evidence record includes
a stable artifact identifier
and content digest.

\begin{definition}[Preregistration Boundary]
  \label{ik.def.preregistration-boundary}
  The \emph{preregistration boundary}
  is the recorded point
  after which the following items
  may not be altered
  within the confirmatory evidence release:
  \begin{enumerate}[label=(\alph*)]
    \item the corpus membership
          and original question wording;
    \item the prediction statements;
    \item the primary classification dispositions;
    \item the candidate frame
          and occurrence regimes;
    \item the probe families
          and horizons;
    \item the canonicalization rules;
    \item the pinned external artifacts; and
    \item the criteria for identifying residuals.
  \end{enumerate}
\end{definition}

The preregistration is assigned a version, timestamp, and content digest
before confirmatory classification begins.

A later correction does not overwrite
the frozen record.
It creates a new release
with an amendment explaining:

\begin{enumerate}[label=(\alph*)]
  \item what changed;
  \item why it changed;
  \item whether the change occurred
        before or after inspecting results; and
  \item which conclusions depend on the change.
\end{enumerate}

New questions,
new probes,
longer horizons,
new frame releases,
and newly discovered occurrence regimes
are permitted after preregistration,
but belong to an exploratory extension set.

They do not silently enter the confirmatory corpus.

\begin{definition}[Result Provenance Record]
  \label{ik.def.result-provenance}
  Every computed kernel result
  is accompanied by:
  \begin{enumerate}[label=(\alph*)]
    \item the question identifier;
    \item the probe-family identifier;
    \item the frame-carrier identifier;
    \item the frame-regime identifier;
    \item the occurrence-regime identifier;
    \item the interpretation-table digest;
    \item the implementation version;
    \item the canonicalization version;
    \item the computed partition;
    \item the threshold partition;
    \item the verdict; and
    \item any supporting witness identifiers.
  \end{enumerate}
\end{definition}

The result record makes explicit
the two constitutive kernel parameters.

A statement that reports ``the schedule kernel'' without identifying
both the probe family and occurrence regime is incomplete.

The evidence bundle also preserves negative results.

A question classified as outside kernel scope,
a policy rejected as a non-partition,
a degenerate kernel,
an unexercised regime,
or a residual remains in the released result set.

\subsection{Qualifying Residuals}

A residual is not simply a question without enough information.

Missing boundaries, missing regimes, missing horizons, or missing frame artifacts
are declaration failures and receive their corresponding dispositions.

A residual remains after the required declarations have been supplied.

\begin{definition}[Faithful Residual]
  \label{ik.def.residual}
  A corpus question is a \emph{faithful residual}
  when:
  \begin{enumerate}[label=(\alph*)]
    \item the compared objects
          and finite carriers are explicit;
    \item the probe family
          and finite horizons are explicit;
    \item the frame environments
          and candidate frame regime are explicit;
    \item the occurrence carrier,
          canonicalization,
          and candidate occurrence regimes are explicit;
    \item interpretation is total
          on the examined artifact;
    \item the question makes
          a genuine interpretive-equivalence demand;
    \item the demand cannot be represented
          by the kernel construction
          without adding,
          erasing,
          or misclassifying
          a distinction present in the question; and
    \item the failure is not repaired
          by selecting another already-declared
          frame or occurrence regime.
  \end{enumerate}
\end{definition}

The faithfulness condition prevents
a residual from being eliminated
merely by encoding the disputed distinction
into a larger carrier
or a specially constructed frame variable.

A representation is not faithful
when it succeeds only by placing
the desired answer into the input.

The following outcomes count
as substantive residuals.

\begin{enumerate}[label=(R\arabic*)]
  \item \emph{Missing comparison machinery.}
        A well-scoped schedule question
        requires a sameness distinction
        that no admitted occurrence
        or frame regime can express.

  \item \emph{Nonfunctional interpretation.}
        The complete declared inputs
        do not determine one interpretation result,
        and the remaining variation
        cannot be represented
        as an explicit frame component
        or probe parameter.

  \item \emph{Unfaithful frame reduction.}
        A frame-sensitive question
        can be represented only by
        erasing a distinction present
        in the source artifacts
        or adding a distinction
        not supplied by them.

  \item \emph{Unrepresentable threshold.}
        The threshold exists
        in the ambient lattice
        \(\operatorname{Eq}(\mathcal{A})\)
        but no occurrence regime
        admitted by the schedule vocabulary
        denotes that partition.

  \item \emph{Invalid component factorization.}
        A declared componentwise frame model
        cannot reproduce the kernel
        of the complete frame tuple,
        and the interaction cannot be expressed
        by the existing frame declaration.

  \item \emph{Persistent scope mismatch.}
        A question remains interpretive
        but cannot be placed
        on one fixed finite comparison artifact
        without changing the object
        whose equivalence is being asked about.
\end{enumerate}

The following outcomes do not,
by themselves,
count as residuals.

\begin{enumerate}[label=(N\arabic*)]
  \item a question whose object boundary is undeclared;
  \item a question whose occurrence
        or frame regime is undeclared;
  \item a question lacking a finite horizon;
  \item a resolution policy
        that fails partition validation;
  \item an inadmissible occurrence regime
        for which a finite witness is returned;
  \item two incomparable thresholds;
  \item a distinct frame pair
        with zero demand under the selected probes;
  \item a regime not exercised
        by the current finite corpus;
  \item a result outside the scope
        of interpretive equivalence; or
  \item an implementation that violates
        a proved algebraic law.
\end{enumerate}

These outcomes may be important findings,
but each has an explanation
within the declared construction.

A regime untouched by the corpus
bears specifically on domain-relative necessity.
It does not show that the regime collapses
or that the broader vocabulary is invalid.

Likewise,
two regimes performing the same work
on the finite corpus
do not establish global equality.
They establish only
that the selected evidence
does not separate them.

The analysis therefore reports
three claims independently:

\begin{enumerate}[label=(\alph*)]
  \item \emph{coverage}:
        how many corpus questions
        receive a non-residual disposition;

  \item \emph{necessity within the corpus}:
        which declared regimes
        are required by at least one question
        or computed threshold; and

  \item \emph{non-collapse within the corpus}:
        which regime pairs
        are separated by at least one finite witness.
\end{enumerate}

No one of these claims entails the others.

Complete coverage does not establish
that every admitted regime is necessary.

A regime exercised by one question
does not establish that it is distinct
from every other regime.

Failure to separate two regimes
on the present corpus
does not establish that they are globally equal.

\begin{remark}[Failure Modes of the Evidence Program]
  \label{ik.rem.evidence-failure}
  The evidence would materially weaken
  the schedule claim of this paper if:
  \begin{enumerate}[label=(\alph*)]
    \item nearly every computed kernel
          is discrete or indiscrete;
    \item changing the occurrence regime
          never changes a kernel
          or admissibility verdict;
    \item no nontrivial threshold is observed;
    \item every question is resolved
          only by declaring another parameter,
          without excluding
          or constraining any proposed comparison;
    \item the complete frame kernel
          always reduces trivially
          to independent component kernels;
    \item the vocabulary cannot express
          thresholds generated
          by ordinary schedule questions; or
    \item faithful residuals remain
          without a common structural diagnosis.
  \end{enumerate}
\end{remark}

The fourth failure mode is especially important.
A framework that answers every question only with
``declare another regime''
may accommodate every outcome without explaining any of them.

The evidence must show
not only that declarations make questions well posed,
but also that the resulting finite construction:

\begin{enumerate}[label=(\alph*)]
  \item rules out invalid comparisons;
  \item distinguishes admissible from inadmissible regime choices;
  \item produces nontrivial thresholds;
  \item exposes incomparable demands;
  \item rejects policies that are not equivalence relations; and
  \item returns finite witnesses for its negative verdicts.
\end{enumerate}

For the frozen corpus and pinned schedule artifacts,
a successful evidence result would establish
that the construction supplies a finite,
parameter-explicit, nondegenerate, and refutable account
of the examined interpretive demands.


% =============================================================================
\section{Related Work}
\label{ik.sec.related}
% =============================================================================

This paper draws on three distinct foundations.

\begin{itemize}
  \item The Structural Explainability series supplies
        the representational separation
        between declared structure,
        interpretive regimes,
        and operationally induced relations.

  \item Order theory supplies
        partition lattices,
        refinement,
        Galois connections,
        and the preservation laws of adjoint maps.

  \item Generalized-ultrametric and observational-equivalence work supplies
        close mathematical and semantic precedents
        for representing families of equivalence relations
        through pairwise demands indexed by observations.
\end{itemize}

The schedule realization additionally depends
on calendaring,
recurrence,
and time-zone standards
that define the artifacts
over which interpretation is performed.

The contribution is not a replacement
for any of these foundations.
It combines them into a finite construction
that transports equivalence requirements
between distinct interpreted carriers.

\subsection{Structural Explainability Foundations}

The present work extends the Structural Explainability sequence
developed in \emph{Neutral Substrates}, \emph{Referential Regimes}, and
\emph{Operational Identity}.

\emph{Neutral Substrates} separates stable referential structure
from causal and normative interpretation~\citep{case2026neutral}.
Its foundational layer preserves the referential commitments required
for shared use under persistent disagreement,
while framework-dependent interpretation
is introduced through attributed extensions.
That separation supplies the architectural premise of this paper.
A schedule declaration and the artifacts to which it refers
may remain fixed while different frame environments
produce different interpreted occurrence environments.
The interpretation is therefore not collapsed
into the represented declaration.

\emph{Referential Regimes} makes sameness relative
to an explicit identity basis and derives a family of transformation-invariant
referential regimes for neutral substrates~\citep{case2026referential}.
This paper inherits the methodological requirement
that a claim of sameness identifies the regime under which it is made.
This paper does not import particular referential regimes as occurrence or frame regimes.
Occurrence and frame regimes are equivalence relations
on the carriers declared by the present comparison artifact.

\emph{Operational Identity} represents declared and implemented rules of sameness
as partitions of one finite record domain and compares them
in the refinement lattice~\citep{case2026operational}.
It establishes the use of finite partition comparison,
refinement witnesses, and disclosure-relative computation
for examining identity-relevant behavior.
The present work uses that finite-partition method
but does not import the operational-identity audit itself.
It does not use implementation-surface registries,
sibling-regime classification,
regime-substitution verdicts,
or disclosure-completeness claims.

Instead, it introduces two distinct carriers:
a carrier of frame environments
and a carrier of occurrence environments.
An interpretation function connects them,
and a probe family determines
which interpreted outputs participate in the comparison.

The companion \emph{Operational Identity software}
provides an executable realization
of the earlier finite-audit method~\citep{case2026oiverification}.
The computational procedure of this paper
is a separate construction for kernels, generated regimes, thresholds, and witnesses.

The broader ontology-engineering literature treats explicit categories,
identity conditions, and persistence commitments
as central modeling concerns~\citep{arp2015building}.
Classification studies likewise emphasize that categories and standards
operate as consequential infrastructure
rather than as neutral labels~\citep{bowkerstar1999}.
The present work is narrower.
It does not propose a general ontology of frames or occurrences.
It assumes finite typed carriers and studies the equivalence relations
induced between them by interpretation.

\subsection{Partition Lattices and Galois Connections}

Equivalence relations on a fixed set form a complete lattice
under refinement~\citep{daveypriestley2002}.

Under the orientation used here,
a finer equivalence relation is lower in the order.
Meet is intersection, and join is the equivalence closure of union.

This lattice structure supplies the occurrence-regime lattice
\[
  \operatorname{Eq}(\mathcal{A})
\]
and the frame-regime lattice
\[
  \operatorname{Eq}(\mathcal{E}).
\]

Galois connections provide paired monotone maps
between ordered structures whose inequalities are equivalent
in opposite domains.
The general construction was introduced by Ore~\citep{ore1944galois}
and is now standard in lattice and order theory~\citep{daveypriestley2002}.

For a left adjoint \(L\) and right adjoint \(K\),
\[
  L(x)
  \sqsubseteq
  y
  \quad\Longleftrightarrow\quad
  x
  \sqsubseteq
  K(y).
\]
The left adjoint preserves joins, and the right adjoint preserves meets.

Abstract interpretation made lattice-theoretic adjunctions
operationally central
in program analysis~\citep{cousot1977abstract}.
There, concrete and abstract semantic domains
are related by maps that construct or approximate
program properties and fixpoints.

The present adjunction has a different role.

The two domains are lattices of equivalence relations over different carriers.

The right adjoint
\[
  K_P:
  \operatorname{Eq}(\mathcal{A})
  \longrightarrow
  \operatorname{Eq}(\mathcal{E})
\]
transports an occurrence-equivalence declaration
to the frame equivalence
that it induces under the probe family.

The left adjoint
\[
  L_P:
  \operatorname{Eq}(\mathcal{E})
  \longrightarrow
  \operatorname{Eq}(\mathcal{A})
\]
constructs the least occurrence regime
required to support a declared frame regime.

The existence and general preservation laws of adjoints are classical.
The contribution here is the explicit probe-generated pair
\[
  L_P
  \dashv
  K_P,
\]
its finite construction from an interpretation table,
and its interpretation
as an exact threshold
between admissible and inadmissible occurrence regimes.

\subsection{Lattice-Valued Generalized Ultrametrics}

Lawvere recast generalized metric spaces
as enriched categorical structures,
showing that numerical distance
is one instance of a broader ordered notion
of comparison~\citep{lawvere1973metric}.

Prie{\ss}-Crampe and Ribenboim
developed generalized ultrametric spaces
whose distance values lie in partially ordered structures
rather than necessarily
in the nonnegative real numbers~\citep{priesscrampe1996generalized}.

In such spaces, the lattice join plays the role
taken by maximum in an ordinary ultrametric inequality.
Distances may therefore express qualitatively different and incomparable demands.

A close mathematical precedent
for Section~\ref{ik.sec.representation}
is Braunfeld's correspondence
between lattice-valued ultrametric spaces
and structures equipped with lattice-indexed families
of equivalence relations~\citep{braunfeld2019lambda}.

In that correspondence, the distance between two points
is constructed from the finest equivalence relation
that identifies them, and the equivalence family
can be recovered by thresholding those distance values.

The broad correspondence between meet-preserving equivalence families
and lattice-valued ultrametrics is therefore prior work.
The representation result of Section~\ref{ik.sec.representation}
is a finite interpretive specialization
of that correspondence.
The specialization has three features specific to the present construction.

First,
the value lattice is itself
the complete partition lattice
\[
  \operatorname{Eq}(\mathcal{A})
\]
of occurrence environments.

Second,
the equivalence family
is not supplied independently.
It is generated from
a finite probe family,
an interpretation function,
and an occurrence regime.

Third,
the pairwise value
\[
  \delta_P(F,F')
\]
has an operational threshold meaning:
it is the least occurrence regime
under which the selected probes
make the two frame environments
interpretively equivalent.

Terminology differs across
the generalized-metric literature.
Some definitions of ultrametric
require separation,
so that zero distance implies equality.

The present construction does not require
that property.
Distinct frame environments may have
\[
  \delta_P(F,F')
  =
  \bot_{\mathcal{A}}
\]
when the probe family produces
identical canonical occurrence environments
under both.

Under terminology that reserves
\emph{ultrametric}
for separating distances,
the pairwise interpretive demand
is therefore a generalized pseudoultrametric
before quotienting by
the zero-demand equivalence.

\subsection{Observational and Testing Equivalence}

Testing semantics defines behavioral equivalence
through the observations that a selected family of tests can make.

De Nicola and Hennessy show how sets of tests
induce distinct process equivalences,
including may-testing and must-testing equivalence~\citep{denicola1984testing}.
Testing equivalence itself refines the older notion of
observational equivalence \citep{morris1968lambda, plotkin1977lcf},
which individuates terms by the contexts that distinguish them.

Two processes are equivalent when the permitted observations
cannot distinguish them under the selected testing interpretation.
The probe family of this paper has a related observational role.

Two frame environments lie in the same interpretive-kernel block
when every selected probe produces occurrence environments
that are equivalent under the declared occurrence regime.

There are, however, two substantive differences.
First, the observation result is not reduced
to a fixed success or failure value.
A probe produces a structured occurrence environment,
possibly containing ordered values,
multiplicity,
boundary information,
and a terminal tag.
Second, equivalence of those results is itself a constitutive parameter.
The same probes may induce different frame equivalences
under different occurrence regimes.

The construction therefore moves from
\[
  \text{equivalence under a family of tests}
\]
to
\[
  \text{equivalence under a family of probes
    and a declared output-equivalence regime}.
\]

The reverse map
has no ordinary analogue
in basic testing equivalence.
Given a declared frame regime,
\(L_P\) computes
the least occurrence regime
under which the probes support that declaration.

The resulting threshold makes explicit which output distinctions
must be suppressed to preserve a requested observational equivalence.

Testing equivalence is used only as the nearest semantic precedent
for probe-relative indistinguishability;
the present setting has no processes or nondeterministic execution,
and no process-equivalence machinery is imported.

\subsection{Canonicalization and Quotient Comparison}

Canonical forms
are a standard method
for making comparison independent
of incidental presentation.
Term-rewriting theory,
for example,
uses normal forms
to replace repeated comparison
over many equivalent representations
with comparison of selected representatives~\citep{baadernipkow1998}.

The schedule construction uses
the same broad strategy.

Raw generated occurrences
are placed in a declared
chronological canonical form.
The occurrence regime
is then applied
to that indexed representation.

Canonicalization is not treated
as a proof that
one representation is intrinsically privileged.
It is part of the declared comparison protocol
and determines the carrier
on which the kernel is computed.

Matching and entity-resolution methods
address a different problem:
finding plausible correspondences
among records
under imperfect,
heterogeneous,
or noisy identifying information~\citep{christen2012datamatching}.

Such methods may properly optimize
over alternative candidate matches.

The rejected matching comparison
of Section~\ref{ik.sec.finite-counterexamples}
is not rejected
because matching is generally invalid.
It is rejected
because minimizing over alternative bijections
does not preserve
the lattice laws required
by this construction.

Separate matchings may witness
comparison under separate occurrence regimes,
while their meet may admit
no single matching
that witnesses both.

The resulting comparison
can fail meet preservation
and even fail
the strong triangle law.

Permutation invariance alone
is therefore insufficient
for the present purpose.

The construction requires
a declared comparison
whose threshold relations
are equivalence relations
and whose induced kernel
preserves meets.

Chronological canonicalization
provides that structure
for the schedule realization
without factorial enumeration
of possible occurrence matchings.

\subsection{Calendar, Recurrence, and Time-Zone Semantics}

Internet calendaring standards
supply the primary representational foundation
for the schedule realization.

RFC~5545 defines
the iCalendar data model,
recurrence rules,
exception dates,
floating date-times,
UTC date-times,
zone-referenced local times,
and embedded time-zone components~\citep{rfc5545}.

A recurring component
may describe an indefinitely extensible
set of occurrences.
A referenced time-zone definition
must provide valid information
for the recurrence instances
to which it applies.

These features motivate
two explicit declarations
in the present model:
a finite observation horizon
and a frame environment
containing the cited time-zone artifact.

JSCalendar provides
an independently defined JSON data model
for calendar events and tasks
and aims to reduce representational ambiguity
while remaining broadly compatible
with iCalendar semantics~\citep{rfc8984}.

These standards determine
how schedules and calendar values
may be represented.
They do not determine
which distinct representations,
artifact releases,
or interpreted occurrence sequences
must count as equivalent
for a particular comparison.

Time-zone interpretation additionally depends
on maintained and versioned external artifacts.

RFC~6557 specifies procedures
for maintaining
the IANA Time Zone Database~\citep{rfc6557}.
RFC~7808 defines a service
for distributing time-zone data
and associated synchronization information
to clients~\citep{rfc7808}.
RFC~9636 specifies
the Time Zone Information Format
used to represent and exchange
time-zone transition data~\citep{rfc9636}.

Together,
these standards establish
that calendaring systems consume
external,
mutable,
versioned interpretive artifacts.

They do not define
an identity regime over releases.

For example,
the standards do not decide
whether two time-zone releases
that have the same current extension
but different identifiers or histories
are the same frame.

They likewise do not determine:

\begin{enumerate}[label=(\alph*)]
  \item which frame environments
        a deployment declaration requires
        to remain equivalent;
  \item which differences
        among generated occurrence environments
        are material;
  \item the least occurrence regime
        supporting a declared frame regime;
  \item whether a pinning
        or fallback practice
        is an equivalence relation
        rather than a selection policy; or
  \item the finite witness
        refuting an unsupported
        frame-equivalence declaration.
\end{enumerate}

Those are the questions
addressed by the schedule interpretation
of this paper.

Taken together,
the related foundations delimit
the contribution precisely.

This paper does not introduce
partition lattices,
Galois connections,
adjoint preservation laws,
lattice-valued generalized ultrametrics,
testing equivalence,
canonical forms,
recurrence syntax,
or time-zone formats.

It contributes:

\begin{enumerate}[label=(\arabic*)]
  \item a finite typed interpretation setting
        with separate occurrence
        and frame carriers;

  \item a kernel indexed jointly
        by a probe family
        and an occurrence regime;

  \item an explicit adjoint construction
        producing the least occurrence regime
        required by a frame regime;

  \item a lattice-valued threshold interpretation
        of that adjunction;

  \item pairwise interpretive demands
        as the generalized-pseudoultrametric
        representation
        of the resulting kernel family;

  \item finite counterexamples
        delimiting canonicalization,
        matching,
        and partition-hood requirements;

  \item a near-linear partition procedure
        with finite witness extraction; and

  \item a schedule realization
        over finite horizons
        and versioned interpretive frames.
\end{enumerate}


% =============================================================================
\section{Limitations}
\label{ik.sec.limitations}
% =============================================================================

The interpretive-kernel construction is exact
relative to a finite declared comparison artifact.

It does not determine
which objects,
probes,
frame environments,
occurrence environments,
canonicalization rules,
or equivalence regimes
a domain ought to declare.

Those choices remain part
of the domain specification.
The construction computes
their consequences.

\subsection{Finite Declared Scope}

The results assume finite carriers
and a finite probe family.

A schedule declaration may denote
an indefinitely extensible recurrence,
but each probe must supply
a finite observation horizon.

The resulting kernel therefore characterizes
the selected finite comparison artifact.
It does not establish equivalence
over all possible probes,
all future horizons,
or every release
of an external frame artifact.

A later probe,
longer horizon,
or enlarged frame carrier
may distinguish environments
that the original artifact
placed in one kernel block.

\subsection{Dependence on Probe Family}

Interpretive equivalence is observational.
Two frame environments may be equivalent
under one probe family
and distinguishable under another.

Consequently, a reported kernel is incomplete
unless it identifies the probe family under which it was computed.

Zero pairwise demand does not establish
that two frame environments are identical.
It establishes only that the selected probes
produce equivalent occurrence environments
under the selected occurrence regime.

\subsection{Dependence on the Occurrence Regime}

The occurrence regime is a constitutive parameter.
The same interpretation table
may induce different frame partitions
under different occurrence regimes.
The construction does not decide
which occurrence distinctions
are substantively appropriate.
It determines whether a proposed occurrence regime
supports a declared frame regime and computes the least regime
that would do so.

Admissibility therefore establishes formal compatibility
with a declaration, not the correctness of the declaration itself.

\subsection{Interpretation and Canonicalization}

The construction assumes a total interpretation function
on the selected finite artifact.
It does not resolve semantic ambiguity
left undeclared by the domain.

Runtime state, provider behavior,
unstated defaults,or incomplete external artifacts
must either be represented as declared frame components
or excluded from the comparison claim.

Canonicalization is likewise declared.
Chronological canonicalization removes incidental generation order,
but it may not silently change the represented denotation.
A different canonicalization defines a different comparison artifact
unless representative independence is proved separately.

The counterexamples of
Section~\ref{ik.sec.finite-counterexamples}
show that permutation invariance alone
does not justify replacing canonicalization
with optimization over alternative matchings.

\subsection{Frame Completeness}

A frame environment may contain multiple interacting artifacts,
including time-zone rules, holiday calendars,
business calendars, and provider data.
The kernel is computed only over the components declared
in the frame tuple.

An omitted component can produce apparent equivalence
that disappears when the environment is completed.

The present work does not establish
that kernels over complete frame tuples
factor into independent kernels over individual components.
Such factorization requires a separate theorem
or finite domain evidence.

\subsection{Exact Rather than Graded Equivalence}

Occurrence and frame regimes
are ordinary equivalence relations.
The value structure
\[
  \operatorname{Eq}(\mathcal{A})
\]
records exact equivalence requirements,
not probabilities,
degrees of similarity,
confidence values,
or fuzzy truth values.

The construction therefore does not directly cover
approximate schedule matching,
probabilistic interpretation,
uncertain external data,
or graded equivalence.

An extension to fuzzy equivalences,
residuated lattices,
or other ordered relation structures
would require new definitions
and new preservation results.

\subsection{Nonseparation}

The pairwise interpretive demand
is generally a pseudoultrametric
rather than a separating ultrametric.
Distinct frame environments may have
\[
  \delta_P(F,F')
  =
  \bot_{\mathcal{A}}.
\]

The construction does not recover
representational frame identity
from observational equality.

Separation can be obtained
only by strengthening the probe family
or by quotienting the frame carrier
by zero demand.

\subsection{No Normative Selection Procedure}

The kernel does not determine which frame release
a system should select, whether an institution should pin
or update an external artifact, or which historical,
legal, or organizational notion of sameness is preferable.
Pinning, latest-compatible selection,
always-latest selection,
and fallback behavior
are resolution policies.
The construction can test
whether a declared relation
is an equivalence relation,
but it does not convert
a selection policy into an identity regime.

\subsection{No Implementation-Completeness Claim}

The finite procedure of Section~\ref{ik.sec.computational-procedure}
operates on the supplied interpretation table and declared carriers.
It does not establish that every deployed implementation path
uses the declared interpreter, that every external dependency
has been disclosed, or that every runtime behavior
is represented in the table.
Those questions belong to implementation conformance,
provenance, and disclosure auditing.
A correct kernel computation over incomplete inputs
remains correct only for those inputs.

\subsection{Computational Cost}

The partition operations are near-linear
in the finite artifact size, but schedule interpretation
may be substantially more expensive.
Long horizons, dense recurrences,
large frame inventories,
or many output-rich probes
can make construction
of the interpretation table
the dominant cost.
The complexity bounds do not make recurrence expansion
constant-time, nor do they prevent the canonical outputs themselves
from becoming large.

\subsection{Evidence Status}

The schedule construction of Section~\ref{ik.sec.schedule-interpretation}
shows that the theory admits a nondegenerate finite realization.
It is not an empirical result about production schedule systems.

The corpus classification,
pinned artifact comparisons,
kernel computations,
threshold comparisons,
and residual analysis
described in
Section~\ref{ik.sec.evidence-plan}
remain separate evidence obligations.

Until those obligations are discharged,
the paper establishes
the mathematical framework
and its executable procedure,
not the prevalence
or practical distribution
of particular schedule kernels.

\subsection{Domain Transfer}

The abstract construction is not restricted to scheduling,
but the present domain realization is.
Applying it elsewhere requires:

\begin{enumerate}
  \item finite typed carriers;
  \item a total interpretation function on the selected artifact;
  \item declared probes;
  \item declared equivalence regimes;
  \item a comparison-preserving canonicalization; and
  \item a faithful account of the external frame environment.
\end{enumerate}

The paper does not establish that every interpretive domain
admits these declarations naturally or that the schedule signature
generalizes unchanged.
A domain in which interpretation is irreducibly probabilistic,
nondeterministic, continuous, or dependent on undeclarable runtime state
may require a different construction.
These limitations delimit rather than negate the principal claim.

Given finite declared carriers, a probe family,
a total interpretation function,
and an occurrence regime,
the interpretive kernel
and its least supporting thresholds
are computable, their algebraic properties are exact,
and failed admissibility claims have finite witnesses.

% =============================================================================
\section{Conclusion}
\label{ik.sec.conclusion}
% =============================================================================

An interpreted system contains a rule of sameness for the interpretive frames,
a rule of sameness for the outputs produced under those frames,
and an interpretation function connecting the two.
The interpretive kernel makes that connection a first-class object.

For a finite probe family \(P\), the map
\[
  K_P:
  \operatorname{Eq}(\mathcal{A})
  \longrightarrow
  \operatorname{Eq}(\mathcal{E})
\]
sends an occurrence regime to the frame equivalence that it induces.

Its left adjoint
\[
  L_P:
  \operatorname{Eq}(\mathcal{E})
  \longrightarrow
  \operatorname{Eq}(\mathcal{A})
\]
sends a frame regime to the least occurrence regime needed to support it.

The adjunction
\[
  L_P(\varphi)
  \sqsubseteq
  \rho
  \quad\Longleftrightarrow\quad
  \varphi
  \sqsubseteq
  K_P(\rho)
\]
turns compatibility between the two declarations into a decidable refinement condition.

It also explains the preservation laws.
Interpretive kernels preserve meets because satisfying several occurrence regimes
means satisfying their common refinement.
Generated occurrence demands preserve joins
because supporting combined frame demands
requires the equivalence closure of their combined output requirements.

Pairwise interpretive demand provides a second view of the same construction.
For two frame environments,
\[
  \delta_P(F,F')
\]
is the least occurrence regime that makes them indistinguishable
under the selected probes.

Thresholding these values recovers the kernel family.
The values satisfy symmetry, zero self-demand,
and the strong triangle inequality in the partition lattice.
They need not separate distinct frames, because observational equality
under a finite probe family does not imply representational identity.

The finite counterexamples delimit the result.
Index-wise comparison does not descend
from unordered presentations
without canonicalization.
Permutation-invariant matching
does not guarantee
meet preservation
or generalized-ultrametric structure.
A deterministic frame-selection policy
does not become an identity regime
unless it is an equivalence relation.
These failures identify the conditions
under which the kernel construction
is well defined and its algebraic guarantees remain valid.

The computational procedure turns the theory into a finite audit.
A cached interpretation table supports kernel construction
through probe signatures.
Generating forests support threshold construction
without enumerating every frame pair.
Union-find computes the required equivalence closures.
Failed admissibility returns a frame pair,
a probe, and two occurrence environments
that witness the unsupported declaration.

The schedule realization shows how the abstract construction
applies when interpretation depends
on mutable, versioned external artifacts.
A schedule declaration may remain fixed
while time-zone,
holiday-calendar,
business-calendar,
or provider artifacts change.
Civil output and resolved-instant output remain separate typed dimensions.
Finite horizons make recurrence interpretation executable.
Chronological canonicalization removes incidental generation order
without introducing favorable matching.

The resulting kernel states which frame environments
the selected schedule probes cannot distinguish
under the declared occurrence regime.
The paper does not determine
which regime is substantively correct.
It does not choose a time-zone release,
holiday authority,
resolution policy,
or institutional rule of sameness.
It computes the consequences
of those declarations
and identifies the least occurrence demand
that each frame declaration imposes.
Its guarantees remain relative
to the declared carriers,
probes,
horizons,
frame tuple,
canonicalization,
and occurrence regime.

The evidence plan therefore remains essential.
The mathematical construction establishes exact finite laws.
The schedule evidence must establish that real comparison artifacts
produce nondegenerate kernels, nontrivial thresholds,
regime-sensitive classifications, and finite witnesses
that matter in practice.

A faithful residual would identify a well-scoped interpretive demand
that the declared construction cannot represent
without adding or erasing a distinction in the source question.
Such residuals would delimit the schedule realization
and may motivate a broader semantic structure.

Within its declared finite scope, the result is complete.
Given two finite carriers, a finite probe family,
a total interpretation function, and an occurrence regime,
the induced frame equivalence is computable.
Given a frame regime, its least supporting occurrence regime is computable.
The two constructions form an adjunction.
Their pairwise demands admit a lattice-valued
generalized-pseudoultrametric representation.
And when a proposed occurrence regime
fails to support a declared frame demand,
the failure has a finite witness.

Interpretive equivalence is a declared, parameter-relative, and executable relation
between rules of sameness on opposite sides of interpretation.

% =============================================================================
\section*{Statements and Declarations}
% =============================================================================

\subsection*{Acknowledgments}
The author thanks the anonymous reviewers of earlier papers in this series
for comments that improved the framing, organization, and presentation of this work.
% This research made use of SciX,
% a scientific literature search and discovery platform,
% whose related-literature tools helped identify relevant work across disciplinary boundaries.

\subsection*{Author Contributions}
The author is the sole contributor to this work
and is responsible for all aspects of the research, authorship, and publication.

% \subsection*{Code Availability}
% A reference implementation that checks the finite constructions of this paper
% is available~\citep{case2026oiverification}.
% It provides two independent checkers of the finite comparison procedure,
% an explicit quadratic reference and the near-linear labelled procedure
% of Theorem~\ref{oi.thm.decidable},
% and confirms they agree on the worked examples (Example~\ref{oi.ex.loc-obj-fork},
% Propositions~\ref{oi.prop.aligned-not-substitution}
% and~\ref{oi.prop.witness-not-alignment},
% Remark~\ref{oi.rem.nonvacuity}, the Section~\ref{oi.sec.worked} example,
% and Proposition~\ref{oi.prop.nonmonotone}) and across small and randomized instances.
% The implementation does not discover
% mechanisms, evaluate a deployed system, or discharge the completeness
% claims of Section~\ref{oi.sec.disclosure}.

\subsection*{Use of AI-Assisted Tools}
AI-assisted tools were used for editing, formatting, and consistency checking.
The author reviewed all suggestions and is solely responsible for the content.

\subsection*{Declaration of Conflicting Interest}
The author declares no competing interests.

\begin{thebibliography}{99}
  \small

  \bibitem[Arp et al.(2015)]{arp2015building}
  R.\ Arp, B.\ Smith, and A.\ D.\ Spear.
  \emph{Building Ontologies with Basic Formal Ontology}.
  MIT Press, 2015.

  \bibitem[Baader and Nipkow(1998)]{baadernipkow1998}
  F.\ Baader and T.\ Nipkow.
  \emph{Term Rewriting and All That}.
  Cambridge University Press, Cambridge, UK, 1998.

  \bibitem[Bowker and Star(1999)]{bowkerstar1999}
  G.\ C.\ Bowker and S.\ L.\ Star.
  \emph{Sorting Things Out: Classification and Its Consequences}.
  MIT Press, Cambridge, MA, 1999.

  \bibitem[Braunfeld(2019)]{braunfeld2019lambda}
  S.\ Braunfeld.
  \(\Lambda\)-ultrametric spaces and lattices of equivalence relations.
  \emph{Algebra Universalis}, 80(3), 2019.
  \href{https://doi.org/10.1007/s00012-019-0606-4}{doi:10.1007/s00012-019-0606-4}.

  \bibitem[Case(2026a)]{case2026neutral}
  D.\ M.\ Case.
  \emph{Neutral Substrates: A Design Constraint for Shared Records Under Persistent Interpretive Disagreement}.
  arXiv preprint \href{https://arxiv.org/abs/2601.14271}{arXiv:2601.14271}, 2026.

  \bibitem[Case(2026b)]{case2026referential}
  D.\ M.\ Case.
  \emph{Referential Regimes: Transformation-Based Identity Profiles for Neutral Substrates}.
  arXiv preprint \href{https://arxiv.org/abs/2601.16152}{arXiv:2601.16152}, 2026.

  \bibitem[Case(2026c)]{case2026operational}
  D.\ M.\ Case.
  \emph{Operational Identity: A Finite Audit of Declared and Implemented Rules of Sameness}.
  arXiv preprint \href{https://arxiv.org/abs/2607.20729}{arXiv:2607.20729}, 2026.

  \bibitem[Case(2026d)]{case2026oiverification}
  D.\ M.\ Case.
  \emph{Structural Explainability Verification: Operational Identity}.
  Software, version 0.1.0, 2026.
  \href{https://doi.org/10.5281/zenodo.21499599}{doi:10.5281/zenodo.21499599}.

  \bibitem[Christen(2012)]{christen2012datamatching}
  P.\ Christen.
  \emph{Data Matching: Concepts and Techniques for Record Linkage, Entity Resolution, and Duplicate Detection}.
  Springer, Berlin, 2012.

  \bibitem[Cousot and Cousot(1977)]{cousot1977abstract}
  P.\ Cousot and R.\ Cousot.
  Abstract interpretation:
  A unified lattice model for static analysis of programs
  by construction or approximation of fixpoints.
  In \emph{Proceedings of the 4th ACM SIGACT-SIGPLAN Symposium on Principles of Programming Languages},
  pages 238--252, 1977.
  \href{https://doi.org/10.1145/512950.512973}{doi:10.1145/512950.512973}.

  \bibitem[Davey and Priestley(2002)]{daveypriestley2002}
  B.\ A.\ Davey and H.\ A.\ Priestley.
  \emph{Introduction to Lattices and Order}.
  Second edition.
  Cambridge University Press, Cambridge, UK, 2002.

  \bibitem[De Nicola and Hennessy(1984)]{denicola1984testing}
  R.\ De Nicola and M.\ C.\ B.\ Hennessy.
  Testing equivalences for processes.
  \emph{Theoretical Computer Science}, 34(1--2):83--133, 1984.
  \href{https://doi.org/10.1016/0304-3975(84)90113-0}{doi:10.1016/0304-3975(84)90113-0}.

  \bibitem[Desruisseaux(2009)]{rfc5545}
  B.\ Desruisseaux, editor.
  \emph{Internet Calendaring and Scheduling Core Object Specification (iCalendar)}.
  RFC 5545, Internet Engineering Task Force, September 2009.
  \href{https://doi.org/10.17487/RFC5545}{doi:10.17487/RFC5545}.

  \bibitem[Douglass and Daboo(2016)]{rfc7808}
  M.\ Douglass and C.\ Daboo.
  \emph{Time Zone Data Distribution Service}.
  RFC 7808, Internet Engineering Task Force, March 2016.
  \href{https://doi.org/10.17487/RFC7808}{doi:10.17487/RFC7808}.

  \bibitem[Jenkins and Stepanek(2021)]{rfc8984}
  N.\ Jenkins and R.\ Stepanek.
  \emph{JSCalendar: A JSON Representation of Calendar Data}.
  RFC 8984, Internet Engineering Task Force, July 2021.
  \href{https://doi.org/10.17487/RFC8984}{doi:10.17487/RFC8984}.

  \bibitem[Lawvere(1973)]{lawvere1973metric}
  F.\ W.\ Lawvere.
  Metric spaces, generalized logic, and closed categories.
  \emph{Rendiconti del Seminario Matematico e Fisico di Milano},
  43:135--166, 1973.
  \href{https://doi.org/10.1007/BF02924844}{doi:10.1007/BF02924844}.

  \bibitem[Lear and Eggert(2012)]{rfc6557}
  E.\ Lear and P.\ Eggert.
  \emph{Procedures for Maintaining the Time Zone Database}.
  BCP 175, RFC 6557, Internet Engineering Task Force, February 2012.
  \href{https://doi.org/10.17487/RFC6557}{doi:10.17487/RFC6557}.

  \bibitem[Morris(1968)]{morris1968lambda}
  J.\ H.\ Morris.
  \emph{Lambda-Calculus Models of Programming Languages}.
  PhD thesis, Massachusetts Institute of Technology, 1968.
  \href{https://hdl.handle.net/1721.1/149376}{https://hdl.handle.net/1721.1/149376}.

  \bibitem[Olson et al.(2024)]{rfc9636}
  A.\ D.\ Olson, P.\ Eggert, and K.\ Murchison.
  \emph{The Time Zone Information Format (TZif)}.
  RFC 9636, Internet Engineering Task Force, October 2024.
  \href{https://doi.org/10.17487/RFC9636}{doi:10.17487/RFC9636}.

  \bibitem[Ore(1944)]{ore1944galois}
  {\O}.\ Ore.
  Galois connexions.
  \emph{Transactions of the American Mathematical Society},
  55:493--513, 1944.
  \href{https://doi.org/10.1090/S0002-9947-1944-0010555-7}{doi:10.1090/S0002-9947-1944-0010555-7}.

  \bibitem[Plotkin(1977)]{plotkin1977lcf}
  G.\ D.\ Plotkin.
  {LCF} considered as a programming language.
  \emph{Theoretical Computer Science},
  5(3):223--255, 1977.
  \href{https://doi.org/10.1016/0304-3975(77)90044-5}{doi:10.1016/0304-3975(77)90044-5}.

  \bibitem[Priess-Crampe and Ribenboim(1996)]{priesscrampe1996generalized}
  S.\ Prie{\ss}-Crampe and P.\ Ribenboim.
  Generalized ultrametric spaces I.
  \emph{Abhandlungen aus dem Mathematischen Seminar der Universit\"at Hamburg},
  66:55--73, 1996.

  \bibitem[Schmidt(1986)]{schmidt1986}
  David A.\ Schmidt.
  \emph{Denotational Semantics: A Methodology for Language Development}.
  Allyn and Bacon, Boston, MA, 1986.


\end{thebibliography}

\end{document}
